\chapter{Doob-Meyer Theorem}
\label{chap:doob_meyer}


This chapter starts with the derivation of a Komlòs lemma, which is a useful tool to extract converging subsequences from bounded sequences of functions.
Then, we give a short review of the properties of the Doob decomposition of an adapted process indexed on a discrete set, and then follows \cite{Beiglböck_Schachermayer_Veliyev_2012} which gives an elementary and short proof of the Doob-Meyer theorem.




\section{Komlòs Lemma}


%technically a more general version with Cesaro sums exists, but it is not needed for this case
%(see "J. Komlòs, A generalization of a problem of Steinhaus, Acta Math. Acad. Sci. Hungar. 18 (1967) 217–229").

\begin{lemma}\label{lem:komlos_convex}
  \leanok
  \lean{komlos_convex}
Let $(f_n)_{n\in\mathbb{N}}$ be a sequence in a vector space $E$ and $\phi : E \to \mathbb{R}_+$ be a function such that $\phi(f_n)$ is a bounded sequence.
For $\delta > 0$, let $S_\delta = \{(f, g) \mid \phi(f)/2 + \phi(g)/2 - \phi((f+g)/2) \ge \delta\}$.
Then there exist $g_n\in convex(f_n,f_{n+1},\cdots)$ such that for all $\delta > 0$, for $N$ large enough and $n, m \ge N$, $(g_n, g_m) \notin S_\delta$.
\end{lemma}

\begin{proof}
  \leanok
Let $B$ be the bound of $(\phi(f_n))_{n\in\mathbb{N}}$.
Then for all $n\in\mathbb{N}$ and $g\in convex(f_n,f_{n+1},\cdots)$ we have $\phi(g)\le B$ by convexity of $\phi$.
Let $r_n = \inf(\phi(g) \mid g\in convex(f_n, f_{n+1},\ldots))$.
By construction $(r_n)_{n\in\mathbb{N}}$ is nondecreasing.
Let $A = \sup_{n \ge 1} r_n$, which is finite (as $A \le B$) and for each $n$ we may pick some $g_n\in convex(f_n, f_{n+1},\ldots)$ such that $\phi(g_n) \le A+1/n$ by $\inf$ and $\sup$ definitions.

Let $\varepsilon \in (0, \delta/4)$.
By properties of $\sup$ there exists $\bar{n}$ such that $r_{\bar{n}} \ge A-\varepsilon$ and such that $\frac{1}{\bar{n}} \le \varepsilon$.
Let $m \ge k \ge \bar{n}$.
We have $(g_k+g_m)/2 \in convex(f_k,f_{k+1},\ldots)$ and it follows since $(r_n)_{n\in\mathbb{N}}$ is nondecreasing that $\phi((g_k+g_m)/2) \ge A - \varepsilon$.
Hence due to the ordering of $m,k,\bar{n}$,
\begin{align*}
  \phi(g_k)/2 + \phi(g_m)/2 - \phi((g_k+g_m)/2)
  &\le 2(A + \frac{1}{\bar{n}}) - 2(A - \varepsilon)
  \\
  &\le 4 \varepsilon
  \\
  &< \delta
  \: .
\end{align*}
Thus, for $n, m \ge \bar{n}$, $(g_n, g_m) \notin S_\delta$.
\end{proof}


\begin{lemma}\label{lem:komlos_norm}
  \leanok
  \lean{komlos_norm}
Let $H$ be a Hilbert space and $(f_n)_{n\in\mathbb{N}}$ a bounded sequence in $H$. Then there exist functions $g_n\in convex(f_n,f_{n+1},\cdots)$ such that $(g_n)_{n\in\mathbb{N}}$ converges in $H$.
\end{lemma}

\begin{proof}
  \leanok
  \uses{lem:komlos_convex}
Consider $\phi : H \to \mathbb{R}_+$ defined by $\phi(f) = \|f\|_2^2$, which is convex.
Then Lemma \ref{lem:komlos_convex} applied to $(f_n)_{n\in\mathbb{N}}$ and $\phi$ gives us functions $g_n\in convex(f_n,f_{n+1},\cdots)$ such that for every $\delta>0$ there exists $N$ such that for $n,m\geq N$, $(g_n,g_m)\notin S_\delta$.
Thus for every $\delta>0$ there exists $N$ such that for $n,m\geq N$,
\begin{align*}
  \|g_n\|_2^2/2 + \|g_m\|_2^2/2 - \|(g_n+g_m)/2\|_2^2
  &< \delta
  \: .
\end{align*}
But the left-hand side is equal to $\|g_n - g_m\|_2^2/4$ by the parallelogram identity, hence $(g_n)_{n\in\mathbb{N}}$ is a Cauchy sequence in $H$ and thus converges in $H$ by completeness.
\end{proof}


\begin{lemma}\label{lem:convex_of_converg_seq_is_converg}\lean{TendstoUniformly_convexTail}\leanok
  Let $(x_n)_{n \in \mathbb{N}}$ be a sequence in a real vector space converging to $x$.
  Let $\mathcal{C}((x_n))$ be the set of sequences $(y_n)_{n \in \mathbb{N}}$ such that for all $n$, $y_n \in convex(x_n, x_{n+1}, \ldots)$.
  Then we have uniform convergence over $\mathcal{C}((x_n))$: for all $\varepsilon > 0$, there exists $\bar{n}$ such that for all $n \ge \bar{n}$, for all $(y_n)_{n \in \mathbb{N}} \in \mathcal{C}((x_n))$, $\Vert y_n - x \Vert \le \varepsilon$;
\end{lemma}

\begin{proof} \leanok
Let $\varepsilon>0$.
By convergence of $x_n$, there exists $\bar{n}$ such that for all $n \ge \bar{n}$, $\Vert x_n-x \Vert \le \varepsilon$.
Let $a_{n, m}$ be convex weights such that $y_n = \sum_{m = n}^{N_n} a_{n, m} x_m$.
By triangular inequality it follows that for $n \ge \bar{n}$,
\begin{align*}
  \Vert y_n - x \Vert
  = \left\Vert \sum_{m = n}^{N_n} a_{n, m} x_m - x \right\Vert
  = \left\Vert \sum_{m = n}^{N_n} a_{n, m} (x_m - x) \right\Vert
  \le \sum_{m = n}^{N_n} a_{n, m} \Vert x_m - x \Vert
  \le \varepsilon
  \: .
\end{align*}
\end{proof}

By convex weights on $\mathbb{N}$, we mean a sequence of non-negative real numbers $(a_n)_{n \in \mathbb{N}}$ with finitely many nonzero entries such that $\sum_{n \in \mathbb{N}} a_n = 1$.

\begin{definition}\label{def:convex_weights_product}
  \leanok
  \lean{Convexity.StdSimplex.bind}
If $(a_m)_{m \in \mathbb{N}}$ are convex weights and $(b^n_m)_{n,m \in \mathbb{N}}$ is such that for all $n$, the $(b^n_m)$ are convex weights, then we denote by $(a_\cdot) * (b^\cdot_\cdot)$ the convex weights defined by $((a_\cdot) * (b^\cdot_\cdot))_m = \sum_{k} a_k b^k_m$.
\end{definition}


\begin{lemma}\label{lem:komlos_convex_weights}
  \lean{komlos_convex_weights}
  \leanok
Let $E$ be a Hilbert space and for $i \in \mathbb{N}$, let $(x_n^{(i)})_{n \in \mathbb{N}}$ be a bounded sequence in $E$. Then there exists a sequence of convex weights $(\lambda^{k,n}_\cdot)_{k, n \in \mathbb{N}}$ with $\lambda^{k,n}_m = 0$ for $m < n$ such that for all $k \in \mathbb{N}$, $\left(\sum_{m \ge n} \left((\lambda^{k,n}_\cdot) * \ldots * (\lambda^{1,\cdot}_\cdot)\right)_m x_m^{(k)}\right)_{n \in \mathbb{N}}$ converges.
\end{lemma}

\begin{proof}
  \uses{lem:komlos_norm,def:convex_weights_product}
  \leanok
First by Lemma~\ref{lem:komlos_norm} applied to $(x_n^{(1)})_{n\in\mathbb{N}}$ in the Hilbert space $E$, there exist $g_n^1 \in convex(x_n^{(1)}, x_{n+1}^{(1)}, \ldots)$ (call its weights $\lambda^{1,n}_n,\cdots,\lambda^{1,n}_{N^1_n}$) such that $g_n^1$ converges to some $g^1$.

Secondly define $\tilde{g}_n^1$, convex combination of $x_n^{(2)}, x_{n+1}^{(2)}, \ldots$ with weights $\lambda^{1,n}_n,\cdots,\lambda^{1,n}_{N^1_n}$.
Applying lemma~\ref{lem:komlos_norm} to $(\tilde{g}_n^1)_{n\in\mathbb{N}}$ gives us $g_n^2 \in convex(\tilde{g}_n^1, \tilde{g}_{n+1}^1, \ldots)$ (call its weights $\lambda^{2,n}_n,\cdots,\lambda^{2,n}_{N^2_n}$) such that $g_n^2$ converges to some $g^2$.
$g_n^2$ is a convex combination of $x_n^{(2)}, x_{n+1}^{(2)}, \ldots$ with weights $(\lambda^{2,n}_\cdot) * (\lambda^{1,\cdot}_\cdot)$.

We continue iterating this process inductively.
At iteration $k$ we have weights $(\lambda^{k,n}_\cdot * \ldots * \lambda^{1,\cdot}_\cdot)$.
We define $\tilde{g}_n^k$ as the convex combination of $x_n^{(k+1)}, x_{n+1}^{(k+1)}, \ldots$ with those weights:
\[ \tilde{g}_n^k = \sum_{l} (\lambda^{k,n}_\cdot * \ldots * \lambda^{1,\cdot}_\cdot)_l x^{(k+1)}_l \]
We apply Lemma~\ref{lem:komlos_norm} to $(\tilde{g}_n^k)_{n\in\mathbb{N}}$ to get $g_n^{k+1} \in convex(\tilde{g}_n^k, \tilde{g}_{n+1}^k, \ldots)$ such that $g_n^{k+1}$ converges to some $g^{k+1}$. We denote its weights by $\lambda^{k+1,n}_n,\cdots,\lambda^{k+1,n}_{N^{k+1}_n}$ and can then write:

\begin{align*}
  g_n^{k+1} & = \sum_{m} \lambda^{k+1,n}_m \tilde{g}_m^k
  = \sum_{m} \lambda^{k+1,n}_m \left( \sum_{l} (\lambda^{k,m}_\cdot * \ldots * \lambda^{1,\cdot}_\cdot)_l x_l^{(k+1)} \right) \\
  & = \sum_{l} \left( \sum_{m} \lambda^{k+1,n}_m (\lambda^{k,m}_\cdot * \ldots * \lambda^{1,\cdot}_\cdot)_l \right) x_l^{(k+1)} \\
  & = \sum_{l} (\lambda^{k+1,n}_{\cdot} * (\lambda^{k,\cdot}_\cdot * \ldots * \lambda^{1,\cdot}_\cdot))_l x_l^{(k+1)}
\end{align*}

We have thus defined, for all $k, n \in \mathbb{N}$, convex weights $(\lambda^{k,n}_m)$ (that are zero for $m < n$) such that $\sum_{m \ge n}((\lambda^{k,n}_\cdot * \ldots * \lambda^{1, \cdot}_\cdot))_m x_m^{(k)}$ converges to $g^k$.
\end{proof}


\begin{lemma}\label{lem:komlos_convex_weights_tendsto}
  \lean{komlos_uniform_convergence}
  \uses{lem:komlos_convex_weights}
  \leanok
Let $E$ be a Hilbert space and for $i \in \mathbb{N}$, let $(x_n^{(i)})_{n \in \mathbb{N}}$ be a bounded sequence in $E$.
Let $(\lambda^{k,n}_\cdot)_{k, n \in \mathbb{N}}$ be convex weights satisfying the conclusion of Lemma~\ref{lem:komlos_convex_weights}, and let $(g^i)_{i\in \mathbb{N}}$ be the sequence of limits of the sums.

Then for every $k \ge i$, the sequence $\left(\sum_{m \ge n} \left((\lambda^{k,n}_\cdot) * \ldots * (\lambda^{1,\cdot}_\cdot)\right)_m x_m^{(i)}\right)_{n \in \mathbb{N}}$ converges to $g^i$, uniformly in $k$.
\end{lemma}

\begin{proof}  \uses{lem:convex_of_converg_seq_is_converg}
Let $i \in \mathbb{N}$.
By Lemma~\ref{lem:convex_of_converg_seq_is_converg}, there is uniform convergence over all convex combinations of the sequence $\left(\sum_{m \ge n} \left((\lambda^{i,n}_\cdot) * \ldots * (\lambda^{1,\cdot}_\cdot)\right)_m x_m^{(i)}\right)_{n \in \mathbb{N}}$ to $g^i$.
All sums $\sum_{m \ge n} \left((\lambda^{k,n}_\cdot) * \ldots * (\lambda^{1,\cdot}_\cdot)\right)_m x_m^{(i)}$ for $k \ge i$ are convex combinations of $\sum_{m \ge n} \left((\lambda^{i,n}_\cdot) * \ldots * (\lambda^{1,\cdot}_\cdot)\right)_m x_m^{(i)}$, hence they converge to $g^i$ uniformly in $k$.
\end{proof}


\begin{lemma}\label{lem:komlos_convex_weights_diagonal}
  \lean{komlos_convex_weights_diagonal}
  \leanok
Let $E$ be a Hilbert space and for $i \in \mathbb{N}$, let $(x_n^{(i)})_{n \in \mathbb{N}}$ be a bounded sequence in $E$.
Then there exists a sequence of convex weights $(\eta^n_\cdot)_{n \in \mathbb{N}}$ with $\eta^n_m = 0$ for $m < n$ such that for all $i \in \mathbb{N}$, the sequence $\left(\sum_{m \ge n} \eta^n_m x_m^{(i)}\right)_{n \in \mathbb{N}}$ converges.
\end{lemma}

\begin{proof}
  \uses{lem:komlos_convex_weights, lem:komlos_convex_weights_tendsto}
Let $(\lambda^{k,n}_\cdot)_{k, n \in \mathbb{N}}$ be convex weights satisfying the conclusion of Lemma~\ref{lem:komlos_convex_weights}, and let $(g^i)_{i\in \mathbb{N}}$ be the sequence of limits of the sums.
Let $\eta^n_m = (\lambda^{n,n}_\cdot * \ldots * \lambda^{1,\cdot}_\cdot)_m$.
We show that for all $i \in \mathbb{N}$, the sequence $\left(\sum_{m \ge n} \eta^n_m x_m^{(i)}\right)_{n \in \mathbb{N}}$ converges to $g^i$.

Let $i \in \mathbb{N}$.
By Lemma~\ref{lem:komlos_convex_weights_tendsto}, for all $\varepsilon > 0$, there exists $\bar{n}$ such that for all $n \ge \bar{n}$, for all $k \ge i$, $\left\Vert\sum_{m \ge n} \left((\lambda^{k,n}_\cdot) * \ldots * (\lambda^{1,\cdot}_\cdot)\right)_m x_m^{(i)} - g^i\right\Vert \le \varepsilon$.
Hence for $n \ge \max(\bar{n}, i)$,
\begin{align*}
  \left\Vert\sum_{m \ge n} \eta^n_m x_m^{(i)} - g^i\right\Vert
  &= \left\Vert\sum_{m \ge n} \left((\lambda^{n,n}_\cdot) * \ldots * (\lambda^{1,\cdot}_\cdot)\right)_m x_m^{(i)} - g^i\right\Vert
  \le \varepsilon
  \: .
\end{align*}
\end{proof}


\begin{lemma}\label{lem:komlos_convex_aux}
  \lean{komlos_convergence_L2}
  \leanok
Let $E$ be a Hilbert space and let $(f_n)_{n \in \mathbb{N}}$ be a sequence in $\Omega \to E$.
For $i \in \mathbb{N}$, set $f_n^{(i)} = f_n \mathbb{1}_{(\Vert f_n \Vert \le i)}$, such that $f_n^{(i)} \in L^2(E)$.
Then there exists a sequence of convex weights $\lambda_n^{n}, \ldots, \lambda_{N_n}^{n}$ such that the functions
$\left(\lambda_n^{n} f_n^{(i)} + \ldots + \lambda_{N_n}^{n} f_{N_n}^{(i)} \right)_{n\in\mathbb{N}}$
converge in $L^2(E)$ for every $i \in \mathbb{N}$.
\end{lemma}

\begin{proof}
  \uses{lem:komlos_convex_weights_diagonal}
Use Lemma~\ref{lem:komlos_convex_weights_diagonal} in the Hilbert space $L^2(E)$ with the sequence of sequences $(f_n^{(i)})$, which are bounded in $L^2(E)$ for each $i \in \mathbb{N}$.
\end{proof}


\begin{lemma}[Komlòs Lemma]\label{lem:komlos}
  \leanok
  \lean{komlos_L1}
Let $(f_n)_{n\in\mathbb{N}}$ be a uniformly integrable sequence of functions $\Omega \to E$, for $E$ a Hilbert space.
Then there exist functions $g_n \in convex(f_n, f_{n+1}, \cdots)$ such that $(g_n)_{n\in\mathbb{N}}$ converges in $L^1$.
\end{lemma}

\begin{proof}
  \uses{lem:komlos_convex_aux}
  For $i,n\in\mathbb{N}$ set $f_{n}^{(i)}:=f_n \mathbb{1}_{(|f_n|\leq i)}$ such that $f_{n}^{(i)}\in L^2$.
  Using \ref{lem:komlos_convex_aux} there exist for every $n$ convex weights $\lambda_n^{n}, \ldots, \lambda_{N_n}^{n}$ such that the functions
  $ \lambda_n^{n} f_n^{(i)} + \ldots+\lambda_{N_n}^{n} f_{N_n}^{(i)}$ converge in $L^2$ for every $i\in\mathbb{N}$.
  By uniform integrability, $\lim_{i\to \infty}\| f^{(i)}_n- f_n\|_1=0$, uniformly with respect to $n$.
  Hence, once again, uniformly with respect to $n$,
  $$ \textstyle\lim_{i\to\infty}\|  (\lambda_n^{n} f_n^{(i)} + \ldots+\lambda_{N_n}^{n} f_{N_n}^{(i)})-(\lambda_n^{n} f_n + \ldots+\lambda_{N_n}^{n} f_{N_n})\|_1= 0.$$
  Thus $(\lambda_n^{n} f_n + \ldots+\lambda_{N_n}^{n} f_{N_n})_{n\geq 1}$  is a Cauchy sequence in $L^1$.
\end{proof}


\paragraph{Komlòs lemma for nonnegative random variables}

\begin{lemma}[Komlòs lemma - nonnegative, a.e. convergence]\label{lem:komlos_ennreal}
  \leanok
  \lean{komlos_ennreal}
Let $(f_n)_{n\in\mathbb{N}}$ be a sequence of random variables with values in $[0, \infty]$.
Then there exist random variables $g_n \in convex( f_n, f_{n+1}, \cdots)$ such that $(g_n)_{n\in\mathbb{N}}$ converges almost surely to a random variable $g$.
\end{lemma}

\begin{proof}
  \uses{lem:komlos_convex}
Let $\phi : (\Omega \to [0, \infty]) \to [0, \infty]$ be defined by $\phi(X) = \mathbb{E}[e^{-X}]$.
Then $\phi$ is convex and $\phi(f_n) \le 1$ for all $n$.
By Lemma~\ref{lem:komlos_convex}, there exist $g_n \in convex( f_n, f_{n+1}, \cdots)$ such that for all $\delta>0$, for $N$ large enough and $n, m \ge N$,
\begin{align*}
  \mathbb{E}[e^{-g_n}]/2 + \mathbb{E}[e^{-g_m}]/2 - \mathbb{E}[e^{-(g_n + g_m)/2}] < \delta
  \: .
\end{align*}

For $\varepsilon > 0$, let $B_\varepsilon = \{(x, y) \in [0, \infty]^2 \mid \vert x - y \vert \ge \varepsilon \text{ and } \min\{x, y\} \le 1/\varepsilon \}$.
Then for all $x, y$,
\begin{align*}
  \left\vert e^{-x} - e^{-y} \right\vert
  &\le \varepsilon + 2 e^{-1/\varepsilon} + 2 \mathbb{1}_{B_\varepsilon}(x, y)
  \: .
\end{align*}
Hence for any pair of random variables $(X, Y)$ with values in $[0, \infty]$,
\begin{align*}
  \mathbb{E}\left[\left\vert e^{-X} - e^{-Y} \right\vert\right]
  &\le \varepsilon + 2 e^{-1/\varepsilon} + 2 P((X, Y) \in B_\varepsilon) \: .
\end{align*}
On the other hand, for $(x, y) \in B_\varepsilon$, there exists $\delta_\varepsilon > 0$ such that
\begin{align*}
  e^{-x}/2 + e^{-y}/2 - e^{-(x + y)/2} \ge \delta_\varepsilon
  \: .
\end{align*}
Thus,
\begin{align*}
  P((X, Y) \in B_\varepsilon)
  &\le \frac{1}{\delta_\varepsilon} \mathbb{E}\left[ e^{-X}/2 + e^{-Y}/2 - e^{-(X + Y)/2} \right]
  \: .
\end{align*}
For $n, m \ge N$ large enough so that we can apply the first inequality of this proof with $\delta = \varepsilon \delta_\varepsilon$, we deduce that
\begin{align*}
  \mathbb{E}\left[\left\vert e^{-g_n} - e^{-g_m} \right\vert\right]
  &\le \varepsilon + 2 e^{-1/\varepsilon} + \frac{2}{\delta_\varepsilon} \mathbb{E}\left[ e^{-g_n}/2 + e^{-g_m}/2 - e^{-(g_n + g_m)/2} \right] \\
  &\le \varepsilon + 2 e^{-1/\varepsilon} + 2 \varepsilon
  \: .
\end{align*}
As $\varepsilon$ is arbitrary, we deduce that $(e^{-g_n})_{n\in\mathbb{N}}$ is a Cauchy sequence in $L^1$ and thus converges in $L^1$ to some random variable $h$.
Therefore, it has a subsequence $(e^{-g_{n_k}})_{k\in\mathbb{N}}$ converging almost surely to $h$.
Finally, the subsequence of $g_n$ converges almost surely to $g = -\log(h)$.
\end{proof}





\section{Doob decomposition in discrete time}


\begin{definition}[Predictable part]\label{def:predictablePart}
  \uses{def:filtration}
  \mathlibok
  \lean{MeasureTheory.predictablePart}
Let $X : \mathbb{N} \to \Omega \to E$ be a process indexed by $\mathbb{N}$, for $E$ a Banach space.
Let $(\mathcal{F}_n)_{n\in\mathbb{N}}$ be a filtration on $\Omega$.
The predictable part of $X$ is the process $A : \mathbb{N} \to \Omega \to E$ defined for $n \ge 0$ by
$$A_n = \sum_{k=0}^{n-1} \mathbb{E}[X_{k+1}-X_k \mid \mathcal{F}_k].$$
\end{definition}


\begin{definition}[Martingale part]\label{def:martingalePart}
  \uses{def:predictablePart}
  \mathlibok
  \lean{MeasureTheory.martingalePart}
Let $X : \mathbb{N} \to \Omega \to E$ be a process indexed by $\mathbb{N}$, for $E$ a Banach space.
Let $(\mathcal{F}_n)_{n\in\mathbb{N}}$ be a filtration on $\Omega$ and let $A$ be the predictable part of $X$ for that filtration.
The martingale part of $X$ is the process $M : \mathbb{N} \to \Omega \to E$ defined by $M_n = X_n - A_n$.
\end{definition}


In what follows, we fix a process $X : \mathbb{N} \to \Omega \to E$ and a filtration $(\mathcal{F}_n)_{n \in \mathbb{N}}$, and denote by $A$ the predictable part of $X$ and by $M$ its martingale part.

\begin{lemma}\label{lem:predictablePart_zero}
  \uses{def:predictablePart}
  \mathlibok
  \lean{MeasureTheory.predictablePart_zero}
We have $A_0 = 0$.
\end{lemma}

\begin{proof}  \uses{def:predictablePart}
By definition.
\end{proof}


\begin{lemma}\label{lem:martingalePart_zero}
  \uses{def:martingalePart}
  \mathlibok
  \lean{MeasureTheory.martingalePart_zero}
$M_0 = X_0$.
\end{lemma}

\begin{proof}  \uses{def:martingalePart, lem:predictablePart_zero}
By definition of the martingale part, $M = X - A$. By Lemma~\ref{lem:predictablePart_zero}, $A_0 = 0$, thus $M_0 = X_0 - A_0 = X_0$.
\end{proof}


\begin{lemma}\label{lem:predictablePart_add_one}
  \uses{def:predictablePart}
  \mathlibok
  \lean{MeasureTheory.predictablePart_add_one}
For any integer $n \ge 0$, $A_{n+1} = A_n + \mathbb{E}[X_{n+1} - X_n \mid \mathcal{F}_n]$.
\end{lemma}

\begin{proof}  \uses{def:predictablePart}
Let $n \in \mathbb{N}$. Then
\begin{align*}
  A_{n+1}
  & = \sum_{k=0}^n \mathbb{E}[X_{k+1}-X_k \mid \mathcal{F}_k] \\
  & = \sum_{k=0}^{n-1} \mathbb{E}[X_{k+1}-X_k \mid \mathcal{F}_k] + \mathbb{E}[X_{n+1}-X_n \mid \mathcal{F}_n] \\
  & = A_n + \mathbb{E}[X_{n+1} - X_n \mid \mathcal{F}_n],
\end{align*}
which concludes the proof.
\end{proof}


\begin{lemma}\label{lem:martingalePart_add_one}
  \uses{def:martingalePart}
$M_{n+1} = M_n + X_{n+1} - X_n - \mathbb{E}[X_{n+1} - X_n \mid \mathcal{F}_n]$.
\end{lemma}

\begin{proof}
  \uses{def:martingalePart, lem:predictablePart_add_one}
Using Lemma~\ref{lem:predictablePart_add_one}, we have for $n \in \mathbb{N}$,
\begin{align*}
  M_{n+1}
  & = X_{n+1} - A_{n+1} \\
  & = X_{n+1} - A_n - \mathbb{E}[X_{n+1} - X_n \mid \mathcal{F}_n] \\
  & = M_n + X_{n+1} - X_n - \mathbb{E}[X_{n+1} - X_n \mid \mathcal{F}_n]
  \: .
\end{align*}
\end{proof}


\begin{lemma}\label{lem:Martingale.predictablePart_eq_zero}
  \uses{def:predictablePart, def:Martingale}
  \mathlibok
  \lean{MeasureTheory.Martingale.predictablePart_eq_zero}
If $X$ is a martingale, then $A = 0$ almost surely.
\end{lemma}

\begin{proof}  \uses{def:predictablePart, def:Martingale}
By the martingale property, each conditional expectation in the definition of $A$ is zero.
\end{proof}


\begin{lemma}\label{lem:isStronglyPredictable.predictablePart_eq}
  \uses{def:predictablePart, def:predictable}
  \mathlibok
  \lean{MeasureTheory.IsStronglyPredictable.predictablePart_eq}
If $X$ is predictable, then $A = X - X_0$ almost surely.
\end{lemma}

\begin{proof}  \uses{def:predictablePart, def:predictable}
Since $X$ is predictable, for all $n \in \mathbb{N}$, $X_{n+1}$ is $\mathcal{F}_n$-measurable and thus $\mathbb{E}[X_{n+1} - X_n \mid \mathcal{F}_n] = X_{n+1} - X_n$ a.s..
We get a telescoping sum in the definition of $A$ and thus $A_n = X_n - X_0$ a.s. for all $n \in \mathbb{N}$.
\end{proof}


\begin{lemma}\label{lem:martingalePart_eq_zero}
  \uses{def:martingalePart, def:predictable}
  \mathlibok
  \lean{MeasureTheory.IsPredictable.martingalePart_eq}
If $X$ is predictable, then $M = X_0$ almost surely.
\end{lemma}

\begin{proof}  \uses{lem:isStronglyPredictable.predictablePart_eq}
By definition of the martingale part, $M = X - A$. By Lemma~\ref{lem:isStronglyPredictable.predictablePart_eq}, $A = X - X_0$ almost surely, thus $M = X_0$ almost surely.
\end{proof}


\begin{lemma}\label{lem:Martingale.martingalePart_eq}
  \uses{def:martingalePart, def:Martingale}
  \mathlibok
  \lean{MeasureTheory.Martingale.martingalePart_eq}
If $X$ is a martingale, then $M = X$ almost surely.
\end{lemma}

\begin{proof}  \uses{lem:Martingale.predictablePart_eq_zero}
By definition of the martingale part, $M = X - A$. By Lemma~\ref{lem:Martingale.predictablePart_eq_zero}, $A = 0$ almost surely, thus $M = X$ almost surely.
\end{proof}


\begin{lemma}\label{lem:predictablePart_smul}
  \uses{def:predictablePart}
  \mathlibok
  \lean{MeasureTheory.predictablePart_smul}
For any scalar $c$, the predictable part of $c X$ is $c A$.
\end{lemma}

\begin{proof}  \uses{def:predictablePart}
Linearity of the conditional expectation.
\end{proof}


\begin{lemma}\label{lem:martingalePart_smul}
  \uses{def:martingalePart}
  \mathlibok
  \lean{MeasureTheory.martingalePart_smul}
For any scalar $c$, the martingale part of $c X$ is $c M$.
\end{lemma}

\begin{proof}  \uses{def:martingalePart, lem:predictablePart_smul}
By definition of the martingale part, $M = X - A$. By Lemma~\ref{lem:predictablePart_smul}, the predictable part of $c X$ is $c A$. Therefore, the martingale part of $c X$ is $c X - c A = c M$.
\end{proof}


\begin{lemma}\label{lem:predictablePart_add}
  \uses{def:predictablePart}
  \mathlibok
  \lean{MeasureTheory.predictablePart_add}
The predictable part of $X + Y$ is the sum of the predictable part of $X$ and the predictable part of $Y$.
\end{lemma}

\begin{proof}  \uses{def:predictablePart}
Linearity of the conditional expectation.
\end{proof}


\begin{lemma}\label{lem:martingalePart_add}
  \uses{def:martingalePart}
  \mathlibok
  \lean{MeasureTheory.martingalePart_add}
The martingale part of $X + Y$ is the sum of the martingale part of $X$ and the martingale part of $Y$.
\end{lemma}

\begin{proof}  \uses{def:martingalePart, lem:predictablePart_add}
By definition of the martingale part, $M = X - A$. By Lemma~\ref{lem:predictablePart_add}, the predictable part of $X + Y$ is $A + B$. Therefore, the martingale part of $X + Y$ is $(X + Y) - (A + B) = M + N$.
\end{proof}


\begin{lemma}\label{lem:adapted_predictablePart}
  \uses{def:predictablePart}
  \mathlibok
  \lean{MeasureTheory.stronglyAdapted_predictablePart}
The predictable part $A$ is adapted to the filtration $(\mathcal{F}_{n+1})_{n \in \mathbb{N}}$.
\end{lemma}

\begin{proof}\mathlibok

\end{proof}


\begin{lemma}\label{lem:predictable_predictablePart}
  \uses{def:predictablePart, def:predictable}
  \mathlibok
  \lean{MeasureTheory.isPredictable_predictablePart}
The predictable part of a process is predictable.
\end{lemma}

\begin{proof}  \uses{lem:predictable_nat_iff, lem:predictablePart_zero, lem:adapted_predictablePart}
By Lemma~\ref{lem:predictable_nat_iff}, the process $A$ is predictable if $A_0$ is $\mathcal{F}_0$-measurable and for all integer $n$, $A_{n+1}$ is $\mathcal{F}_n$-measurable. As $A_0 = 0$ from Lemma~\ref{lem:predictablePart_zero}, it is $\mathcal{F}_0$-measurable. Lemma~\ref{lem:adapted_predictablePart} allows to conclude the proof.
\end{proof}


\begin{lemma}\label{lem:martingale_martingalePart}
  \uses{def:martingalePart, def:Martingale}
  \mathlibok
  \lean{MeasureTheory.martingale_martingalePart}
Suppose that the filtration is $\sigma$-finite and that $X$ is adapted, with $X_n$ integrable for all $n$.
Then the martingale part of $X$ is a martingale.
\end{lemma}

\begin{proof}
\end{proof}


\begin{lemma}\label{lem:Submartingale.monotone_predictablePart}
  \uses{def:predictablePart, def:Submartingale}
  \mathlibok
  \lean{MeasureTheory.Submartingale.monotone_predictablePart}
The predictable part of a real-valued submartingale is an almost surely nondecreasing process.
\end{lemma}

\begin{proof}  \uses{def:Submartingale, lem:condExp_sub_nonneg, lem:predictablePart_add_one}
Let $X$ be a submartingale and let $A$ be its predictable part. Then for all $n \geq 0$, from Lemma~\ref{lem:condExp_sub_nonneg} we have that almost surely
\begin{align*}
  A_{n+1} &= A_n + \mathbb{E}\left[ X_{n+1} - X_n | \mathcal{F}_n \right]
  \ge A_n
  \: .
\end{align*}
The first equality comes from Lemma~\ref{lem:predictablePart_add_one}. As $\mathbb{N}$ is countable, we deduce that almost surely, for all $n \in \mathbb{N}$, $A_{n+1} \ge A_n$.
Thus, $(A_n)_{n \in \mathbb{N}}$ is almost surely nondecreasing.
\end{proof}


\begin{lemma}\label{lem:Submartingale.predictablePart_nonneg}
  \uses{def:predictablePart, def:Submartingale}
  \mathlibok
  \lean{MeasureTheory.Submartingale.predictablePart_nonneg}
The predictable part of a real-valued submartingale is almost surely nonnegative.
\end{lemma}

\begin{proof}  \uses{lem:Submartingale.monotone_predictablePart, lem:predictablePart_zero}
By Lemma~\ref{lem:Submartingale.monotone_predictablePart}, the predictable part $A$ is almost surely nondecreasing. By Lemma~\ref{lem:predictablePart_zero}, $A_0 = 0$. Therefore, $A_n \ge 0$ almost surely for all $n \in \mathbb{N}$.
\end{proof}


\begin{lemma}\label{lem:centering_basic}
  \uses{def:predictablePart, def:martingalePart}
  \mathlibok
Fake lemma: import this lemma when using the basic properties of the Doob decomposition.
\end{lemma}

\begin{proof}\leanok
  \uses{lem:predictablePart_zero, lem:martingalePart_zero, lem:predictablePart_add, lem:martingalePart_add, lem:predictablePart_smul, lem:martingalePart_smul, lem:Martingale.predictablePart_eq_zero, lem:isStronglyPredictable.predictablePart_eq, lem:martingalePart_eq_zero, lem:Martingale.martingalePart_eq, lem:Submartingale.monotone_predictablePart, lem:Submartingale.predictablePart_nonneg, lem:predictable_predictablePart, lem:martingale_martingalePart}

\end{proof}


\begin{lemma}\label{lem:IsStoppingTime_leastGT_predictablePart}
  \uses{def:IsStoppingTime,def:predictablePart, def:leastGT}
Let $X$ be a real adapted process and let $A$ be its predictable part. Let $c \in \mathbb{R}$.
The hitting time $\tau_{A_{n + 1} > c}$ of Definition~\ref{def:leastGT} is a stopping time.
\end{lemma}

\begin{proof}
  \uses{lem:predictable_predictablePart, lem:centering_basic}
Since $A_{n}$ is predictable, $A_{n + 1}$ is adapted.
The hitting time of an adapted process is a stopping time (we use the discrete time version of that result here, not the full Début theorem).
\end{proof}


\begin{lemma}\label{lem:leastGT_predictablePart_le}
  \uses{def:predictablePart, def:leastGT}
Let $X$ be a real adapted process and let $A$ be its predictable part. Let $c \in \mathbb{R}$.
Then $A_{\tau_{A_{n + 1} > c}} \le c$ and if $\tau_{A_{n + 1} > c} < \infty$ then $A_{\tau_{A_{n + 1} > c} + 1} > c$.
\end{lemma}

\begin{proof}
  \uses{lem:centering_basic}

\end{proof}


\begin{lemma}\label{lem:leastGT_predictablePart_sub_ge}
  \uses{def:predictablePart, def:leastGT}
Let $X$ be a real adapted process and let $A$ be its predictable part.
Let $a, b \in \mathbb{R}$ with $a \le b$.
If $\tau_{A_{n + 1} > b} < \infty$ then $A_{\tau_{A_{n + 1} > b}+1} - A_{\tau_{A_{n + 1} > a}} \ge b - a$.
\end{lemma}

\begin{proof}
  \uses{lem:leastGT_predictablePart_le}

\end{proof}


\begin{lemma}\label{lem:uniformIntegrable_predictablePart_aux1}
  \uses{def:predictablePart, def:martingalePart}
Let $T \in \mathbb{N}$ and let $X$ be an adapted process with $X_n$ integrable for all $n$ and such that $X_T = 0$.
Then for a stopping time $\tau \le T$, almost surely, $M_\tau = -\mathbb{E}[A_T \mid \mathcal{F}_\tau]$ and $X_\tau = A_\tau - \mathbb{E}[A_T \mid \mathcal{F}_\tau]$.
\end{lemma}

\begin{proof}
  \uses{def:martingalePart, lem:martingale_martingalePart, lem:centering_basic}
By definition and since $X_T = 0$, $M_T = X_T - A_T = - A_T$.
Since $M$ is a martingale it follows by optional sampling that for any stopping time $\tau \le T$
\begin{align*}
  M_\tau = \mathbb{E}[M_T \mid \mathcal{F}_\tau] = -\mathbb{E}[A_T \mid \mathcal{F}_\tau]
  \: .
\end{align*}
And then $X_\tau = M_\tau + A_\tau = A_\tau - \mathbb{E}[A_T \mid \mathcal{F}_\tau]$.
\end{proof}


TODO: define $\tau^T_{A_{n+1}>c}$ as the hitting time of $A_{n+1} > c$ but with the convention that it is equal to $T$ if the hitting time is greater than $T$ (\texttt{hittingBtwn}).

In the next lemmas, we write $\tau^T(c) = \tau^T_{A_{n+1}>c}$ for brevity.

\begin{lemma}\label{lem:predictablePart_gt_iff_leastGT_lt}
  \uses{def:predictablePart}
Suppose that $X$ is a submartingale and let $T \ge 1$, $c \in \mathbb{R}$.
Then $A_T > c \iff \tau^T(c) < T$.
\end{lemma}

\begin{proof}
  \uses{lem:leastGT_lt_iff, lem:Submartingale.monotone_predictablePart, lem:centering_basic}
By Lemma~\ref{lem:leastGT_lt_iff}, $\tau^T(c) < T \iff \exists s < T, A_{s+1} > c$, which is equivalent to $\exists s \in \{1, \ldots, T\}, A_s > c$.
By monotonicity of $A$ (Lemma~\ref{lem:Submartingale.monotone_predictablePart}), this is equivalent to $A_T > c$.
\end{proof}


\begin{lemma}\label{lem:uniformIntegrable_predictablePart_aux2}
  \uses{def:predictablePart, def:leastGT}
Suppose that $X$ is a submartingale with $X_T = 0$.
Then
\begin{align*}
  \mathbb{E}\left[A_T \mathbb{I}\{A_T > c\} \right]
  \le c \mathbb{P}(\tau^T(c) < T) - \int_{\tau^T(c) < T} X_{\tau^T(c)} dP
  \: .
\end{align*}
\end{lemma}

\begin{proof}
  \uses{lem:predictablePart_gt_iff_leastGT_lt, lem:leastGT_predictablePart_le, lem:uniformIntegrable_predictablePart_aux1}
By Lemma~\ref{lem:predictablePart_gt_iff_leastGT_lt},
\begin{align*}
  \mathbb{E}\left[A_T \mathbb{I}\{A_T > c\} \right]
  = \mathbb{E}\left[A_T \mathbb{I}\{\tau^T(c) < T\} \right]
\end{align*}
Since that last event is $\mathcal{F}_{\tau^T(c)}$-measurable, we can apply the tower property of conditional expectation to get
\begin{align*}
  \mathbb{E}\left[A_T \mathbb{I}\{\tau^T(c) < T\} \right]
  = \int_{\tau^T(c) < T} \mathbb{E}[A_T \mid \mathcal{F}_{\tau^T(c)}] dP
\end{align*}
Now by Lemma~\ref{lem:uniformIntegrable_predictablePart_aux1} and Lemma~\ref{lem:leastGT_predictablePart_le},
\begin{align*}
  \mathbb{E}\left[A_T \mathbb{I}\{\tau^T(c) < T\} \right]
  &= \int_{\tau^T(c) < T} A_{\tau^T(c)} - X_{\tau^T(c)} dP
  \\
  &\le c \mathbb{P}(\tau^T(c) < T) - \int_{\tau^T(c) < T} X_{\tau^T(c)} dP
  \: .
\end{align*}
\end{proof}


\begin{lemma}\label{lem:uniformIntegrable_predictablePart_aux3}
  \uses{def:predictablePart, def:leastGT}
Suppose that $X$ is a submartingale with $X_T = 0$.
Then
\begin{align*}
  \mathbb{P}(\tau^T(c) < T)
  &\le - \frac{2}{c} \int_{\tau^T(c/2) < T} X_{\tau^T(c/2)} dP
  \: .
\end{align*}
\end{lemma}

\begin{proof}
  \uses{lem:leastGT_predictablePart_sub_ge, lem:uniformIntegrable_predictablePart_aux1}
Notice that $\{\tau^T(c)<T\}\subseteq \{\tau^T(c/2)<T\}$, thus
\begin{align*}
  \int_{\tau^T(c/2)<T} -X_{\tau^T(c/2)}dP
  &=\int_{\tau^T(c/2)<T}\mathbb{E}[A_T \mid \mathcal{F}_{\tau^T(c/2)}] - A_{\tau^T(c/2)} dP
  \\
  &= \int_{\tau^T(c/2)<T}A_T - A_{\tau^T(c/2)}dP
  \\
  &\geq \int_{\tau^T(c)<T}A_T - A_{\tau^T(c/2)}dP
  \\
  \intertext{(over the event $\{\tau^T(c)<T\}$ $A_T\geq c$ and $A_{\tau^T(c/2)}\leq c/2$, thus $A_T - A_{\tau^T(c/2)}\geq c/2$)}
  &\ge \frac{c}{2}P(\tau^T(c)<T)
  \: .
\end{align*}
\end{proof}


\begin{lemma}\label{lem:uniformIntegrable_predictablePart_aux4}
  \uses{def:predictablePart, def:leastGT}
Suppose that $X$ is a submartingale with $X_T = 0$.
Then
\begin{align*}
  \mathbb{E}\left[A_T \mathbb{I}\{A_T > c\} \right]
  \le - 2 \int_{\tau^T(c/2) < T} X_{\tau^T(c/2)} dP - \int_{\tau^T(c) < T} X_{\tau^T(c)} dP
  \: .
\end{align*}
\end{lemma}

\begin{proof}
  \uses{lem:uniformIntegrable_predictablePart_aux2, lem:uniformIntegrable_predictablePart_aux3}
Put together the bounds of Lemma~\ref{lem:uniformIntegrable_predictablePart_aux2} and Lemma~\ref{lem:uniformIntegrable_predictablePart_aux3}.
\end{proof}


\begin{lemma}\label{lem:uniformIntegrable_predictablePart_aux5}
  \uses{def:predictablePart, def:leastGT}
Suppose that $X$ is a submartingale with $X_T = 0$.
Then
\begin{align*}
  P(\tau^T(c) < T) \le -\frac{\mathbb{E}[X_0]}{c} \: .
\end{align*}
\end{lemma}

\begin{proof}
  \uses{lem:predictablePart_gt_iff_leastGT_lt, lem:martingale_martingalePart}
Starting with Lemma~\ref{lem:predictablePart_gt_iff_leastGT_lt},
\begin{align*}
  P(\tau^T(c)<T)
  =P(A_T>c)
  \stackrel{Markov}{\leq}\frac{\mathbb{E}[A_T]}{c}
  =-\frac{\mathbb{E}[M_T]}{c}
  \stackrel{mg}{=}-\frac{\mathbb{E}[X_0]}{c}
  \: .
\end{align*}
\end{proof}





\section{Doob-Meyer decomposition}



For uniqueness of Doob-Meyer Decomposition we will need theorem
\ref{thm:IsLocalMartingale.eq_zero_of_predictable_finiteVariation}.

We now start the construction for the existence part.


\begin{definition}[Dyadics]\label{def:dyadics}
For $T>0$, let $\mathcal{D}_n^T = \left\lbrace \frac{k}{2^n}T \mid k=0,\cdots 2^n\right\rbrace$ be the set of dyadics at scale $n$ and let $\mathcal{D}^T=\bigcup_{n\in\mathbb{N}}\mathcal{D}_n^T$ be the set of all dyadics of $[0,T]$.
\end{definition}

Let $S : [0,T] \to \Omega \to \mathbb{R}$ be a cadlag submartingale of class D on $[0,T]$.



\begin{definition}[S, A, M]\label{def:SAM}
  \uses{def:dyadics, def:predictablePart}
For $n \in \mathbb{N}$, the restriction of $S$ to $\mathcal{D}_n^T$ is a discrete time submartingale $S^n : \mathbb{N} \to \Omega \to \mathbb{R}$ with respect to the filtration $\mathcal{F}^n_k = \mathcal{F}_{k2^{-n}T}$, with $S^n_k = S_{k2^{-n}T}$ (and constant equal to $S_T$ for $k > 2^n$; similarly for the filtration), and we can apply the construction of the previous section to it.
Let $A^n$ and $M^n$ be the predictable and martingale parts of $S^n$.
\end{definition}


\begin{lemma}\label{lem:A_uniform_integrable}
  \uses{def:SAM}
  The sequence $(A^n_{2^n})_{n\in\mathbb{N}}$ is uniformly integrable (bounded in $L^1$ norm).
\end{lemma}

Remark: $A^n_{2^n}$ is the predictable part of $S$ at time $T$ for the discrete time filtration given by the dyadics at scale $n$.

\begin{proof}
  \uses{lem:uniformIntegrable_predictablePart_aux4, lem:uniformIntegrable_predictablePart_aux5, def:classD, lem:predictablePart_add, lem:martingalePart_add, lem:Martingale.predictablePart_eq_zero}
WLOG $S^n_{2^n} = S_T=0$ and $S_t\leq 0$ (else consider $S_t-\mathbb{E}\left[S_T\vert\mathcal{F}_{t}\right]$).

We write $\tau_n(c)$ for the hitting time $\tau^T_{A^n_{k+1}>c}$.

By Lemma~\ref{lem:uniformIntegrable_predictablePart_aux4},
$$
\int_{(A^n_{2^n}>c)} A^n_{2^n} dP
\le -2 \int_{\tau_n(c/2)< 2^n} S^n_{\tau_n(c/2)} dP - \int_{\tau_n(c) < 2^n} S^n_{\tau_n(c)} dP.
$$
On the other hand, by Lemma~\ref{lem:uniformIntegrable_predictablePart_aux5},
$$
P(\tau_n(c)<2^n) \le -\frac{\mathbb{E}[S_0]}{c}
$$
which goes to $0$ uniformly in $n$ as $c$ goes to infinity.

The integrals in the upper bound on $\int_{(A^n_{2^n}>c)} A^n_{2^n} dP$ are integrals of uniformly integrable random variables (by the class D assumption) over sets whose probability goes to zero uniformly.
Therefore, these integrals go to zero uniformly in $n$ as $c$ goes to infinity.

This implies that $\int_{(A^n_{2^n}>c)} A^n_{2^n} dP$ goes to 0 uniformly in $n$ as $c \to +\infty$.
Hence, the $L^1$ norm is uniformly bounded.
\end{proof}


\begin{lemma}\label{lem:M_uniform_integrable}
  \uses{def:SAM}
  The sequence $(M^n_{2^n})_{n\in\mathbb{N}}$ is uniformly integrable (bounded in $L^1$ norm).
\end{lemma}

\begin{proof}
  \uses{lem:A_uniform_integrable, lem:ClassD.uniformIntegrable}
$M^n_{2^n} = S_{2^n} - A^n_{2^n}$, also $S$ is of class $D$ hence uniformly integrable and $A^n_{2^n}$ is uniformly integrable by Lemma~\ref{lem:A_uniform_integrable}.
\end{proof}


\begin{lemma}\label{lem:M_n_cadlag_mg}
  \uses{def:SAM}
The martingale on $[0, T]$ defined by $t \mapsto \mathbb{E}[M^n_{2^n}\vert\mathcal{F}_t]$ admits a modification which is a cadlag martingale.
\end{lemma}

\begin{proof}
  \uses{thm:Martingale.exists_cadlag_modification}
By theorem \ref{thm:Martingale.exists_cadlag_modification}
\end{proof}


\begin{definition}\label{def:barM}
  \uses{lem:M_n_cadlag_mg}
For $t\in[0,T]$ let $\overline{M}^n_t$ be the cadlag modification of $t \mapsto \mathbb{E}[M^n_{2^n} \mid \mathcal{F}_t]$ from lemma \ref{lem:M_n_cadlag_mg}.
\end{definition}


\begin{lemma}\label{lem:M_cal_converges_L1}
  \uses{def:barM}
  There exists an $M : \Omega \to \mathbb{R}$ and convex weights $\lambda^n_n,\cdots,\lambda^n_{N_n}$ such that
  $\mathcal{M}^n_T\stackrel{L^1}{\rightarrow}M$, where $\mathcal{M}^n_T := \lambda^n_n \overline{M}^n_T+ \cdots + \lambda^n_{N_n} \overline{M}^{N_n}_T$~.
\end{lemma}

\begin{proof}
  \uses{lem:M_uniform_integrable,lem:komlos}
  By lemma \ref{lem:M_uniform_integrable} $(M^n_T)_{n\in\mathbb{N}}$ is uniformly bounded in $L^1$, thus by lemma \ref{lem:komlos} there are convex weights $\lambda^n_n,\cdots,\lambda^n_{N_n}$ such that
  $\mathcal{M}^n_T\stackrel{L^1}{\rightarrow}M$, where $\mathcal{M}^n_T := \lambda^n_n \overline{M}^n_T+ \cdots + \lambda^n_{N_n} \overline{M}^{N_n}_T$~.
\end{proof}


\begin{definition}\label{def:calM}
  \uses{lem:M_cal_converges_L1, def:barM}
For $t\in[0,T]$ let $\mathcal{M}^n_t=\lambda^n_n \overline{M}^n_t+\cdots +\lambda^n_{N_n} \overline{M}^{N_n}_t$ be the convex combination of the $\overline{M}^n_t$'s with the weights from Lemma~\ref{lem:M_cal_converges_L1} and let $M$ be the limit of $\mathcal{M}^n_T$ from the same lemma.
\end{definition}


\begin{lemma}\label{lem:M_cal_cadlag}
  \uses{def:calM}
  $\mathcal{M}^n$ is cadlag.
\end{lemma}

\begin{proof}
  \uses{lem:M_n_cadlag_mg,lem:M_cal_converges_L1}
  By construction and \ref{lem:M_n_cadlag_mg}
\end{proof}


\begin{definition}\label{def:M}
  \uses{def:calM}
Let $M_t$ be a cadlag modification of $\mathbb{E}[M \mid \mathcal{F}_t]$, obtained by applying Theorem~\ref{thm:Martingale.exists_cadlag_modification} to the martingale $t \mapsto \mathbb{E}[M \mid \mathcal{F}_t]$.
\end{definition}


\begin{lemma}\label{lem:M1_komlos}
  \uses{def:calM, def:M}
  For every $t\in[0,T]$ we have $\mathcal{M}^n_t\stackrel{L^1}{\rightarrow}M_t$.
\end{lemma}

\begin{proof}
  \uses{lem:M_cal_converges_L1}
By Jensen's inequality, the tower lemma and lemma \ref{lem:M_cal_converges_L1}
\begin{gather}\nonumber
  \mathbb{E}[|\mathcal{M}^n_t-M_t|]
  = \mathbb{E}[|\mathbb{E}[\mathcal{M}^n_T-M\vert\mathcal{F}_t]|]
  \le \mathbb{E}[|\mathcal{M}^n_T-M|]\rightarrow0 \: ,
  \\
  \Rightarrow \mathcal{M}^n_t\stackrel{L^1}{\rightarrow} M_t,\quad \forall t\in[0,T].\label{equation_DM_e7}
\end{gather}
\end{proof}


\begin{definition}\label{def:barA}
  \uses{def:SAM}
For $t\in[0,T]$ let $\overline{A}^n_t$ be process defined by
\begin{align*}
  \overline{A}^n_s
  := \sum_{m < 2^n} A^n_{m+1} \mathbb{1}_{(m2^{-n}T, (m+1)2^{-n}T]}(s)
  \: .
\end{align*}

\end{definition}


\begin{lemma}\label{lem:barA_eq}
  \uses{def:barA, def:SAM}
For $s \in \mathcal{D}^T_n$, $\overline{A}^n_s = A^n_m$ where $m$ is such that $s = m2^{-n}T$.
\end{lemma}

\begin{proof}

\end{proof}


\begin{lemma}\label{lem:leftContinuous_barA}
  \uses{def:barA}
$\overline{A}^n$ is left continuous.
\end{lemma}

\begin{proof}
  \uses{def:barA}

\end{proof}


\begin{lemma}\label{lem:predictable_barA}
  \uses{def:barA, def:predictable}
$\overline{A}^n$ is predictable.
\end{lemma}

\begin{proof}
  \uses{lem:leftContinuous_barA, def:barA, lem:Adapted.isPredictable_of_leftContinuous}
Since $\overline{A}^n$ is left continuous and adapted, it is predictable (Lemma~\ref{lem:Adapted.isPredictable_of_leftContinuous}).
\end{proof}


\begin{definition}\label{def:calA}
  \uses{def:barA, def:calM}
Let $\mathcal{A}^n := \lambda^n_n \overline{A}^n+\cdots +\lambda^n_{N_n}\overline{A}^{N_n}$, in which the $\lambda^n_n,\cdots,\lambda^n_{N_n}$ are the convex weights from Lemma~\ref{lem:M_cal_converges_L1}.
\end{definition}


\begin{lemma}\label{lem:leftContinuous_calA}
  \uses{def:calA}
$\mathcal{A}^n$ is left continuous.
\end{lemma}

\begin{proof}
  \uses{def:calA, lem:leftContinuous_barA}

\end{proof}


\begin{lemma}\label{lem:predictable_calA}
  \uses{def:calA, def:predictable}
$\mathcal{A}^n$ is predictable.
\end{lemma}

\begin{proof}
  \uses{lem:leftContinuous_calA, def:calA, lem:predictable_barA}
Either use that it is left continuous and adapted, or that it is a convex combination of predictable processes.
\end{proof}


\begin{lemma}\label{lem:calA_mono}
  \uses{def:calA}
$\mathcal{A}^n$ is non-decreasing on $[0,T]$.
\end{lemma}

\begin{proof}
  \uses{lem:Submartingale.monotone_predictablePart, def:calA}

\end{proof}


\begin{definition}\label{def:A}
  \uses{def:M}
Let $A_t = S_t - M_t$ for $t \in [0, T]$.
\end{definition}


\begin{lemma}\label{lem:cadlag_A}
  \uses{def:A}
$A$ is cadlag.
\end{lemma}

\begin{proof}
  \uses{def:A, def:M}
Since $S$ and $M$ are cadlag, their difference $A = S - M$ is also cadlag.
\end{proof}


\begin{lemma}\label{lem:tendsto_calA_A_L1}
  \uses{def:calA, def:A}
For every $t\in \mathcal{D}^T$, $\mathcal{A}^n_t$ converges to $A_t$ in $L^1$.
\end{lemma}

\begin{proof}
  \uses{lem:M1_komlos}
By Lemma \ref{lem:M1_komlos}, for all $t\in\mathcal{D}^T$,
\begin{align*}
  \mathcal{A}^n_t
  = S_t - \mathcal{M}^n_t
  \stackrel{L^1}{\rightarrow} S_t-M_t
  = A_t
  \: .
\end{align*}
\end{proof}


\begin{lemma}\label{lem:tendsto_calA_A_ae}
  \uses{def:calA, def:A}
For every $t\in \mathcal{D}^T$, there exists a subsequence $(k_n)$ such that $\mathcal{A}^{k_n}_t$ converges to $A_t$ almost surely.
\end{lemma}

\begin{proof}
    \uses{lem:tendsto_calA_A_L1}
By Lemma \ref{lem:tendsto_calA_A_L1}, we have convergence in $L^1$ for every $t\in \mathcal{D}^T$, which implies that there exists a subsequence $(k_n)$ such that $\mathcal{A}^{k_n}_t$ converges to $A_t$ almost surely.
\end{proof}

\begin{lemma}\label{lem:A_mono_dyadics}
  \uses{def:A}
$A$ is almost surely non-decreasing on $\mathcal{D}^T$.
\end{lemma}

\begin{proof}
  \uses{lem:tendsto_calA_A_ae, lem:calA_mono}
Since $\mathcal{D}^T$ is countable, to prove that $A$ is almost surely non-decreasing on $\mathcal{D}^T$ it is enough to show that for all $s \le t$ in $\mathcal{D}^T$, $A_s \le A_t$ almost surely.

Let $s \le t$ be two elements of $\mathcal{D}^T$.
By Lemma \ref{lem:tendsto_calA_A_ae}, there exists a subsequence $(k_n)$ such that $\mathcal{A}^{k_n}_t$ converges to $A_t$ almost surely.
The subsequence $\mathcal{A}_s^{k_n}$ converges to $A_s$ in $L^1$ and thus it has a further subsequence $(k'_n)$ that converges almost surely for $s$ and $t$.
By Lemma \ref{lem:calA_mono}, $\mathcal{A}^{k'_n}_s \le \mathcal{A}^{k'_n}_t$ for all $n$, thus by taking the limit we get $A_s \le A_t$ almost surely.
\end{proof}


% \begin{lemma}\label{lem:A_cal_conv_A_on_D_T}
%   \uses{def:calA, def:A}
% There exists a sequence $k_n$ such that almost surely, $\lim_n \mathcal{A}^{k_n}_t(\omega) = A_t(\omega)$ for every $t\in\mathcal{D}^T$.
% \end{lemma}

% \begin{proof}
%   \uses{lem:tendsto_calA_A_L1}
% By Lemma \ref{lem:tendsto_calA_A_L1}, for all $t\in\mathcal{D}^T$,
% \begin{align*}
%   \mathcal{A}^n_t
%   \stackrel{L^1}{\rightarrow} A_t
%   \: .
% \end{align*}
% Since $\mathcal{D}^T$ is countable we can arrange the elements as $(t_n)_{n\in\mathbb{N}}$.
% Given $t_0\in\mathcal{D}^T$ there exists a subsequence $k^{0}_n$ for which $\mathcal{A}^{k^{0}_n}_{t_0}$ converges to $A_{t_0}$ over the set $\Omega\setminus E_{0}$ where $P(E_{0})=0$.
% Suppose we have a sequence $k^m_n$ for which $\mathcal{A}^{k^j_n}_{t_j}$ converges to $A_{t_j}$ over the set $\Omega\setminus E_{m}$ where $P(E_{m})=0$ for each $j=0,\cdots,m$.
% From this subsequence we may extract a new subsequence $k^{m+1}_n$ for which $\mathcal{A}^{k^{m+1}_n}_{t_{m+1}}$ converges to $A_{t_{m+1}}$ over the set $\Omega\setminus E_{m+1}$ where $P(E_{m+1})=0$.
% By construction over this subsequence the convergence for $t_0,\cdots,t_m$ still applies.
% With a diagonal argument we obtain the final result with $E=\bigcup_n E_n$.
% \end{proof}


\begin{lemma}\label{lem:A_increasing}
  \uses{def:A}
  $(A_t)_{t\in[0,T]}$ is almost surely non-decreasing.
\end{lemma}

\begin{proof}
  \uses{lem:A_mono_dyadics, lem:cadlag_A}
$A$ is almost surely non-decreasing on $\mathcal{D}^T$ by Lemma~\ref{lem:A_mono_dyadics}.
Since $A$ is cadlag (thus right-continuous) by Lemma~\ref{lem:cadlag_A}, it follows that $A$ must be non-decreasing on $[0,T]$.
\end{proof}


\begin{lemma}\label{lem:jump_A}
  \uses{def:A}
For $q,k \in \mathbb{N}$, let $\tau_{q, k}$ be the $q$-th time that the process $A_t$ has a jump higher than $1/(k+1)$ (this is a hitting time).
Then if $A(\omega)$ is discontinuous at $t$, then there exists $q,k$ such that $\tau_{q, k}(\omega) = t$.
\end{lemma}

\begin{proof}
  \uses{lem:A_increasing}
Since $A$ is non-decreasing it has finitely many jumps of size higher than $1/(k+1)$ for each $k$.
Thus, if $A$ is discontinuous at $t$, there exists $k$ such that the jump of $A$ at $t$ is higher than $1/(k+1)$, and thus there exists $q$ such that $\tau_{q, k} = t$.
\end{proof}


\begin{lemma}\label{lem:incr_fun_lim_right_cont_limsup_ineq}
  If $f_n, f : [0, 1] \rightarrow \mathbb{R}$ are increasing functions such that $f$ is right continuous and
  $\lim_n f_n(t) = f (t)$ for $t \in\mathcal{D}^T$, then  $\limsup_n  f_n(t) \leq f (t)$ for all $t \in [0, T]$.
\end{lemma}

\begin{proof}
  Let $t\in[0,T]$ and $s\in\mathcal{D}^T$ such that $t<s$. We have
  $$
  \limsup_n f_n(t)\leq \limsup_n f_n(s)=f(s).
  $$
  Since the above is true uniformly in $s$ in particular since $f$ is right-continuous
  $$
  \limsup_n f_n(t)\leq\lim_{\stackrel{s\rightarrow t^+}{s\in\mathcal{D}^T}}f(s)=f(t).
  $$
\end{proof}


\begin{lemma}\label{lem:incr_fun_lim_right_cont_lim_eq}
  If $f_n, f : [0, 1] \rightarrow \mathbb{R}$ are increasing functions such that $f$ is right continuous and
  $\lim_n f_n(t) = f (t)$ for $t \in\mathcal{D^T}$, if $f$ is continuous in $t\in[0,T]$ then $\lim_n  f_n(t) = f (t)$.
\end{lemma}

\begin{proof}
  \uses{lem:incr_fun_lim_right_cont_limsup_ineq}
  By lemma \ref{lem:incr_fun_lim_right_cont_limsup_ineq} it is enough to show that $\liminf_n f_n(t)\geq f(t)$.
  Let $s\in\mathcal{D}^T$ such that $t>s$. We have
  $$
  \liminf_n f_n(t)\geq \liminf_n f_n(s)=f(s).
  $$
  Since the above is true uniformly in $s$ in particular since $f$ is continuous in $t$
  $$
  \liminf_n f_n(t)\geq\lim_{\stackrel{s\rightarrow t^-}{s\in\mathcal{D}^T}}f(s)=f(t).
$$
\end{proof}


\begin{lemma}\label{lem:lim_Exp_A_n_tau_is_Exp_A_tau}
  Let $\tau$ be an $(\mathcal{F}_t)_{t\in[0,T]}$ stopping time. We have $\lim_n\mathbb{E}[A^n_\tau]=\mathbb{E}[A_\tau]$.
\end{lemma}

\begin{proof}
  \uses{lem:M_n_cadlag_mg, def:classD}
  Let $\sigma_n:=\inf\left(t\in\mathcal{D}^T_n\vert t>\tau\right)$. By construction of $A^n$ we have $A^n_\tau=A^n_{\sigma_n}$.
  Also $\sigma_n\searrow\tau$. Since $S$ is of class $D$ and cadlag we have
  \begin{align*}
    \mathbb{E}[A^n_\tau]&=\mathbb{E}[A^n_{\sigma_n}]=\mathbb{E}[S_{\sigma_n}]-\mathbb{E}[M^n_{\sigma_n}]=\mathbb{E}[S_{\sigma_n}]-\mathbb{E}[M^n_0]=\\
    &=\mathbb{E}[S_{\sigma_n}]-\mathbb{E}[S_0]\rightarrow \mathbb{E}[S_\tau]-\mathbb{E}[M_0]=\mathbb{E}[S_\tau]-\mathbb{E}[M_\tau]=\mathbb{E}[A_\tau].
  \end{align*}
\end{proof}


\begin{lemma}\label{lem:limsup_A_n_tau_is_A_tau_ae}
  Let $\tau$ be an $(\mathcal{F}_t)_{t\in[0,T]}$ stopping time. We have $\limsup_n \mathcal{A}_\tau^n = A_\tau$.
\end{lemma}

\begin{proof}
  \uses{lem:lim_Exp_A_n_tau_is_Exp_A_tau, lem:incr_fun_lim_right_cont_limsup_ineq, lem:Submartingale.monotone_predictablePart}
Firstly we notice that $\liminf_n \mathbb{E}[A_\tau^n]  \leq \limsup_n  \mathbb{E}  [\mathcal{A}_\tau^n  ]  \leq \mathbb{E}[\limsup_n  \mathcal{A}_\tau^n  ]  \leq \mathbb{E}[ A_\tau ]$,
where the first inequality is justified by the definition of limsup and liminf and the fact that
$$
\sup_{k\geq n}\mathbb{E}[\mathcal{A}^k_\tau]\geq \sum_{m=k}^{N_k}\lambda^k_m\mathbb{E}[A^m_\tau]\geq \sum_{m=k}^{N_k}\lambda^k_m\inf_{j\geq n}\mathbb{E}[A^j_\tau]=\inf_{k\geq n}\mathbb{E}[A^k_\tau]
$$
the third inequality by \ref{lem:incr_fun_lim_right_cont_limsup_ineq}.
Let's prove the second inequality: observe that
$$
\mathcal{A}^n_\tau= A_1+\mathcal{A}^n_\tau-A_1\leq A_1+(\mathcal{A}^n_\tau-A_1)_+,
$$
thus it follows that $\mathcal{A}^n_\tau - (\mathcal{A}^n_\tau-A_1)_+\leq A_1$; since $A_1$ is an integrable guardian the inverse Fatou Lemma may be applied to show together with limsup properties that
\begin{align*}
  \limsup_n\mathbb{E}[\mathcal{A}^n_\tau]+0 &= \limsup_n\mathbb{E}[\mathcal{A}^n_\tau]+\liminf_n-\mathbb{E}[(\mathcal{A}^n_\tau-A_1)_+] \leq \limsup_n\mathbb{E}[\mathcal{A}^n_\tau-(\mathcal{A}^n_\tau-A_1)_+]\leq\\
  &\leq \mathbb{E}[\limsup_n\mathcal{A}^n_\tau-(\mathcal{A}^n_\tau-A_1)_+]\leq \mathbb{E}[\limsup_n\mathcal{A}^n_\tau]-\mathbb{E}[\liminf_n(\mathcal{A}^n_\tau-A_1)_+]\leq\mathbb{E}[\limsup_n\mathcal{A}^n_\tau],
  \end{align*}
where the first equality is justified by the fact that $\mathcal{A}^n_\tau\leq\mathcal{A}^n_1\rightarrow A_1$ almost surely.
Due to lemma \ref{lem:lim_Exp_A_n_tau_is_Exp_A_tau} and \ref{lem:incr_fun_lim_right_cont_limsup_ineq} the first sequence of inequalities is a sequence of equalities, thus
we know that $A_\tau- \limsup_n \mathcal{A}_\tau^n $ is an a.s. nonnegative function with null expected value, and thus it must be almost everywhere null.
\end{proof}


\begin{theorem}[Doob-Meyer decomposition]\label{thm:Doob_Meyer}
  Let $S = (S_t )_{0\leq t\leq T}$ be a cadlag submartingale of class $D$.
  Then, $S$ can be written in a unique way in the form  $S = M + A$ where $M$ is a cadlag martingale and $A$ is a predictable increasing process starting at $0$.
\end{theorem}

\begin{proof}
  \uses{lem:A_increasing, lem:limsup_A_n_tau_is_A_tau_ae, lem:incr_fun_lim_right_cont_lim_eq, lem:Adapted.isPredictable_of_leftContinuous}
  By construction $M$ is a cadlag martingale and $A_0=0$ and by lemma \ref{lem:A_increasing} $A$ is increasing. It suffices to show that $A$ is predictable.
  $A^n,\mathcal{A}^n$ are left continuous and adapted, and thus they are predictable (Lemma~\ref{lem:Adapted.isPredictable_of_leftContinuous}).
  It is enough to show that $\omega-a.e.$, $\forall t\in[0,T]$, $\limsup_n\mathcal{A}^n_t(\omega)=A_t(\omega)$.

  By lemma \ref{lem:incr_fun_lim_right_cont_lim_eq} that is true for any continuity point of $A$. Since $A$ is increasing it can only have a finite amount of jumps larger than $1/k$ for any $k\in\mathbb{N}$.
  Consider now $\tau_{q,k}$ the family of stopping times equal to the $q$-th time that the process $A_t$ has a jump higher than $1/k$. This is a countable family.
  Given a time $t$ and a trajectory $\omega$ there are only two possibilities: either $A$ is continuous or not at time $t$ along $\omega$.
  If $A$ is continuous at time $t$ we have $\limsup_n\mathcal{A}^n_t(\omega)=A_t(\omega)$, if it jumps there exists $q(\omega),k(\omega)$ such that $t=\tau_{q(\omega),k(\omega)}(\omega)$.
  Due to lemma \ref{lem:limsup_A_n_tau_is_A_tau_ae} we know that $\limsup_n A^n_{\tau_{q,k}} = A_{\tau_{q,k}}$ for each $q,k$ almost surely. Thus, since it is an intersection of a countable amount of almost sure
  events $\forall\omega\in\Omega'$ with $P(\Omega')=1$, for each $q,k$ $\limsup_n A^n_{\tau_{q,k}}(\omega) = A_{\tau_{q,k}}(\omega)$ ($\omega$ does not depend upon $q,k$).
  Consequently, $\forall\omega\in\Omega'$ we have $\limsup_n\mathcal{A}^n_t(\omega)=\limsup_n\mathcal{A}^n_{\tau_{q(\omega),k(\omega)}}(\omega)=A_{\tau_{q(\omega),k(\omega)}}(\omega)=A_t(\omega)$
\end{proof}



\section{Local version of the Doob-Meyer decomposition}


\begin{theorem}[Doob-Meyer decomposition]\label{thm:local_doobMeyer}
  \uses{def:IsLocalSubmartingale, def:predictable, def:IsLocalMartingale}
  \leanok
  \lean{ProbabilityTheory.IsLocalSubmartingale.doob_meyer}
An adapted process $X$ is a cadlag local submartingale iff $X = M + A$ where $M$ is a cadlag local martingale and $A$ is a predictable, cadlag, locally integrable and increasing process starting at $0$.
The processes $M$ and $A$ are uniquely determined by $X$ a.s.
\end{theorem}

\begin{proof}
  \uses{lem:IsLocalSubmartingale.locally_classD, thm:Doob_Meyer}


\leanok
\end{proof}


\begin{lemma}[Adding a bottom-measurable constant to a martingale]
  \label{lem:Martingale.add_const_fun}
  \lean{MeasureTheory.Martingale.add_const_fun}
Let \(Z\) be an integrable \(\mathcal F_0\)-measurable random variable.
Assume the filtration is sigma-finite for the measure \(P\).  If \(M\)
is a martingale with respect to \(\mathcal F\) and \(P\), then the shifted
process \(t \mapsto M_t+Z\) is again a martingale.
\end{lemma}

\begin{proof}
The time-constant process \(t\mapsto Z\) is a martingale under the
sigma-finite-filtration hypothesis: adaptedness is exactly the
\(\mathcal F_0\)-measurability of \(Z\) transported along
\(\mathcal F_0\subseteq\mathcal F_t\), integrability is the assumed
integrability of \(Z\), and its conditional expectations are unchanged.
Adding this constant martingale to \(M\) preserves the martingale
property.
\end{proof}

\begin{lemma}\leanok
[Running supremum after subtracting a time-constant random variable]
  \label{lem:runningSup_norm_sub_const_le}
  \lean{ProbabilityTheory.runningSup_norm_sub_const_le}
For any process \(A\), random variable \(Z\), time \(t\), and sample
point \(\omega\),
\[
  \sup_{s\le t}\|A_s(\omega)-Z(\omega)\|_e
  \le
  \sup_{s\le t}\|A_s(\omega)\|_e+\|Z(\omega)\|_e .
\]
\end{lemma}

\begin{proof}

\leanok
For each \(s\le t\), the triangle inequality gives
\[
  \|A_s(\omega)-Z(\omega)\|_e
  \le \|A_s(\omega)\|_e+\|Z(\omega)\|_e .
\]
Taking the supremum over all \(s\le t\) on the left and bounding each
\(\|A_s(\omega)\|_e\) by the running supremum of \(A\) gives the claimed
inequality.
\end{proof}

\begin{lemma}\leanok
[Integrable running supremum after subtracting a bottom-measurable random variable]
  \label{lem:HasIntegrableSup.sub_const_fun}
  \uses{lem:runningSup_norm_sub_const_le}
  \lean{ProbabilityTheory.HasIntegrableSup.sub_const_fun}
Assume the time index and filtration satisfy the usual hypotheses needed
to turn strong progressivity into strong measurability of the running
supremum process.  Let \(A\) have integrable running supremum, and let
\(Z\) be an \(\mathcal F_0\)-measurable random variable such that
\(\|Z\|_e\) is integrable.  If \(A\) is strongly progressive, then
\(t\mapsto A_t-Z\) also has integrable running supremum.
\end{lemma}

\begin{proof}

\leanok
The constant process \(t\mapsto Z\) is strongly progressive because
\(\mathcal F_0\subseteq\mathcal F_t\) for every \(t\).  Hence
\(A-Z\) is strongly progressive.  Under the usual hypotheses, strong
progressivity gives strong measurability of the running-supremum process
for \(A-Z\).

For each \(t\) and \(\omega\), Lemma~\ref{lem:runningSup_norm_sub_const_le}
gives
\[
  \sup_{s\le t}\|A_s(\omega)-Z(\omega)\|_e
  \le
  \sup_{s\le t}\|A_s(\omega)\|_e+\|Z(\omega)\|_e .
\]
The right-hand side is integrable by the integrable running supremum of
\(A\) and the assumed integrability of \(\|Z\|_e\).  Domination therefore
gives integrability of the running supremum of \(A-Z\).
\end{proof}

\begin{lemma}\leanok
[Locally integrable running supremum after subtracting the initial value]
  \label{lem:HasLocallyIntegrableSup.sub_initial}
  \uses{lem:HasIntegrableSup.sub_const_fun}
  \lean{ProbabilityTheory.HasLocallyIntegrableSup.sub_initial}
Assume the same usual hypotheses as
Lemma~\ref{lem:HasIntegrableSup.sub_const_fun}.  If \(A\) is strongly
progressive and has locally integrable running supremum, then
\(t\mapsto A_t-A_0\) has locally integrable running supremum.
\end{lemma}

\begin{proof}

\leanok
Let \((\tau_n)\) be a localizing sequence witnessing that \(A\) has locally
integrable running supremum.  For each \(n\), set
\[
  Z_n(\omega)=\mathbf 1_{\{0<\tau_n(\omega)\}} A_0(\omega).
\]
Strong progressivity of \(A\) implies that \(Z_n\) is
\(\mathcal F_0\)-measurable.  Its norm is integrable because, at time
\(0\), it is exactly the running supremum of the stopped localized process
for \(A\).

Apply Lemma~\ref{lem:HasIntegrableSup.sub_const_fun} to the \(n\)-th
stopped process and \(Z_n\).  Pointwise, the stopped version of
\(A-A_0\) is the stopped process for \(A\) shifted by this \(Z_n\), after
splitting according to whether \(0<\tau_n(\omega)\).  Therefore the same
localizing sequence witnesses locally integrable running supremum for
\(A-A_0\).
\end{proof}

\begin{lemma}

\leanok
[Adding the initial value of a locally integrable process to a local martingale]
  \label{lem:IsLocalMartingale.add_initial_of_hasLocallyIntegrableSup}
  \uses{lem:Martingale.add_const_fun}
  \lean{ProbabilityTheory.IsLocalMartingale.add_initial_of_hasLocallyIntegrableSup}
Assume the time index is second countable, Borel, and pseudo-metrizable,
the measure is finite, and the filtration admits the approximating
sequences needed for optional stopping.  If \(M\) is a local martingale
and \(A\) is strongly progressive with locally integrable running
supremum, then \(t\mapsto M_t+A_0\) is a local martingale.
\end{lemma}

\begin{proof}

\leanok
Choose a localizing sequence for \(M\) and a localizing sequence for the
locally integrable running supremum of \(A\), and localize by their
pointwise minimum.  For the \(n\)-th stopped process let
\[
  S_n=\{0<\tau_n\}, \qquad
  Z_n(\omega)=\mathbf 1_{S_n}(\omega)A_0(\omega).
\]
Stability of martingales under stopping by the second localizing sequence
shows that the stopped and truncated version of \(M\) is a martingale.
The random variable \(Z_n\) is \(\mathcal F_0\)-measurable because \(A\)
is strongly progressive, and it is integrable because the stopped
localized running supremum of \(A\) is integrable at time \(0\).

By Lemma~\ref{lem:Martingale.add_const_fun}, adding the time-constant
process \(Z_n\) preserves the martingale property for this stopped
process.  The stopped process for \(M+A_0\) is pointwise equal to this
sum, after splitting according to whether \(0<\tau_n(\omega)\).  Hence
the minimum localizing sequence witnesses that \(M+A_0\) is a local
martingale.
\end{proof}

\begin{theorem}\leanok
[Normalized local Doob-Meyer decomposition]\label{thm:local_doobMeyer_normalized}
  \uses{thm:local_doobMeyer, lem:IsLocalMartingale.add_initial_of_hasLocallyIntegrableSup, lem:HasLocallyIntegrableSup.sub_initial}
  \lean{ProbabilityTheory.IsLocalSubmartingale.doob_meyer_normalized}
Assume the usual finite-measure and time-index hypotheses needed by the
two normalization closure lemmas: the local-martingale shift hypotheses
from Lemma~\ref{lem:IsLocalMartingale.add_initial_of_hasLocallyIntegrableSup}
and the running-supremum measurability hypotheses from
Lemma~\ref{lem:HasLocallyIntegrableSup.sub_initial}.  If \(X\) is a
càdlàg local submartingale, then there exist processes \(M\) and \(A\)
such that \(X=M+A\), \(M\) is a càdlàg local martingale, and \(A\) is
strongly predictable, strongly progressive, càdlàg, locally integrable,
nondecreasing, and satisfies \(A_0=0\).
\end{theorem}

\begin{proof}

\leanok
Apply the local Doob--Meyer decomposition to get \(X=M+A\), where \(M\)
is a càdlàg local martingale and \(A\) is strongly predictable, hence
strongly progressive, càdlàg, locally integrable, and nondecreasing.
Define the time-constant process \(C_t(\omega)=A_0(\omega)\), then set
\[
  \widetilde A_t=A_t-C_t,
  \qquad
  \widetilde M_t=M_t+C_t .
\]
Then \(X=\widetilde M+\widetilde A\) by pointwise algebra and
\(\widetilde A_0=0\).  Subtracting the same random variable from every
time does not change monotonicity of the paths, and subtracting a
time-constant càdlàg path from a càdlàg path is again càdlàg.  The
time-constant process \(C\) is strongly predictable because it is
\(\mathcal F_0\)-measurable, and the predictable sigma-algebra contains
sets of the form \(T\times B\) for \(B\in\mathcal F_0\).  Therefore
\(\widetilde A=A-C\) is strongly predictable.  Its strong progressivity is
then either inherited from predictability or obtained by subtracting the
corresponding strongly progressive processes.

It remains to justify two closure facts.  First, \(\widetilde M=M+C\) is a
local martingale by
Lemma~\ref{lem:IsLocalMartingale.add_initial_of_hasLocallyIntegrableSup},
applied to \(M\), \(A\), strong progressivity of \(A\), and the locally
integrable running supremum of \(A\).

Second, \(\widetilde A=A-C\) has locally integrable running supremum by
Lemma~\ref{lem:HasLocallyIntegrableSup.sub_initial}.  This is the
formal version of using the same localizing sequence as for \(A\), bounding
the shifted running supremum by
Lemma~\ref{lem:runningSup_norm_sub_const_le}, and transporting strong
measurability of the running-supremum process through subtraction of the
time-constant initial value.  Under the usual hypotheses stated in the
theorem, these two closure facts are available in the ambient context, and
the normalized decomposition follows.
\end{proof}


\begin{lemma}\label{lem:isStronglyPredictable_predictablePart}
  \uses{thm:local_doobMeyer_normalized}
  \lean{ProbabilityTheory.IsLocalSubmartingale.isStronglyPredictable_predictablePart}
  \leanok
For the normalized choice of Doob--Meyer decomposition, the predictable
part is strongly predictable.
\end{lemma}

\begin{proof}

\leanok
Define the choice-based martingale and predictable parts using the
normalized existence theorem.  The assertion is the strongly predictable
component of the chosen witness.
\end{proof}


\begin{lemma}\label{lem:predictablePart_bot_eq_zero}
  \uses{thm:local_doobMeyer_normalized}
  \lean{ProbabilityTheory.IsLocalSubmartingale.predictablePart_bot_eq_zero}
  \leanok
For the normalized choice of Doob--Meyer decomposition, the predictable
part satisfies \(A_0=0\).
\end{lemma}

\begin{proof}

\leanok
Define the choice-based martingale and predictable parts using the
normalized existence theorem.  The assertion is then the final component
of the chosen witness.
\end{proof}


\begin{lemma}\label{lem:IsStronglyPredictable_stoppedProcess_indicator}
  \uses{def:IsStoppingTime}
  \lean{MeasureTheory.IsStronglyPredictable.stoppedProcess_indicator}
Let \(U\) be a strongly predictable real-valued process and let \(\tau\) be
a stopping time.  Then the stopped process
\[
  (t,\omega)\longmapsto
  \mathbf 1_{\{\tau(\omega)>0\}}\,U_{t\wedge\tau(\omega)}(\omega)
\]
is strongly predictable.
\end{lemma}

\begin{proof}
It is enough to show that the map
\[
  \Phi(t,\omega)=(t\wedge\tau(\omega),\omega)
\]
is measurable from the predictable sigma-algebra to itself and that
indicator truncation by \(\{\tau>0\}\) preserves predictable measurability.
For the stopping map, check the generators of the predictable
sigma-algebra.  The preimage of \((i,\infty)\times A\), with
\(A\in\mathcal F_i\), is
\[
  (i,\infty)\times \bigl(A\cap\{i<\tau\}\bigr),
\]
which is predictable because \(\{i<\tau\}\in\mathcal F_i\).  The preimage of
\(\{0\}\times A\), with \(A\in\mathcal F_0\), is the union of
\(\{0\}\times A\) and \(\iota\times(A\cap\{\tau=0\})\); this is predictable
because \(\{\tau=0\}\in\mathcal F_0\) and rectangles over
\(\mathcal F_0\) are predictable.  Composing the predictable version of
\(U\) with \(\Phi\) gives predictability of \(U^\tau\).  Finally,
\(\{\tau>0\}\in\mathcal F_0\), so multiplying by its indicator is
multiplication by the predictable set \(\iota\times\{\tau>0\}\).
\end{proof}


\begin{lemma}\label{lem:LocallyBoundedVariationOn_stoppedProcess_indicator}
  \uses{def:LocallyBoundedVariationOn}
  \lean{ProbabilityTheory.locallyBoundedVariationOn_stoppedProcess_indicator}
  \leanok
Let \(f:\iota\to\mathbb R\) have locally bounded variation on all compact
intervals, and let \(a\in\iota_\infty\).  Then the path
\[
  t\longmapsto
  \mathbf 1_{\{0<a\}}\,f(t\wedge a)
\]
also has locally bounded variation on all compact intervals.
\end{lemma}

\begin{proof}

\leanok
If \(a\le 0\), the stopped-indicator path is identically zero.  Otherwise,
on a compact interval \([r,s]\) the stopped path is obtained from \(f\) by
reading \(f\) along the monotone map \(t\mapsto t\wedge a\), and then
remaining constant after \(a\).  Every partition of \([r,s]\) is sent to an
ordered finite list of points in \([r,s]\), with possible repetitions.
Deleting repetitions can only remove zero increments.  Hence the variation
of the stopped path on \([r,s]\) is bounded by the variation of \(f\) on
\([r,s]\), which is finite by the locally bounded-variation hypothesis.
\end{proof}


\begin{lemma}\label{lem:stoppedProcess_indicator_bound_on_Icc}
  \uses{def:stoppedProcess}
  \lean{MeasureTheory.stoppedProcess_indicator_bound_on_Icc}
Let \(N\) be a real-valued process, let \(\tau\) be a stopping time candidate
with values in \(\iota_\infty\), and fix a deterministic horizon
\([\bot,t]\).  Suppose that a sample path is bounded on this horizon by a
nonnegative deterministic constant \(C\):
\[
  \|N_s(\omega)\|\le C\qquad(s\in[\bot,t]).
\]
Then the stopped/indicator path
\[
  s\longmapsto
  \mathbf 1_{\{\bot<\tau(\omega)\}}\,
  N_{s\wedge\tau(\omega)}(\omega)
\]
is also bounded by \(C\) on \([\bot,t]\).
\end{lemma}

\begin{proof}
Fix \(s\in[\bot,t]\).  If \(\bot<\tau(\omega)\), then
\((s\wedge\tau(\omega))\) lies between \(\bot\) and \(t\): the lower bound
uses both \(\bot\le s\) and \(\bot<\tau(\omega)\), while the upper bound
comes from \(s\le t\).  The original path bound applies at this stopped
time.  If \(\bot<\tau(\omega)\) fails, the indicator is zero, so the
conclusion is \(0\le C\).
\end{proof}


\begin{lemma}\label{lem:stoppedProcess_indicator_bound_on_Icc_of_pre_stop_bound}
  \uses{def:stoppedProcess}
  \lean{MeasureTheory.stoppedProcess_indicator_bound_on_Icc_of_pre_stop_bound}
Let \(N\) be a real-valued process and let \(\tau\) be an
\(\iota_\infty\)-valued cutoff.  Fix a deterministic horizon \([\bot,t]\)
and a constant \(C\ge0\).  Suppose that for a fixed sample point \(\omega\),
the original path is bounded by \(C\) at every horizon time strictly before
\(\tau(\omega)\), and that the stop value is also bounded by \(C\) whenever
the indicator branch is active and \(\tau(\omega)\) is finite and lies at or
before \(t\):
\[
  r\in[\bot,t],\quad r<\tau(\omega)
  \quad\Longrightarrow\quad
  \|N_r(\omega)\|\le C,
\]
and
\[
  \bot<\tau(\omega),\quad \tau(\omega)\ne\top,\quad
  \tau(\omega)\le (t:\iota_\infty)
  \quad\Longrightarrow\quad
  \|N_{\tau(\omega)}(\omega)\|\le C .
\]
Then the stopped/indicator path
\[
  s\longmapsto
  \mathbf 1_{\{\bot<\tau(\omega)\}}\,
  N_{s\wedge\tau(\omega)}(\omega)
\]
is bounded by \(C\) on \([\bot,t]\).
\end{lemma}

\begin{proof}
Fix \(s\in[\bot,t]\).  If \(\bot<\tau(\omega)\) is false, then the indicator
is zero and the conclusion follows from \(0\le C\).

Assume \(\bot<\tau(\omega)\).  If \(s<\tau(\omega)\), stopping has not yet
occurred at \(s\), the indicator is one, and the pre-stop bound applies to
\(s\).  If \(s\ge\tau(\omega)\), the stopped coordinate is \(\tau(\omega)\).
Since \(s\le t\), this gives \(\tau(\omega)\le t\); since
\(\tau(\omega)\le s<\top\), the stopping value is finite.  The stop-value
bound therefore applies.  The equality case \(s=\tau(\omega)\) is covered by
the same stop-value bound.  These cases prove the bound for every
\(s\in[\bot,t]\).
\end{proof}

\begin{lemma}\label{lem:ae_stoppedProcess_indicator_bound_on_Icc_of_pre_stop_bound}
  \uses{lem:stoppedProcess_indicator_bound_on_Icc_of_pre_stop_bound}
  \lean{MeasureTheory.ae_stoppedProcess_indicator_bound_on_Icc_of_pre_stop_bound}
Let \(N\), \(\tau\), \(t\), and \(C\ge0\) be as in the preceding lemma.  If
the pre-stop bound and the stop-value bound hold almost surely, then the
stopped/indicator process is bounded by \(C\) on \([\bot,t]\) almost surely:
\[
  \forall s\in[\bot,t],\quad
  \bigl\|
    \mathbf 1_{\{\bot<\tau\}}N_{s\wedge\tau}
  \bigr\|\le C
  \quad\text{a.s.}
\]
\end{lemma}

\begin{proof}
Intersect the almost-sure event carrying the pre-stop bound with the
almost-sure event carrying the stop-value bound.  On a sample point in this
intersection, apply
Lemma~\ref{lem:stoppedProcess_indicator_bound_on_Icc_of_pre_stop_bound}.
This gives the stopped/indicator horizon bound pointwise on that full-measure
event, which is exactly the desired almost-sure conclusion.
\end{proof}

\begin{lemma}\label{lem:leastGT_norm_pre_stop_bound_on_Icc}
  \uses{def:leastGT}
  \lean{MeasureTheory.leastGT_norm_pre_stop_bound_on_Icc}
Let \(N\) be a real-valued process and fix a sample point \(\omega\), a
deterministic horizon \([\bot,t]\), and a level \(C\).  Put
\[
  \tau_C(\omega)=\inf\{s:\|N_s(\omega)\|>C\}.
\]
Then before this first strict crossing, the path is bounded by \(C\): if
\(r\in[\bot,t]\) and \(r<\tau_C(\omega)\), then
\[
  \|N_r(\omega)\|\le C .
\]
\end{lemma}

\begin{proof}
The definition of \(\tau_C\) is the hitting time `leastGT` of the open set
\((C,\infty)\) by the process \(s\mapsto\|N_s\|\), started at \(\bot\).
If \(r<\tau_C(\omega)\), then \(r\) is not a hitting time witness; otherwise
the hitting time would be at most \(r\).  Hence it is not true that
\(C<\|N_r(\omega)\|\), and linearity of the real order gives
\(\|N_r(\omega)\|\le C\).
\end{proof}

\begin{lemma}\label{lem:leastGT_norm_stop_value_bound_on_Icc_of_continuousOn}
  \uses{def:leastGT, lem:leastGT_lt_iff, lem:Filter.nhdsWithin_Iio_self_neBot_of_bot_lt}
  \lean{MeasureTheory.leastGT_norm_stop_value_bound_on_Icc_of_continuousOn}
Assume now that the time index is densely ordered with the order topology.
Let \(N(\cdot,\omega)\) be continuous on \([\bot,t]\), and put
\(\tau_C(\omega)=\inf\{s:\|N_s(\omega)\|>C\}\).  If
\(\bot<\tau_C(\omega)\), \(\tau_C(\omega)\ne\top\), and
\(\tau_C(\omega)\le t\), then the finite stopped value does not overshoot:
\[
  \|N_{\tau_C(\omega)}(\omega)\|\le C .
\]
\end{lemma}

\begin{proof}
Let \(a=\tau_C(\omega)\), viewed as a finite time.  Suppose, for contradiction,
that \(C<\|N_a(\omega)\|\).  By continuity of the path on \([\bot,t]\), the
strict inequality holds throughout a sufficiently small neighborhood of \(a\)
inside the horizon.  Since the order is dense and \(\bot<a\),
Lemma~\ref{lem:Filter.nhdsWithin_Iio_self_neBot_of_bot_lt} gives a point
\(s<a\) in this neighborhood.  Then \(C<\|N_s(\omega)\|\), so
Lemma~\ref{lem:leastGT_lt_iff} implies \(\tau_C(\omega)<a\), contradicting
the equality \(a=\tau_C(\omega)\).  Therefore the stopped value is at most
\(C\).
\end{proof}

\begin{lemma}\label{lem:ae_leastGT_norm_pre_stop_stop_value_bound_on_Icc_of_continuousOn}
  \uses{lem:leastGT_norm_pre_stop_bound_on_Icc, lem:leastGT_norm_stop_value_bound_on_Icc_of_continuousOn}
  \lean{MeasureTheory.ae_leastGT_norm_pre_stop_stop_value_bound_on_Icc_of_continuousOn}
Let \(N\) be a real-valued process with every sample path continuous on
\([\bot,t]\), and let
\[
  \tau_C(\omega)=\inf\{s:\|N_s(\omega)\|>C\}.
\]
Then the two bound hypotheses needed by
Lemma~\ref{lem:ae_stoppedProcess_indicator_bound_on_Icc_of_pre_stop_bound}
hold almost surely for this stopping level: before \(\tau_C\) the path is
bounded by \(C\), and at any active finite stopped value lying before \(t\)
the stopped value is also bounded by \(C\).
\end{lemma}

\begin{proof}
Both conclusions are pointwise in \(\omega\).  The pre-stop bound is
Lemma~\ref{lem:leastGT_norm_pre_stop_bound_on_Icc}.  The stop-value bound is
Lemma~\ref{lem:leastGT_norm_stop_value_bound_on_Icc_of_continuousOn}, applied
to the pathwise continuity hypothesis at the same sample point.  Turning these
pointwise implications into almost-sure statements uses the full-measure event
containing every sample point.
\end{proof}


\begin{lemma}\label{lem:eVariationOn_stoppedProcess_indicator_le_Icc}
  \uses{def:stoppedProcess}
  \lean{BoundedVariationOn.eVariationOn_stoppedProcess_indicator_le_Icc}
For a real-valued path \(f\), a level \(a\in\iota_\infty\), and a
deterministic ordered interval \([\bot,t]\), the total variation of the stopped/indicator
path
\[
  s\longmapsto \mathbf 1_{\{\bot<a\}} f(s\wedge a)
\]
on \([\bot,t]\) is at most the total variation of \(f\) on \([\bot,t]\).
\end{lemma}

\begin{proof}
If \(\bot<a\) is false, the stopped/indicator path is identically zero and
the inequality is immediate.  Otherwise every monotone finite partition of
\([\bot,t]\) is sent by \(s\mapsto s\wedge a\) to a monotone finite list of
points still lying in \([\bot,t]\).  The sum of stopped increments along the
original partition is therefore bounded by the variation sum of \(f\) along
that image partition.  Taking the supremum over partitions gives the stated
inequality.
\end{proof}


\begin{lemma}\label{lem:stoppedProcess_indicator_variation_bound_on_Icc}
  \uses{def:stoppedProcess, lem:eVariationOn_stoppedProcess_indicator_le_Icc}
  \lean{BoundedVariationOn.stoppedProcess_indicator_variation_bound_on_Icc}
Let \(N\) be a real-valued process and assume that, for a fixed sample point
\(\omega\), \(N(\cdot,\omega)\) has bounded variation on \([\bot,t]\) with
total variation at most \(V\).  Then the stopped/indicator path
\[
  s\longmapsto
  \mathbf 1_{\{\bot<\tau(\omega)\}}\,
  N_{s\wedge\tau(\omega)}(\omega)
\]
also has bounded variation on \([\bot,t]\), and its total variation on this
interval is at most \(V\).
\end{lemma}

\begin{proof}
Apply
Lemma~\ref{lem:eVariationOn_stoppedProcess_indicator_le_Icc}.  Since the
original variation on \([\bot,t]\) is finite, the stopped variation is finite
as well.  The same inequality, followed by monotonicity of
\(\mathrm{toReal}\) for finite extended nonnegative reals, transfers the
deterministic real bound \(V\).
\end{proof}

\begin{lemma}\label{lem:eVariationOn_stoppedProcess_indicator_le_closed_pre_stop_Icc}
  \uses{def:stoppedProcess}
  \lean{BoundedVariationOn.eVariationOn_stoppedProcess_indicator_le_closed_pre_stop_Icc}
Fix a sample point \(\omega\), a cutoff \(\tau(\omega)\in\iota_\infty\),
and a deterministic horizon \([\bot,t]\).  Let
\[
  S_{\tau,t}(\omega)
  =
  \{r\in[\bot,t] : (r:\iota_\infty)\le \tau(\omega)\}
\]
be the closed pre-stop part of the horizon, including the finite stopping
value when it lies before \(t\).  Then the total variation on \([\bot,t]\)
of the stopped/indicator path
\[
  s\longmapsto
  \mathbf 1_{\{\bot<\tau(\omega)\}}\,
  N_{s\wedge\tau(\omega)}(\omega)
\]
is at most the total variation of the original path on
\(S_{\tau,t}(\omega)\).
\end{lemma}

\begin{proof}
If \(\bot<\tau(\omega)\) is false, the indicator is zero and the stopped path
has zero variation, so the claim follows from nonnegativity of variation.

Assume \(\bot<\tau(\omega)\).  Given a monotone finite partition of
\([\bot,t]\), map every partition time \(s\) to
\[
  g(s)=(s\wedge\tau(\omega))_{\mathrm{finite}}.
\]
The map \(g\) is monotone.  For \(s\in[\bot,t]\), the point \(g(s)\) still
belongs to \([\bot,t]\), because it is no larger than \(s\le t\) and no
smaller than \(\bot\); moreover \((g(s):\iota_\infty)\le\tau(\omega)\).
Thus the image partition lies in \(S_{\tau,t}(\omega)\).  On the active
indicator branch the stopped/indicator increment along the original
partition is exactly the original-path increment along the image partition.
The variation sum of the stopped path is therefore bounded by the variation
of the original path on \(S_{\tau,t}(\omega)\).  Taking the supremum over all
partitions gives the result.
\end{proof}

\begin{lemma}\label{lem:stoppedProcess_indicator_variation_bound_on_Icc_of_closed_pre_stop_bound}
  \uses{def:stoppedProcess, lem:eVariationOn_stoppedProcess_indicator_le_closed_pre_stop_Icc}
  \lean{BoundedVariationOn.stoppedProcess_indicator_variation_bound_on_Icc_of_closed_pre_stop_bound}
Assume that, for a fixed sample point \(\omega\), the original path has
bounded variation on the closed pre-stop horizon set
\[
  S_{\tau,t}(\omega)
  =
  \{r\in[\bot,t] : (r:\iota_\infty)\le \tau(\omega)\}
\]
and that this variation is at most \(V\).  Then the stopped/indicator path
has bounded variation on \([\bot,t]\), and its variation on \([\bot,t]\) is
also at most \(V\).
\end{lemma}

\begin{proof}
Apply
Lemma~\ref{lem:eVariationOn_stoppedProcess_indicator_le_closed_pre_stop_Icc}.
Finiteness of the original variation on \(S_{\tau,t}(\omega)\) implies
finiteness of the stopped variation.  Monotonicity of
\(\mathrm{toReal}\) for finite extended nonnegative reals then transfers the
real bound \(V\).
\end{proof}

\begin{lemma}\label{lem:ae_stoppedProcess_indicator_bound_variation_on_Icc_of_pre_stop_bound}
  \uses{lem:ae_stoppedProcess_indicator_bound_on_Icc_of_pre_stop_bound, lem:stoppedProcess_indicator_variation_bound_on_Icc_of_closed_pre_stop_bound}
  \lean{MeasureTheory.ae_stoppedProcess_indicator_bound_variation_on_Icc_of_pre_stop_bound}
Let \(N\), \(\tau\), \(t\), \(C\), and \(V\) be fixed.  Suppose \(C\ge0\).
Assume the pre-stop horizon bound and finite stop-value bound from
Lemma~\ref{lem:ae_stoppedProcess_indicator_bound_on_Icc_of_pre_stop_bound}
hold almost surely.  Also assume that almost surely the original path has
bounded variation at most \(V\) on the closed pre-stop horizon set
\[
  S_{\tau,t}(\omega)
  =
  \{r\in[\bot,t] : (r:\iota_\infty)\le \tau(\omega)\}.
\]
Then the stopped/indicator process has, almost surely, both the horizon bound
by \(C\) and the variation bound by \(V\) on \([\bot,t]\).
\end{lemma}

\begin{proof}
For the horizon bound, apply
Lemma~\ref{lem:ae_stoppedProcess_indicator_bound_on_Icc_of_pre_stop_bound}
to the almost-sure pre-stop and stop-value bounds.  For the variation bound,
intersect the almost-sure event carrying the closed pre-stop variation input
and apply
Lemma~\ref{lem:stoppedProcess_indicator_variation_bound_on_Icc_of_closed_pre_stop_bound}
pointwise.  The two almost-sure conclusions form exactly the desired pair of
stopped/indicator bounds.
\end{proof}

\begin{lemma}\label{lem:ae_stoppedProcess_indicator_bound_variation_on_Icc}
  \uses{lem:stoppedProcess_indicator_bound_on_Icc, lem:stoppedProcess_indicator_variation_bound_on_Icc}
  \lean{MeasureTheory.ae_stoppedProcess_indicator_bound_variation_on_Icc}
Let \(N\) be a real-valued process, let \(\tau\) be an
\(\iota_\infty\)-valued time, and fix a deterministic horizon \([\bot,t]\).
Assume that \(C\ge0\), that \(N\) is bounded by \(C\) on \([\bot,t]\) almost
surely, and that \(N\) has bounded variation on \([\bot,t]\) with total
variation at most \(V\) almost surely.  Then the stopped/indicator process
\[
  s\longmapsto
  \mathbf 1_{\{\bot<\tau(\omega)\}}\,
  N_{s\wedge\tau(\omega)}(\omega)
\]
has the corresponding almost-sure horizon bound by \(C\) and almost-sure
variation bound by \(V\) on \([\bot,t]\).
\end{lemma}

\begin{proof}
Intersect the two almost-sure events carrying the original horizon bound and
the original variation bound.  On a sample point in the first event, apply
Lemma~\ref{lem:stoppedProcess_indicator_bound_on_Icc}.  On a sample point in
the second event, apply
Lemma~\ref{lem:stoppedProcess_indicator_variation_bound_on_Icc}.  These two
pointwise conclusions are exactly the two almost-sure stopped/indicator
hypotheses.
\end{proof}

\begin{lemma}\label{lem:stoppedProcess_indicator_continuousOn_Icc}
  \lean{MeasureTheory.stoppedProcess_indicator_continuousOn_Icc}
Let \(N\) be a real-valued process, let \(\tau\) be an
\(\iota_\infty\)-valued time, fix a sample point \(\omega\), and fix a
deterministic horizon \([\bot,t]\).  If the path
\(s\mapsto N_s(\omega)\) is continuous on \([\bot,t]\), then the
stopped/indicator path
\[
  s\longmapsto
  \mathbf 1_{\{\bot<\tau(\omega)\}}\,
  N_{s\wedge\tau(\omega)}(\omega)
\]
is continuous on \([\bot,t]\).
\end{lemma}

\begin{proof}
Split according to whether \(\bot<\tau(\omega)\).  If not, the indicator is
identically zero, so the stopped/indicator path is continuous.

Assume \(\bot<\tau(\omega)\).  If \(\tau(\omega)=\top\), then stopping does
nothing on the deterministic horizon and the stopped path is just the
original continuous path.  Otherwise write \(\tau(\omega)=a\) with
\(\bot<a\).  On \([\bot,t]\), the stopped time coordinate is the continuous
order map \(s\mapsto \min(s,a)\).  This map sends \([\bot,t]\) into itself
when \(a\ge t\), and sends \([\bot,t]\) into \([\bot,a]\subseteq[\bot,t]\)
when \(a\le t\).  Composing the assumed continuity of
\(s\mapsto N_s(\omega)\) on \([\bot,t]\) with this stopped-coordinate map
therefore gives continuity of the stopped path.  The indicator is the
constant value \(1\) on this branch, so it does not affect continuity.
\end{proof}


\begin{lemma}\label{lem:Filtration.measurable_prod_mk_predictable_past}
  \lean{MeasureTheory.Filtration.measurable_prod_mk_predictable_past}
Let \((\mathcal F_s)_{s\in\iota}\) be a filtration on \(\Omega\).  If
\(t>0\), then the deterministic section map
\[
  \Omega\longrightarrow \iota\times\Omega,\qquad
  \omega\longmapsto (t,\omega)
\]
from \(\Omega\) equipped with the sigma-algebra
\(\bigvee_{s<t}\mathcal F_s\) to \(\iota\times\Omega\) equipped with the
predictable sigma-algebra is measurable.
\end{lemma}

\begin{proof}
Unfold the predictable sigma-algebra as the sigma-algebra generated by the
sets \(\{0\}\times A\), with \(A\in\mathcal F_0\), and
\((r,\infty)\times A\), with \(A\in\mathcal F_r\).  It is enough to check
the inverse image of each generator under \(\omega\mapsto(t,\omega)\).
Since \(t>0\), the section of \(\{0\}\times A\) is empty.  For
\((r,\infty)\times A\), the section is empty if \(t\le r\), and is exactly
\(A\) if \(r<t\).  In the nonempty case \(A\in\mathcal F_r\), and
\(\mathcal F_r\) is one of the sigma-algebras included in
\(\bigvee_{s<t}\mathcal F_s\).  Thus every generator has a measurable
section, proving the claimed measurability.
\end{proof}


\begin{lemma}\label{lem:IsStronglyPredictable.stronglyMeasurable_past}
  \uses{lem:Filtration.measurable_prod_mk_predictable_past}
  \lean{MeasureTheory.IsStronglyPredictable.stronglyMeasurable_past}
Let \(X\) be a strongly predictable real-valued process.  For every
deterministic time \(t>0\), the random variable \(X_t\) is strongly
measurable with respect to the past sigma-algebra
\(\bigvee_{s<t}\mathcal F_s\).
\end{lemma}

\begin{proof}
Strong predictability says that the uncurried map
\((s,\omega)\mapsto X_s(\omega)\) is strongly measurable for the predictable
sigma-algebra on \(\iota\times\Omega\).  Compose this map with the
deterministic section \(\omega\mapsto(t,\omega)\).  By
Lemma~\ref{lem:Filtration.measurable_prod_mk_predictable_past}, the section
is measurable from \(\bigvee_{s<t}\mathcal F_s\) into the predictable
sigma-algebra.  The composition is precisely \(\omega\mapsto X_t(\omega)\),
so \(X_t\) is strongly measurable with respect to the past sigma-algebra.
\end{proof}


\begin{theorem}\label{thm:Martingale.eq_zero_of_predictable_finiteVariation_discrete}
  \uses{def:Martingale}
  \lean{MeasureTheory.Martingale.eq_zero_of_predictable_finiteVariation_discrete}
Let \(M\) be a martingale indexed by \(\mathbb N\), with \(M_0=0\).
Assume that \(M\) is strongly predictable.  Then \(M_n=0\) almost surely
for every \(n\).
\end{theorem}

\begin{proof}
This is the discrete predictable-martingale uniqueness theorem.  A strongly
predictable process indexed by \(\mathbb N\) has \(M_{n+1}\)
\(\mathcal F_n\)-measurable for every \(n\), and the martingale identity gives
\[
  \mathbb E[M_{n+1}\mid\mathcal F_n]=M_n .
\]
Since \(M_{n+1}\) is already \(\mathcal F_n\)-measurable, the conditional
expectation is \(M_{n+1}\).  Hence \(M_{n+1}=M_n\) almost surely, and induction
from \(M_0=0\) gives the result.  In Lean this is exactly the Mathlib theorem
\(\mathrm{MeasureTheory.Martingale.eq\_zero\_of\_predictable'}\), followed by
the normalization \(M_0=0\).
\end{proof}


\begin{lemma}\label{lem:Martingale.eq_zero_of_predictable_finiteVariation_past_measurable_value}
  \uses{def:Martingale}
  \lean{MeasureTheory.Martingale.eq_zero_of_predictable_finiteVariation_past_measurable_value}
Let \(M\) be a true martingale and fix a deterministic time \(t\).  Suppose
that \(M_t\) is strongly measurable with respect to the past sigma-algebra
\(\bigvee_{s<t}\mathcal F_s\).  If
\[
  \mathbb E[M_t\mid \bigvee_{s<t}\mathcal F_s]=0
\]
almost surely, then \(M_t=0\) almost surely.
\end{lemma}

\begin{proof}
The past sigma-algebra is a sub-sigma-algebra of the ambient measurable
space, because every \(\mathcal F_s\) is.  Since \(M_t\) is integrable and
strongly measurable with respect to this sub-sigma-algebra, the conditional
expectation of \(M_t\) onto the past sigma-algebra is \(M_t\) itself.  The
assumed conditional-expectation equality is therefore exactly the desired
almost-sure equality \(M_t=0\).
\end{proof}


\begin{lemma}\label{lem:Martingale.eq_zero_of_predictable_finiteVariation_past_measurable_zero_increment}
  \uses{def:Martingale}
  \lean{MeasureTheory.Martingale.eq_zero_of_predictable_finiteVariation_past_measurable_zero_increment}
Let \(M\) be a true martingale.  If \(s\le t\) and the increment
\(M_t-M_s\) is strongly \(\mathcal F_s\)-measurable, then
\(M_t-M_s=0\) almost surely.
\end{lemma}

\begin{proof}
The martingale identity gives
\[
  \mathbb E[M_t\mid\mathcal F_s]=M_s
\]
almost surely.  Also \(M_s\) is already \(\mathcal F_s\)-measurable, so its
conditional expectation onto \(\mathcal F_s\) is \(M_s\).  By linearity of
conditional expectation,
\[
  \mathbb E[M_t-M_s\mid\mathcal F_s]=0
\]
almost surely.  Since the increment \(M_t-M_s\) is itself strongly
\(\mathcal F_s\)-measurable and integrable, this conditional expectation is
the increment.  Hence \(M_t-M_s=0\) almost surely.
\end{proof}


\begin{lemma}\label{lem:condExp_ae_eq_zero_of_ae_eq_zero}
  \lean{MeasureTheory.condExp_ae_eq_zero_of_ae_eq_zero}
If a real random variable \(X\) is zero almost surely, then its conditional
expectation onto any sub-\(\sigma\)-algebra is zero almost surely.
\end{lemma}

\begin{proof}
Replace \(X\) by the zero random variable inside conditional expectation using
almost-sure congruence of conditional expectation.  The conditional
expectation of the zero random variable is zero.
\end{proof}


\begin{lemma}\label{lem:LocallyBoundedVariationOn.exists_tendsto_left_univ}
  \uses{def:LocallyBoundedVariationOn}
  \lean{LocallyBoundedVariationOn.exists_tendsto_left_univ}
Let the time index have a bottom element and the order topology.  If a path
\(f\) has locally bounded variation on all of time, then for every
deterministic time \(t\) the left limit of \(f\) at \(t\) exists; equivalently,
there is an \(l\) such that \(f(s)\to l\) as \(s\uparrow t\).
\end{lemma}

\begin{proof}
Apply the local bounded-variation hypothesis to the compact interval
\([0,t]\).  A bounded-variation function on an ordered interval has a left
limit at the right endpoint.  The left-neighborhood of \(t\) inside
\((-\infty,t]\) is the same as the left-neighborhood of \(t\), so this gives
the desired limit.
\end{proof}

\begin{lemma}\label{lem:Martingale.stronglyMeasurable_leftLim_past_of_left_approach}
  \uses{def:Martingale, def:LocallyBoundedVariationOn, lem:LocallyBoundedVariationOn.exists_tendsto_left_univ}
  \lean{MeasureTheory.Martingale.stronglyMeasurable_leftLim_past_of_left_approach}
Let \(M\) be a real-valued true martingale and fix a deterministic time
\(t\).  Suppose each sample path has locally bounded variation.  Let
\((u_n)\) be deterministic times with \(u_n<t\) and \(u_n\to t\) from the
left.  Then the left-limit random variable
\[
  \omega\longmapsto M_{t-}(\omega)
\]
is strongly measurable with respect to the past sigma-algebra
\(\bigvee_{s<t}\mathcal F_s\).
\end{lemma}

\begin{proof}
For each \(n\), the martingale is strongly adapted, so \(M_{u_n}\) is
strongly \(\mathcal F_{u_n}\)-measurable.  Since \(u_n<t\),
\(\mathcal F_{u_n}\) is one of the sigma-algebras included in
\(\bigvee_{s<t}\mathcal F_s\), hence every \(M_{u_n}\) is strongly
measurable with respect to the past sigma-algebra.

For a fixed sample point \(\omega\), local bounded variation gives a left
limit at \(t\), and the theorem
Lemma~\ref{lem:LocallyBoundedVariationOn.exists_tendsto_left_univ} identifies
the Lean left-limit \(M_{t-}(\omega)\).  Composing this left-limit convergence
with \(u_n\to t\) from the left gives
\[
  M_{u_n}(\omega)\to M_{t-}(\omega).
\]
Thus the random variables \(M_{u_n}\) converge pointwise to
\(\omega\mapsto M_{t-}(\omega)\).  A pointwise sequential limit of strongly
measurable real random variables is strongly measurable, so the left-limit
random variable is past-measurable.
\end{proof}

\begin{lemma}\label{lem:Martingale.stronglyMeasurable_jump_leftLim_past_of_left_approach}
  \uses{def:Martingale, def:LocallyBoundedVariationOn, lem:IsStronglyPredictable.stronglyMeasurable_past, lem:Martingale.stronglyMeasurable_leftLim_past_of_left_approach}
  \lean{MeasureTheory.Martingale.stronglyMeasurable_jump_leftLim_past_of_left_approach}
Let \(M\) be a strongly predictable real-valued true martingale with locally
bounded-variation paths.  Fix \(t>\bot\), and let \(u_n<t\) tend to \(t\)
from the left.  Then the jump random variable
\[
  \omega\longmapsto M_t(\omega)-M_{t-}(\omega)
\]
is strongly measurable with respect to the past sigma-algebra
\(\bigvee_{s<t}\mathcal F_s\).
\end{lemma}

\begin{proof}
Strong predictability gives strong measurability of \(M_t\) with respect to
the past sigma-algebra by
Lemma~\ref{lem:IsStronglyPredictable.stronglyMeasurable_past}.  The preceding
lemma gives strong measurability of the left-limit random variable
\(M_{t-}\) with respect to the same sigma-algebra.  Strong measurability is
stable under subtraction for real-valued functions, so their difference is
past-measurable.
\end{proof}

\begin{lemma}\label{lem:Filtration.ae_eq_zero_of_past_setIntegral_eq_zero}
  \lean{MeasureTheory.Filtration.ae_eq_zero_of_past_setIntegral_eq_zero}
Let \((\mathcal F_s)\) be a linearly ordered filtration with bottom time, and
fix \(t>\bot\).  Let \(X\) be an integrable real random variable which is
strongly measurable with respect to the past sigma-algebra
\[
  \mathcal F_{t-}:=\bigvee_{s<t}\mathcal F_s .
\]
Assume that for every \(s<t\) and every \(A\in\mathcal F_s\),
\[
  \int_A X\,dP=0 .
\]
Then \(X=0\) almost surely.
\end{lemma}

\begin{proof}
Let
\[
  \mathcal C=\{A:\exists s<t,\ A\in\mathcal F_s\}.
\]
Because the time index is linearly ordered, \(\mathcal C\) is a
\(\pi\)-system: if \(A\in\mathcal F_s\) and \(B\in\mathcal F_r\), then
\(A\cap B\) is measurable in \(\mathcal F_{\max(s,r)}\), and
\(\max(s,r)<t\).  Moreover
\(\mathcal F_{t-}\) is exactly the sigma-algebra generated by
\(\mathcal C\).

First extend the zero-integral hypothesis from \(\mathcal C\) to every
\(\mathcal F_{t-}\)-measurable set by Dynkin's \(\pi\)-\(\lambda\)
induction.  The empty set is immediate.  For complements, use
\(\int_{A^c}X\,dP=\int X\,dP-\int_A X\,dP\); the total integral is zero
because \(\Omega\in\mathcal F_\bot\) and \(\bot<t\).  For pairwise disjoint
countable unions, countable additivity of the Bochner integral gives the sum
of the zero integrals.

Now \(X\) is integrable and \(\mathcal F_{t-}\)-strongly measurable, and its
integral over every \(\mathcal F_{t-}\)-measurable set is zero.  The standard
uniqueness theorem for Bochner integrals on a sub-sigma-algebra therefore
gives \(X=0\) almost surely.
\end{proof}

\begin{lemma}\label{lem:Martingale.ae_eq_leftLim_of_left_approach_past_setIntegral_zero}
  \uses{def:Martingale, def:LocallyBoundedVariationOn, lem:Martingale.stronglyMeasurable_jump_leftLim_past_of_left_approach, lem:Filtration.ae_eq_zero_of_past_setIntegral_eq_zero}
  \lean{MeasureTheory.Martingale.ae_eq_leftLim_of_left_approach_past_setIntegral_zero}
Let \(M\) be a strongly predictable true martingale with locally
bounded-variation paths.  Fix \(t>\bot\), and let \(u_n<t\) tend to \(t\)
from the left.  Suppose the jump
\[
  J_t:=M_t-M_{t-}
\]
is integrable, and suppose that for every \(s<t\) and every
\(A\in\mathcal F_s\),
\[
  \int_A J_t\,dP=0 .
\]
Then the deterministic jump at \(t\) vanishes:
\[
  M_t=M_{t-}
\]
almost surely.
\end{lemma}

\begin{proof}
By
Lemma~\ref{lem:Martingale.stronglyMeasurable_jump_leftLim_past_of_left_approach},
the jump \(J_t\) is strongly measurable with respect to
\(\bigvee_{s<t}\mathcal F_s\).  Apply
Lemma~\ref{lem:Filtration.ae_eq_zero_of_past_setIntegral_eq_zero} to
\(X=J_t\), using the assumed integrability and zero integrals on every
earlier-filtration set.  This gives \(J_t=0\) almost surely, which is exactly
\(M_t=M_{t-}\) after rewriting the definition of \(J_t\).
\end{proof}

\begin{lemma}\label{lem:Martingale.setIntegral_increment_eq_zero_of_measurableSet}
  \uses{def:Martingale}
  \lean{MeasureTheory.Martingale.setIntegral_increment_eq_zero_of_measurableSet}
Let \(M\) be a real-valued martingale.  If \(s\le u\le v\) and
\(A\in\mathcal F_s\), then
\[
  \int_A (M_v-M_u)\,dP=0 .
\]
\end{lemma}

\begin{proof}
The martingale identity gives
\[
  \mathbf E[M_v-M_u\mid\mathcal F_s]=0
\]
because \(s\le u\le v\).  Equivalently, multiply the future increment
\(M_v-M_u\) by the bounded \(\mathcal F_s\)-measurable random variable
\(\mathbf 1_A\).  The product is integrable since \(M_v-M_u\) is integrable
and \(\mathbf 1_A\) is bounded.  The conditional-expectation pull-out
identity then gives
\[
  \int \mathbf 1_A(M_v-M_u)\,dP=0 .
\]
This is the same as the set integral over \(A\).
\end{proof}

\begin{lemma}\label{lem:Martingale.setIntegral_jump_leftLim_eq_zero_of_left_approach_of_dominated}
  \uses{def:Martingale, def:LocallyBoundedVariationOn, lem:LocallyBoundedVariationOn.exists_tendsto_left_univ, lem:Martingale.setIntegral_increment_eq_zero_of_measurableSet}
  \lean{MeasureTheory.Martingale.setIntegral_jump_leftLim_eq_zero_of_left_approach_of_dominated}
Let \(M\) be a real-valued martingale whose paths have locally bounded
variation.  Fix \(t\), a deterministic sequence \(u_n<t\) tending to \(t\)
from the left, a strict-past time \(s<t\), and \(A\in\mathcal F_s\).
Assume \(u_n\ge s\) eventually.  Assume also that, on the restricted measure
\(P|_A\), the increments \(M_t-M_{u_n}\) are dominated by an integrable real
bound.  Then
\[
  \int_A (M_t-M_{t-})\,dP=0 .
\]
\end{lemma}

\begin{proof}
Work with the restricted measure \(P|_A\).  For every sample point, local
bounded variation gives a left limit at \(t\), and composing the left-limit
convergence with \(u_n\to t\) from the left gives
\[
  M_t(\omega)-M_{u_n}(\omega)\longrightarrow M_t(\omega)-M_{t-}(\omega).
\]
The assumed integrable domination lets us apply dominated convergence on
\(P|_A\), hence
\[
  \int_A (M_t-M_{u_n})\,dP
    \longrightarrow
  \int_A (M_t-M_{t-})\,dP .
\]
On the other hand, for all sufficiently large \(n\) we have \(s\le u_n<t\).
Lemma~\ref{lem:Martingale.setIntegral_increment_eq_zero_of_measurableSet}
therefore gives
\[
  \int_A (M_t-M_{u_n})\,dP=0
\]
eventually.  The same sequence of real numbers has limit \(0\) and the
displayed jump integral, so uniqueness of limits gives the desired zero
set integral.
\end{proof}

\begin{lemma}\label{lem:Tendsto.eventually_const_le_of_nhdsWithin_Iio}
  \lean{Filter.Tendsto.eventually_const_le_of_nhdsWithin_Iio}
Let \(u_n\) be a deterministic sequence tending to \(t\) from the left.
If \(s<t\), then \(s\le u_n\) for all sufficiently large \(n\).
\end{lemma}

\begin{proof}
The interval \((s,\infty)\) is a neighbourhood of \(t\) in the order
topology.  Since the within-neighbourhood filter at \(t\) along
\((-\infty,t)\) is finer than the ordinary neighbourhood filter, the
left-convergence of \(u_n\) implies \(u_n\in(s,\infty)\) eventually.  This
is stronger than the required eventual inequality \(s\le u_n\).
\end{proof}

\begin{lemma}\label{lem:norm_sub_le_two_mul_of_Icc_bound}
  \lean{MeasureTheory.norm_sub_le_two_mul_of_Icc_bound}
Fix a deterministic horizon \([\bot,t]\).  Suppose a path of \(M\) is bounded
there by a constant \(C\):
\[
  \|M_r(\omega)\|\le C \qquad (r\in[\bot,t]).
\]
If \(u\in[\bot,t]\), then
\[
  \|M_t(\omega)-M_u(\omega)\|\le 2C .
\]
\end{lemma}

\begin{proof}
Both \(t\) and \(u\) belong to \([\bot,t]\), so the horizon bound gives
\(\|M_t(\omega)\|\le C\) and \(\|M_u(\omega)\|\le C\).  The triangle
inequality gives
\[
  \|M_t(\omega)-M_u(\omega)\|
    \le \|M_t(\omega)\|+\|M_u(\omega)\|
    \le C+C=2C .
\]
\end{proof}

\begin{lemma}\label{lem:Martingale.ae_eq_leftLim_of_left_approach_past_setIntegral_zero_of_bound}
  \uses{def:Martingale, def:LocallyBoundedVariationOn, lem:Tendsto.eventually_const_le_of_nhdsWithin_Iio, lem:norm_sub_le_two_mul_of_Icc_bound, lem:Martingale.setIntegral_jump_leftLim_eq_zero_of_left_approach_of_dominated, lem:Martingale.ae_eq_leftLim_of_left_approach_past_setIntegral_zero}
  \lean{MeasureTheory.Martingale.ae_eq_leftLim_of_left_approach_past_setIntegral_zero_of_bound}
Let \(M\) be a strongly predictable true martingale whose paths have locally
bounded variation.  Fix \(t>\bot\), and let \(u_n<t\) tend to \(t\) from the
left.  Assume the jump \(M_t-M_{t-}\) is integrable.  Assume also that, almost
surely, \(M\) is bounded by a deterministic constant \(C\) on the horizon
\([\bot,t]\).  Then the deterministic jump at \(t\) vanishes:
\[
  M_t=M_{t-}
\]
almost surely.
\end{lemma}

\begin{proof}
Apply
Lemma~\ref{lem:Martingale.ae_eq_leftLim_of_left_approach_past_setIntegral_zero}.
It remains to prove that for every \(s<t\) and every
\(A\in\mathcal F_s\),
\[
  \int_A (M_t-M_{t-})\,dP=0 .
\]

Fix such \(s\) and \(A\).  By
Lemma~\ref{lem:Tendsto.eventually_const_le_of_nhdsWithin_Iio}, the
left-approaching sequence is eventually after \(s\).  On the restricted
measure \(P|_A\), use the constant function \(2C\) as a dominator.  It is
integrable because \(P\) is finite, hence so is the restricted measure.  The
pathwise bound follows from
Lemma~\ref{lem:norm_sub_le_two_mul_of_Icc_bound}: for each \(n\), the time
\(u_n\) lies in \([\bot,t]\) because \(\bot\le u_n\) and \(u_n<t\), so the
assumed horizon bound gives
\[
  \|M_t-M_{u_n}\|\le 2C
\]
almost surely, and therefore also almost surely under \(P|_A\).

Now
Lemma~\ref{lem:Martingale.setIntegral_jump_leftLim_eq_zero_of_left_approach_of_dominated}
applies with this eventual ordering and this constant dominator.  It gives
the required zero set integral for the fixed \(s\) and \(A\).  Since \(s\)
and \(A\) were arbitrary, the generated-past jump-removal wrapper gives
\(M_t=M_{t-}\) almost surely.
\end{proof}

\begin{lemma}\label{lem:Martingale.integrable_jump_leftLim_of_left_approach_of_bound}
  \uses{def:Martingale, def:LocallyBoundedVariationOn, lem:Martingale.stronglyMeasurable_jump_leftLim_past_of_left_approach, lem:norm_sub_le_two_mul_of_Icc_bound}
  \lean{MeasureTheory.Martingale.integrable_jump_leftLim_of_left_approach_of_bound}
Let \(M\) be a strongly predictable true martingale whose paths have locally
bounded variation.  Fix \(t>\bot\), and let \(u_n<t\) tend to \(t\) from the
left.  Suppose that, almost surely, \(M\) is bounded by a deterministic
constant \(C\ge0\) on the horizon \([\bot,t]\).  Then the deterministic jump
\[
  \omega\longmapsto M_t(\omega)-M_{t-}(\omega)
\]
is integrable.
\end{lemma}

\begin{proof}
Let \(J_t=M_t-M_{t-}\).  By
Lemma~\ref{lem:Martingale.stronglyMeasurable_jump_leftLim_past_of_left_approach},
the random variable \(J_t\) is strongly measurable with respect to the
generated strict-past sigma-algebra, hence it is a.e. strongly measurable for
the ambient measure.

It remains to dominate its norm by an integrable constant.  Work on the
almost-sure event on which the horizon bound holds.  For each \(n\), the time
\(u_n\) belongs to \([\bot,t]\), so
Lemma~\ref{lem:norm_sub_le_two_mul_of_Icc_bound} gives
\[
  \|M_t-M_{u_n}\|\le 2C .
\]
Local bounded variation gives the left limit at \(t\), and composing that
left-limit convergence with \(u_n\to t\) from the left gives
\[
  M_t-M_{u_n}\longrightarrow M_t-M_{t-}=J_t .
\]
Taking norms preserves convergence.  Since the closed interval
\((-\infty,2C]\) is closed in \(\mathbb R\), the limit also satisfies
\(\|J_t\|\le 2C\).

The constant function \(2C\) is integrable because the measure is finite.
The a.e. strong measurability of \(J_t\) and the a.e. bound
\(\|J_t\|\le 2C\) therefore imply that \(J_t\) is integrable.
\end{proof}

\begin{lemma}\label{lem:Martingale.ae_eq_leftLim_of_left_approach_past_setIntegral_zero_of_bound_no_integrability}
  \uses{def:Martingale, def:LocallyBoundedVariationOn, lem:Martingale.integrable_jump_leftLim_of_left_approach_of_bound, lem:Martingale.ae_eq_leftLim_of_left_approach_past_setIntegral_zero_of_bound}
  \lean{MeasureTheory.Martingale.ae_eq_leftLim_of_left_approach_past_setIntegral_zero_of_bound_no_integrability}
Let \(M\) be a strongly predictable true martingale whose paths have locally
bounded variation.  Fix \(t>\bot\), and let \(u_n<t\) tend to \(t\) from the
left.  Suppose that, almost surely, \(M\) is bounded by a deterministic
constant \(C\ge0\) on the horizon \([\bot,t]\).  Then the deterministic jump
at \(t\) vanishes:
\[
  M_t=M_{t-}
\]
almost surely.
\end{lemma}

\begin{proof}
First apply
Lemma~\ref{lem:Martingale.integrable_jump_leftLim_of_left_approach_of_bound}
to obtain integrability of \(M_t-M_{t-}\) from the same horizon bound and the
same left-approaching sequence.  Then all hypotheses of
Lemma~\ref{lem:Martingale.ae_eq_leftLim_of_left_approach_past_setIntegral_zero_of_bound}
are available, so that lemma gives \(M_t=M_{t-}\) almost surely.
\end{proof}

\begin{lemma}\label{lem:Martingale.ae_eq_leftLim_stoppedProcess_indicator_of_left_approach_of_bound}
  \uses{def:Martingale, def:LocallyBoundedVariationOn, lem:Martingale.ae_eq_leftLim_of_left_approach_past_setIntegral_zero_of_bound_no_integrability, lem:stoppedProcess_indicator_bound_on_Icc}
  \lean{MeasureTheory.Martingale.ae_eq_leftLim_stoppedProcess_indicator_of_left_approach_of_bound}
Let \(M\) be a strongly predictable true martingale with cadlag paths and
locally bounded variation.  Let \(\tau\) be a stopping time, and form the
usual stopped/indicator localization
\[
  Z_s(\omega)
  =
  \mathbf 1_{\{\bot<\tau(\omega)\}}
  M_{s\wedge\tau(\omega)}(\omega).
\]
Fix \(t>\bot\), and let \(u_n<t\) tend to \(t\) from the left.  Suppose that,
almost surely, the original process \(M\) is bounded by a deterministic
constant \(C\ge0\) on \([\bot,t]\).  Then the stopped process has no
deterministic jump at \(t\):
\[
  Z_t = Z_{t-}
\]
almost surely.
\end{lemma}

\begin{proof}
The stopped/indicator process \(Z\) is again a true martingale: apply the
stopped-process martingale theorem to \(M\), using right-continuity from the
cadlag-path hypothesis and the stopping-time property of \(\tau\).  It is
strongly predictable by the stopped-process predictability lemma.  Its paths
have locally bounded variation because stopping and multiplying by the
localizing indicator preserve locally bounded variation pathwise.

The deterministic horizon bound transfers from \(M\) to \(Z\) on the same
almost-sure event.  Indeed, for \(s\in[\bot,t]\), either
\(\bot<\tau(\omega)\), in which case \(Z_s(\omega)\) is an original value of
\(M\) at a time \(s\wedge\tau(\omega)\le t\), or the indicator is zero.  This
is precisely Lemma~\ref{lem:stoppedProcess_indicator_bound_on_Icc}.

All hypotheses of
Lemma~\ref{lem:Martingale.ae_eq_leftLim_of_left_approach_past_setIntegral_zero_of_bound_no_integrability}
are therefore available for \(Z\), with the same left-approaching sequence
and constant \(C\).  Applying that lemma gives \(Z_t=Z_{t-}\) almost surely.
\end{proof}

\begin{lemma}\label{lem:Martingale.ae_eq_leftLim_localizingSequence_stoppedProcess_indicator_of_left_approach_of_bound}
  \uses{def:Martingale, def:LocallyBoundedVariationOn, lem:Martingale.ae_eq_leftLim_stoppedProcess_indicator_of_left_approach_of_bound}
  \lean{MeasureTheory.Martingale.ae_eq_leftLim_localizingSequence_stoppedProcess_indicator_of_left_approach_of_bound}
Let \(M\) be as in the preceding lemma, and let \((\tau_n)\) be a localizing
sequence.  Fix \(t>\bot\), a deterministic sequence \(u_k<t\) tending to \(t\)
from the left, and deterministic constants \(C_n\ge0\).  Assume that, for
each \(n\), \(M\) is bounded by \(C_n\) on \([\bot,t]\) almost surely.  Then
for every \(n\), the stopped/indicator localization at \(\tau_n\) has no
deterministic jump at \(t\):
\[
  Z^n_t = Z^n_{t-}
\]
almost surely, where
\[
  Z^n_s(\omega)
  =
  \mathbf 1_{\{\bot<\tau_n(\omega)\}}
  M_{s\wedge\tau_n(\omega)}(\omega).
\]
\end{lemma}

\begin{proof}
Fix \(n\).  The localizing-sequence hypothesis says that \(\tau_n\) is a
stopping time.  Apply
Lemma~\ref{lem:Martingale.ae_eq_leftLim_stoppedProcess_indicator_of_left_approach_of_bound}
with \(\tau=\tau_n\), the bound \(C_n\), and the \(n\)-th almost-sure horizon
bound.  Since \(n\) was arbitrary, this proves the asserted family of
almost-sure jump-removal statements.
\end{proof}

\begin{lemma}\label{lem:stoppedProcess_indicator_eventuallyEq_left_of_lt}
  \uses{def:stoppedProcess}
  \lean{MeasureTheory.stoppedProcess_indicator_eventuallyEq_left_of_lt}
Let \(M\) be a real-valued process and let \(\tau\) be a stopping time value.
Fix a sample point \(\omega\) and a deterministic time \(t\).  If
\((t:\overline{\kappa})<\tau(\omega)\), then the stopped/indicator process
\[
  Z_s(\omega)
  =
  \mathbf 1_{\{\bot<\tau(\omega)\}}M_{s\wedge\tau(\omega)}(\omega)
\]
agrees with the original path \(s\mapsto M_s(\omega)\) eventually in the
left-neighborhood filter of \(t\):
\[
  Z_s(\omega)=M_s(\omega)
  \qquad\text{for all \(s<t\) sufficiently close to \(t\).}
\]
\end{lemma}

\begin{proof}
Let
\[
  S=\{s:\kappa\mid (s:\overline{\kappa})<\tau(\omega)\}.
\]
This set is an open neighborhood of \(t\), because it is the preimage of an
open lower interval in \(\overline{\kappa}\) under the continuous coercion
\(\kappa\to\overline{\kappa}\), and the hypothesis puts \(t\) in it.  Hence
\(S\) belongs to the left-neighborhood filter \(\mathcal N[<]t\).  For any
\(s\in S\), the stopped time is still after \(s\), so
\[
  s\wedge\tau(\omega)=s .
\]
Also \(\bot<\tau(\omega)\), since \(\bot\le t\) and \(t<\tau(\omega)\).
Thus the indicator is one and the stopped/indicator value is exactly
\(M_s(\omega)\).  This proves the eventual equality in
\(\mathcal N[<]t\).
\end{proof}

\begin{lemma}\label{lem:leftLim_stoppedProcess_indicator_eq_of_lt}
  \uses{def:stoppedProcess, def:LocallyBoundedVariationOn, lem:stoppedProcess_indicator_eventuallyEq_left_of_lt}
  \lean{MeasureTheory.leftLim_stoppedProcess_indicator_eq_of_lt}
In the setting of the preceding lemma, suppose additionally that the original
sample path has locally bounded variation.  If
\((t:\overline{\kappa})<\tau(\omega)\), then the left limit of the
stopped/indicator path at \(t\) is the original left limit:
\[
  Z_{t-}(\omega)=M_{t-}(\omega).
\]
\end{lemma}

\begin{proof}
Local bounded variation of the original path gives a left limit at \(t\), and
the project's left-limit choice theorem identifies
\(\operatorname{leftLim}(s\mapsto M_s(\omega),t)\) as the limit of the
original path along \(\mathcal N[<]t\).  By
Lemma~\ref{lem:stoppedProcess_indicator_eventuallyEq_left_of_lt}, the stopped
path and the original path are eventually equal along that same filter.
Therefore the stopped path also tends to the original left limit.

The same eventual convergence gives existence of a left limit for the stopped
path.  Applying the left-limit choice theorem to the stopped path gives
convergence to
\(\operatorname{leftLim}(s\mapsto Z_s(\omega),t)\).  Limits in \(\mathbb R\)
are unique, so the stopped left limit equals the original left limit.
\end{proof}

\begin{lemma}\label{lem:Martingale.ae_eq_leftLim_on_event_of_left_approach_of_bound}
  \uses{def:Martingale, def:LocallyBoundedVariationOn, lem:Martingale.ae_eq_leftLim_stoppedProcess_indicator_of_left_approach_of_bound, lem:leftLim_stoppedProcess_indicator_eq_of_lt}
  \lean{MeasureTheory.Martingale.ae_eq_leftLim_on_event_of_left_approach_of_bound}
Let \(M\) be a strongly predictable true martingale with cadlag paths and
locally bounded variation, and let \(\tau\) be a stopping time.  Fix
\(t>\bot\) and a deterministic sequence \(u_n<t\) tending to \(t\) from the
left.  Suppose that, almost surely, \(M\) is bounded by a deterministic
constant \(C\ge0\) on \([\bot,t]\).  Then on the event where the stopping time
lies strictly after \(t\), the original process has no deterministic jump:
\[
  (t:\overline{\kappa})<\tau(\omega)
  \quad\Longrightarrow\quad
  M_t(\omega)=M_{t-}(\omega)
\]
almost surely.
\end{lemma}

\begin{proof}
Apply
Lemma~\ref{lem:Martingale.ae_eq_leftLim_stoppedProcess_indicator_of_left_approach_of_bound}
to the stopped/indicator process \(Z\).  Work on the full-measure event where
that stopped jump-removal identity holds.  Now fix a sample point with
\((t:\overline{\kappa})<\tau(\omega)\).  At the terminal time \(t\), stopping
does nothing and the indicator is one, so \(Z_t(\omega)=M_t(\omega)\).  By
Lemma~\ref{lem:leftLim_stoppedProcess_indicator_eq_of_lt}, the stopped left
limit \(Z_{t-}(\omega)\) is the original left limit \(M_{t-}(\omega)\).  The
stopped jump-removal identity \(Z_t=Z_{t-}\) therefore rewrites to
\(M_t=M_{t-}\) on the event \(\{t<\tau\}\).
\end{proof}

\begin{lemma}\label{lem:ProbabilityTheory.IsLocalizingSequence.eventually_exists_gt}
  \lean{ProbabilityTheory.IsLocalizingSequence.eventually_exists_gt}
Let \((\tau_n)\) be a localizing sequence.  For every deterministic time
\(t\), almost surely there exists \(n\) such that
\[
  (t:\overline{\kappa})<\tau_n(\omega).
\]
\end{lemma}

\begin{proof}
The definition of a localizing sequence includes
\(\tau_n(\omega)\to\top\) along \(n\to\infty\), almost surely.  Work on this
full-measure convergence event.  The set
\[
  \{a:\overline{\kappa}\mid (t:\overline{\kappa})<a\}
\]
is an open neighborhood of \(\top\) in \(\overline{\kappa}\).  Therefore, for
all sufficiently large \(n\), \(\tau_n(\omega)\) belongs to this set.  Choosing
one such \(n\) gives \((t:\overline{\kappa})<\tau_n(\omega)\).
\end{proof}

\begin{lemma}\label{lem:Martingale.ae_eq_leftLim_localizingSequence_of_left_approach_of_bound}
  \uses{def:Martingale, def:LocallyBoundedVariationOn, lem:Martingale.ae_eq_leftLim_on_event_of_left_approach_of_bound, lem:ProbabilityTheory.IsLocalizingSequence.eventually_exists_gt}
  \lean{MeasureTheory.Martingale.ae_eq_leftLim_localizingSequence_of_left_approach_of_bound}
Let \(M\) be as in the preceding lemma, and let \((\tau_n)\) be a localizing
sequence.  Fix \(t>\bot\), a deterministic sequence \(u_k<t\) tending to \(t\)
from the left, and deterministic constants \(C_n\ge0\).  Assume that, for
each \(n\), \(M\) is bounded by \(C_n\) on \([\bot,t]\) almost surely.  Then
the original process has no deterministic jump at \(t\):
\[
  M_t=M_{t-}
\]
almost surely.
\end{lemma}

\begin{proof}
For each \(n\), apply
Lemma~\ref{lem:Martingale.ae_eq_leftLim_on_event_of_left_approach_of_bound}
with the stopping time \(\tau_n\), the constant \(C_n\), and the corresponding
almost-sure horizon bound.  This gives, for each \(n\), an almost-sure
implication
\[
  (t:\overline{\kappa})<\tau_n(\omega)
  \Longrightarrow
  M_t(\omega)=M_{t-}(\omega).
\]
Intersect these countably many full-measure events.  The localizing-sequence
cover lemma says that almost surely there exists some \(n\) with
\((t:\overline{\kappa})<\tau_n(\omega)\).  On that sample point, use the
corresponding implication to obtain \(M_t(\omega)=M_{t-}(\omega)\).
\end{proof}

\begin{lemma}\label{lem:Martingale.ae_eq_leftLim_of_left_approach_past_measurable_increments}
  \uses{def:Martingale, def:LocallyBoundedVariationOn, lem:Martingale.eq_zero_of_predictable_finiteVariation_past_measurable_zero_increment, lem:LocallyBoundedVariationOn.exists_tendsto_left_univ}
  \lean{MeasureTheory.Martingale.ae_eq_leftLim_of_left_approach_past_measurable_increments}
Let \(M\) be a real-valued true martingale.  Fix a deterministic time \(t\),
and suppose each sample path has locally bounded variation.  Let
\((u_n)\) be deterministic times tending to \(t\) from the left.  Assume
\(u_n<t\) for every \(n\), and assume that each increment
\(M_t-M_{u_n}\) is strongly \(\mathcal F_{u_n}\)-measurable.  Then the jump
at \(t\) vanishes along this announcing sequence:
\[
  M_t=M_{t-}
\]
almost surely, where \(M_{t-}\) is represented in Lean by the left limit
\(\operatorname{leftLim}(s\mapsto M_s(\omega),t)\).
\end{lemma}

\begin{proof}
For each \(n\), apply
Lemma~\ref{lem:Martingale.eq_zero_of_predictable_finiteVariation_past_measurable_zero_increment}
to the increment \(M_t-M_{u_n}\).  The hypotheses \(u_n<t\) and strong
\(\mathcal F_{u_n}\)-measurability give
\[
  M_t-M_{u_n}=0
\]
almost surely.  Since there are only countably many \(n\), intersect these
almost-sure events and work on the event where \(M_{u_n}=M_t\) for every
\(n\).

On that event, the sequence \(M_{u_n}(\omega)\) is the constant sequence with
value \(M_t(\omega)\).  On the other hand,
Lemma~\ref{lem:LocallyBoundedVariationOn.exists_tendsto_left_univ} gives the
left limit of the path \(s\mapsto M_s(\omega)\), and composing this left-limit
convergence with \(u_n\to t\) from the left shows
\[
  M_{u_n}(\omega)\to M_{t-}(\omega).
\]
Uniqueness of limits in \(\mathbb R\) therefore gives
\(M_t(\omega)=M_{t-}(\omega)\).  This proves the asserted almost-sure
equality.
\end{proof}


\begin{lemma}\label{lem:BoundedVariationOn.sum_norm_sub_le_toReal_eVariationOn}
  \uses{def:LocallyBoundedVariationOn}
  \lean{BoundedVariationOn.sum_norm_sub_le_toReal_eVariationOn}
Let \(f\) be a real-valued function of bounded variation on a set \(S\).
For every finite nondecreasing chain \(u_0,\ldots,u_n\) in \(S\),
\[
  \sum_{k<n}|f(u_{k+1})-f(u_k)|
  \le \operatorname{Var}_S(f).
\]
\end{lemma}

\begin{proof}
The extended variation \(\operatorname{eVariationOn}(f,S)\) is the supremum
over all such finite ordered chains of the sum of the corresponding
extended distances.  The chosen chain is one admissible term in that
supremum, so its extended-distance sum is bounded by the extended variation.
Since \(f\) has bounded variation, the extended variation is finite; passing
to real values identifies each extended distance with
\(|f(u_{k+1})-f(u_k)|\) and gives the stated inequality.
\end{proof}


\begin{lemma}\label{lem:BoundedVariationOn.sq_increment_sum_le_uniform_bound}
  \uses{lem:BoundedVariationOn.sum_norm_sub_le_toReal_eVariationOn}
  \lean{BoundedVariationOn.sq_increment_sum_le_uniform_bound}
Let \(f\) have bounded variation on \(S\).  If along a finite ordered chain
all increments satisfy
\[
  |f(u_{k+1})-f(u_k)|\le \delta
\]
with \(\delta\ge0\), then
\[
  \sum_{k<n}|f(u_{k+1})-f(u_k)|^2
  \le
  \delta\,\operatorname{Var}_S(f).
\]
\end{lemma}

\begin{proof}
Each summand satisfies
\[
  |f(u_{k+1})-f(u_k)|^2
  \le
  \delta\,|f(u_{k+1})-f(u_k)|
\]
by the uniform increment bound and nonnegativity of the absolute value.
Summing and applying
Lemma~\ref{lem:BoundedVariationOn.sum_norm_sub_le_toReal_eVariationOn}
gives the result.
\end{proof}


\begin{lemma}\label{lem:BoundedVariationOn.sq_increment_sum_tendsto_zero_of_uniform_bound}
  \uses{lem:BoundedVariationOn.sq_increment_sum_le_uniform_bound}
  \lean{BoundedVariationOn.sq_increment_sum_tendsto_zero_of_uniform_bound}
Let \(f\) have bounded variation on \(S\).  Along a sequence of finite
ordered chains, suppose the largest increment is bounded by deterministic
numbers \(\delta_n\ge0\) with \(\delta_n\to0\).  Then the corresponding
square-increment sums tend to \(0\).
\end{lemma}

\begin{proof}
The square-increment sums are nonnegative.  By
Lemma~\ref{lem:BoundedVariationOn.sq_increment_sum_le_uniform_bound}, the
\(n\)-th sum is bounded above by
\(\delta_n\operatorname{Var}_S(f)\).  Since the variation is a fixed finite
real number and \(\delta_n\to0\), the upper bound tends to \(0\).  The squeeze
theorem gives the conclusion.
\end{proof}


\begin{lemma}\label{lem:Martingale.integral_mul_increment_eq_zero_of_stronglyMeasurable}
  \uses{def:Martingale}
  \lean{MeasureTheory.Martingale.integral_mul_increment_eq_zero_of_stronglyMeasurable}
Let \(M\) be a true martingale.  If \(s\le t\), \(Z\) is a real random
variable strongly measurable with respect to \(\mathcal F_s\), and
\(Z(M_t-M_s)\) is integrable, then
\[
  \mathbb E[Z(M_t-M_s)] = 0.
\]
\end{lemma}

\begin{proof}
The martingale property gives
\(\mathbb E[M_t-M_s\mid\mathcal F_s]=0\).  Since \(Z\) is
\(\mathcal F_s\)-measurable, it can be pulled out of conditional expectation:
\[
  \mathbb E[Z(M_t-M_s)\mid\mathcal F_s]
  =
  Z\,\mathbb E[M_t-M_s\mid\mathcal F_s]
  =
  0.
\]
Integrating both sides gives the result.
\end{proof}

\begin{lemma}\label{lem:Martingale.integral_increment_mul_increment_eq_zero}
  \uses{def:Martingale, lem:Martingale.integral_mul_increment_eq_zero_of_stronglyMeasurable}
  \lean{MeasureTheory.Martingale.integral_increment_mul_increment_eq_zero}
Let \(M\) be a true martingale.  For deterministic times
\(a\le b\le c\le d\), assume
\((M_b-M_a)(M_d-M_c)\) is integrable.  Then
\[
  \mathbb E[(M_b-M_a)(M_d-M_c)] = 0.
\]
\end{lemma}

\begin{proof}
The earlier increment \(M_b-M_a\) is \(\mathcal F_c\)-measurable because
\(a,b\le c\) and \(M\) is adapted.  Apply
Lemma~\ref{lem:Martingale.integral_mul_increment_eq_zero_of_stronglyMeasurable}
to the later increment \(M_d-M_c\).
\end{proof}

\begin{lemma}\label{lem:Martingale.integral_mul_increment_eq_zero}
  \uses{lem:Martingale.integral_increment_mul_increment_eq_zero}
  \lean{MeasureTheory.Martingale.integral_mul_increment_eq_zero}
Under the same hypotheses as
Lemma~\ref{lem:Martingale.integral_increment_mul_increment_eq_zero}, disjoint
deterministic martingale increments are orthogonal in expectation.
\end{lemma}

\begin{proof}
This is the convenient named wrapper used by later arguments.  It is exactly
Lemma~\ref{lem:Martingale.integral_increment_mul_increment_eq_zero} with the
product-integrability hypothesis kept explicit.
\end{proof}

\begin{lemma}\label{lem:Finset.sum_range_sub_consecutive}
  \lean{Finset.sum_range_sub_consecutive}
Let \(f:\mathbb N\to G\) be a sequence in an additive commutative group.
For every \(n\),
\[
  \sum_{i=0}^{n-1} (f_{i+1}-f_i)=f_n-f_0 .
\]
\end{lemma}

\begin{proof}
Induct on \(n\).  The case \(n=0\) is empty.  In the successor step, split
off the last summand and use the induction hypothesis; the remaining group
identity is \((f_n-f_0)+(f_{n+1}-f_n)=f_{n+1}-f_0\).
\end{proof}

\begin{lemma}\label{lem:Martingale.partition_increment_sum_eq_terminal}
  \uses{lem:Finset.sum_range_sub_consecutive}
  \lean{MeasureTheory.Martingale.partition_increment_sum_eq_terminal}
Let \(M\) be a martingale with \(M_\bot=0\).  If a deterministic finite
partition \(u_0,\ldots,u_n\) satisfies \(u_0=\bot\) and \(u_n=t\), then
\[
  \sum_{i=0}^{n-1} (M_{u_{i+1}}-M_{u_i}) = M_t
\]
as random variables.
\end{lemma}

\begin{proof}
Apply Lemma~\ref{lem:Finset.sum_range_sub_consecutive} pointwise to
\(f_i=M_{u_i}(\omega)\).  The endpoint assumptions rewrite
\(f_n-f_0\) as \(M_t(\omega)-M_\bot(\omega)=M_t(\omega)\).
\end{proof}

\begin{lemma}\label{lem:Martingale.integral_partition_cross_increment_eq_zero}
  \uses{lem:Martingale.integral_increment_mul_increment_eq_zero}
  \lean{MeasureTheory.Martingale.integral_partition_cross_increment_eq_zero}
Let \(M\) be a true martingale and let \(u:\mathbb N\to\kappa\) be a
monotone deterministic partition.  If \(i<j\) and the product of adjacent
increments is integrable, then
\[
  \int (M_{u_{i+1}}-M_{u_i})(M_{u_{j+1}}-M_{u_j})\,dP=0 .
\]
\end{lemma}

\begin{proof}
Monotonicity gives
\(u_i\le u_{i+1}\le u_j\le u_{j+1}\).  Apply
Lemma~\ref{lem:Martingale.integral_increment_mul_increment_eq_zero} to the
two disjoint deterministic increments.
\end{proof}

\begin{lemma}\label{lem:Martingale.integrable_partition_increment_mul_increment_of_bound}
  \uses{def:Martingale}
  \lean{MeasureTheory.Martingale.integrable_partition_increment_mul_increment_of_bound}
Assume \(P\) is finite and \(M\) is a real-valued true martingale.  If
\(0\le C\) and, almost surely, \(|M_{u_k}|\le C\) for every partition index
\(k\), then every product
\[
  (M_{u_{i+1}}-M_{u_i})(M_{u_{j+1}}-M_{u_j})
\]
of adjacent partition increments is integrable.
\end{lemma}

\begin{proof}
The martingale's adapted strongly measurable values make each increment
strongly measurable, hence so is the product.  On the almost-sure bounded
event, every adjacent increment has norm at most \(2C\), so the product has
norm at most \(4C^2\).  A strongly measurable random variable bounded by a
constant on a finite measure space is integrable.
\end{proof}

\begin{lemma}\label{lem:Martingale.integrable_partition_value_mul_increment_of_bound}
  \uses{def:Martingale}
  \lean{MeasureTheory.Martingale.integrable_partition_value_mul_increment_of_bound}
Assume \(P\) is finite and \(M\) is a real-valued true martingale.  If
\(0\le C\) and, almost surely, \(|M_{u_k}|\le C\) for every partition index
\(k\), then each product of a partition value with the next adjacent increment,
\[
  M_{u_i}(M_{u_{i+1}}-M_{u_i}),
\]
is integrable.
\end{lemma}

\begin{proof}
The martingale's adapted strongly measurable values make both
\(M_{u_i}\) and \(M_{u_{i+1}}-M_{u_i}\) strongly measurable, hence their
product is a.e. strongly measurable.  On the almost-sure bounded event,
\(|M_{u_i}|\le C\) and
\(|M_{u_{i+1}}-M_{u_i}|\le |M_{u_{i+1}}|+|M_{u_i}|\le 2C\).  The product is
therefore bounded by \(2C^2\), an integrable constant on a finite measure
space.
\end{proof}

\begin{lemma}\label{lem:Martingale.integrable_partition_sq_increment_of_bound}
  \uses{lem:Martingale.integrable_partition_increment_mul_increment_of_bound}
  \lean{MeasureTheory.Martingale.integrable_partition_sq_increment_of_bound}
Under the same bounded-partition hypotheses, each squared adjacent increment
\((M_{u_{i+1}}-M_{u_i})^2\) is integrable.
\end{lemma}

\begin{proof}
This is the previous product-integrability lemma with the two increment
indices both equal to \(i\).
\end{proof}

\begin{lemma}\label{lem:Martingale.integrable_partition_sq_increment_sum_of_bound}
  \uses{lem:Martingale.integrable_partition_sq_increment_of_bound}
  \lean{MeasureTheory.Martingale.integrable_partition_sq_increment_sum_of_bound}
Under the same bounded-partition hypotheses, every finite sum of squared
adjacent increments is integrable.
\end{lemma}

\begin{proof}
Finite sums of integrable random variables are integrable, and each summand is
integrable by Lemma~\ref{lem:Martingale.integrable_partition_sq_increment_of_bound}.
\end{proof}

\begin{lemma}\label{lem:Martingale.integral_partition_sq_increment_sum_of_bound}
  \uses{lem:Martingale.integrable_partition_sq_increment_of_bound}
  \lean{MeasureTheory.Martingale.integral_partition_sq_increment_sum_of_bound}
Under the same bounded-partition hypotheses,
\[
  \int \sum_{i=0}^{n-1}(M_{u_{i+1}}-M_{u_i})^2\,dP
  =
  \sum_{i=0}^{n-1}\int (M_{u_{i+1}}-M_{u_i})^2\,dP .
\]
\end{lemma}

\begin{proof}
Use linearity of the Bochner integral over a finite sum.  The integrability
side conditions are supplied by
Lemma~\ref{lem:Martingale.integrable_partition_sq_increment_of_bound}.
\end{proof}

\begin{lemma}\label{lem:Martingale.integral_partition_cross_increment_eq_zero_of_bound}
  \uses{lem:Martingale.integral_partition_cross_increment_eq_zero, lem:Martingale.integrable_partition_increment_mul_increment_of_bound}
  \lean{MeasureTheory.Martingale.integral_partition_cross_increment_eq_zero_of_bound}
Under the same bounded-partition hypotheses, all ordered mixed terms in a
monotone deterministic partition have zero integral:
\[
  \int (M_{u_{i+1}}-M_{u_i})(M_{u_{j+1}}-M_{u_j})\,dP=0
  \qquad (i<j).
\]
\end{lemma}

\begin{proof}
The boundedness lemma gives the product-integrability hypothesis required by
Lemma~\ref{lem:Martingale.integral_partition_cross_increment_eq_zero}, and
that orthogonality lemma gives the zero integral.
\end{proof}

\begin{lemma}\label{lem:Martingale.integrable_sq_terminal_of_ae_bound}
  \uses{def:Martingale}
  \lean{MeasureTheory.Martingale.integrable_sq_terminal_of_ae_bound}
Assume \(P\) is finite and \(M\) is a real-valued true martingale.  If
\(|M_t|\le C\) almost surely, then \(M_t^2\) is integrable.
\end{lemma}

\begin{proof}
The terminal value \(M_t\) is strongly measurable by the martingale
hypotheses, so \(M_t^2\) is a.e. strongly measurable.  Replace \(C\) by
\(\max(C,0)\); the square is then pointwise bounded by this constant square on
the almost-sure event.  Boundedness by an integrable constant on a finite
measure space gives integrability.
\end{proof}

\begin{lemma}\label{lem:Martingale.integral_sq_terminal_eq_partition_sq_sum_of_bound}
  \uses{lem:Martingale.integrable_partition_value_mul_increment_of_bound, lem:Martingale.integral_mul_increment_eq_zero_of_stronglyMeasurable, lem:Martingale.integrable_sq_terminal_of_ae_bound}
  \lean{MeasureTheory.Martingale.integral_sq_terminal_eq_partition_sq_sum_of_bound}
Let \(M\) be a real-valued true martingale with \(M_\bot=0\), and let
\(u_0,\ldots,u_n\) be a monotone deterministic finite partition from
\(\bot\) to \(t\).  If the partition values are uniformly bounded almost
surely, then
\[
  \int M_t^2\,dP
  =
  \sum_{i=0}^{n-1}\int (M_{u_{i+1}}-M_{u_i})^2\,dP .
\]
\end{lemma}

\begin{proof}
Induct on the number \(n\) of adjacent increments.  The case \(n=0\) follows
from \(u_0=t=\bot\) and \(M_\bot=0\).  In the successor step write
\[
  M_{u_{n+1}} = M_{u_n} + (M_{u_{n+1}}-M_{u_n}).
\]
After squaring and integrating, the first two terms are integrable by the
terminal-square and squared-increment integrability lemmas, while the cross
term
\[
  M_{u_n}(M_{u_{n+1}}-M_{u_n})
\]
is integrable by
Lemma~\ref{lem:Martingale.integrable_partition_value_mul_increment_of_bound}
and has zero integral by
Lemma~\ref{lem:Martingale.integral_mul_increment_eq_zero_of_stronglyMeasurable},
because \(M_{u_n}\) is measurable with respect to the earlier filtration and
\(u_n\le u_{n+1}\).  The induction hypothesis identifies the integral of
\(M_{u_n}^2\) with the sum of the previous squared-increment integrals, and
the successor partition sum is obtained by adding the last diagonal term.
\end{proof}

\begin{lemma}\label{lem:MeasureTheory.ae_eq_zero_of_integral_sq_eq_zero}
  \lean{MeasureTheory.ae_eq_zero_of_integral_sq_eq_zero}
If a real random variable \(X\) has integrable square and
\(\int X^2\,dP=0\), then \(X=0\) almost surely.
\end{lemma}

\begin{proof}
The function \(X^2\) is nonnegative.  The nonnegative-integral criterion
therefore gives \(X^2=0\) almost surely, and then \(X=0\) almost surely
pointwise on that event.
\end{proof}

\begin{lemma}\label{lem:Martingale.integral_sq_terminal_eq_zero_of_refining_partitions}
  \uses{lem:Martingale.integral_sq_terminal_eq_partition_sq_sum_of_bound, lem:BoundedVariationOn.sq_increment_sum_le_uniform_bound}
  \lean{MeasureTheory.Martingale.integral_sq_terminal_eq_zero_of_refining_partitions}
Let \(M\) be a real-valued true martingale with \(M_\bot=0\), and fix
\(t>\bot\).  Suppose that for each \(n\) we have a monotone deterministic
partition \(u^n_0,\ldots,u^n_{m_n}\) of \([\bot,t]\).  Assume \(0\le C\)
and \(0\le V\), assume that \(|M_s|\le C\) almost surely for every partition
point, and assume that, almost surely, the total variation of each path on
\([\bot,t]\) is bounded by the deterministic constant \(V\).  If the
pathwise square-increment sums
\[
  F_n(\omega)=\sum_{i=0}^{m_n-1}
    (M_{u^n_{i+1}}(\omega)-M_{u^n_i}(\omega))^2
\]
tend to \(0\) almost surely, then \(M_t^2\) is integrable and
\[
  \int M_t^2\,dP=0 .
\]
\end{lemma}

\begin{proof}
For every \(n\), apply
Lemma~\ref{lem:Martingale.integral_sq_terminal_eq_partition_sq_sum_of_bound}
to the partition \(u^n\).  This gives
\[
  \int M_t^2\,dP = \int F_n\,dP
\]
for all \(n\).  The finite sum \(F_n\) is a.e. strongly measurable because its
summands are squared martingale increments.  On the almost-sure event where
the values are bounded by \(C\) and the variation is bounded by \(V\), every
increment has norm at most \(2C\), and
Lemma~\ref{lem:BoundedVariationOn.sq_increment_sum_le_uniform_bound} gives a
uniform deterministic bound \(F_n\le 2CV\).  The constant \(2CV\) is integrable
on a finite measure space.  Dominated convergence applied to the a.e. limit
\(F_n\to0\) therefore gives \(\int F_n\,dP\to0\).  Since the left-hand side
\(\int M_t^2\,dP\) is independent of \(n\), it must be \(0\).  Integrability of
\(M_t^2\) follows from the same bounded-terminal-value hypothesis.
\end{proof}

\begin{lemma}\label{lem:Martingale.integral_sq_terminal_eq_zero_of_bounded_continuous_finiteVariation_of_variation_bound}
  \uses{def:Martingale, lem:Martingale.integral_sq_terminal_eq_zero_of_refining_partitions, lem:BoundedVariationOn.sq_increment_sum_tendsto_zero_of_uniform_bound}
  \lean{MeasureTheory.Martingale.integral_sq_terminal_eq_zero_of_bounded_continuous_finiteVariation_of_variation_bound}
Let \(M\) be a real-valued true martingale with \(M_\bot=0\), and fix
\(t>\bot\).  Suppose that \(M\) is bounded by a deterministic constant \(C\)
on the deterministic horizon \([\bot,t]\), and suppose that the total
variation of each path on \([\bot,t]\) is almost surely bounded by a
deterministic constant \(V\).  Suppose moreover that there are deterministic
monotone partitions
\[
  \bot=u^n_0\le u^n_1\le\cdots\le u^n_{m_n}=t
\]
and nonnegative numbers \(\delta_n\to0\) such that, almost surely, every
partition increment of every path satisfies
\[
  |M_{u^n_{i+1}}-M_{u^n_i}|\le \delta_n
  \qquad (0\le i<m_n).
\]
Then \(M_t^2\) is integrable and
\[
  \int M_t^2\,dP=0 .
\]
\end{lemma}

\begin{proof}
Apply
Lemma~\ref{lem:BoundedVariationOn.sq_increment_sum_tendsto_zero_of_uniform_bound}
pathwise on the almost-sure event where the variation is bounded by \(V\) and
where the partition increment bound holds.  This gives almost-sure convergence
of the square-increment sums along the chosen deterministic partitions.  The
bounded-value hypothesis gives the a.e. partition-value bound required by
Lemma~\ref{lem:Martingale.integral_sq_terminal_eq_zero_of_refining_partitions},
and the deterministic variation bound supplies its domination constant.  That
lemma then yields both integrability of \(M_t^2\) and the identity
\(\int M_t^2\,dP=0\).
\end{proof}

\begin{lemma}\label{lem:Martingale.integral_sq_terminal_eq_zero_of_bounded_continuous_finiteVariation_of_variation_bound_eventual_modulus}
  \uses{def:Martingale, lem:Martingale.integral_sq_terminal_eq_zero_of_refining_partitions, lem:BoundedVariationOn.sq_increment_sum_tendsto_zero_of_uniform_bound}
  \lean{MeasureTheory.Martingale.integral_sq_terminal_eq_zero_of_bounded_continuous_finiteVariation_of_variation_bound_eventual_modulus}
Let \(M\) be a real-valued true martingale with \(M_\bot=0\), and fix
\(t>\bot\).  Suppose \(M\) is bounded by a deterministic constant \(C\) on
\([\bot,t]\), and suppose the total variation of each path on \([\bot,t]\) is
almost surely bounded by a deterministic constant \(V\).  Let
\(\bot=u^n_0\le\cdots\le u^n_{m_n}=t\) be deterministic monotone partitions.
Assume that, for almost every sample point \(\omega\), there is a
nonnegative sequence \(\delta^\omega_n\to0\) such that
\[
  |M_{u^n_{i+1}}(\omega)-M_{u^n_i}(\omega)|
  \le \delta^\omega_n
  \qquad (0\le i<m_n)
\]
for every \(n\).  Then \(M_t^2\) is integrable and
\[
  \int M_t^2\,dP=0 .
\]
\end{lemma}

\begin{proof}
Work on the almost-sure event where the horizon bound, the deterministic
variation bound, and the pathwise increment modulus all hold.  For a fixed
sample point in this event, apply
Lemma~\ref{lem:BoundedVariationOn.sq_increment_sum_tendsto_zero_of_uniform_bound}
with the sample-point dependent sequence \(\delta^\omega_n\).  This proves
that the square-increment sums along the deterministic partitions converge to
zero almost surely.  The horizon bound still gives the partition-value bound
needed by
Lemma~\ref{lem:Martingale.integral_sq_terminal_eq_zero_of_refining_partitions},
and the deterministic variation bound supplies the dominated-convergence
constant.  Applying that refining-partition lemma gives the desired
integrability and zero integral.
\end{proof}

\begin{lemma}\label{lem:Martingale.integral_sq_terminal_eq_zero_of_bounded_continuous_finiteVariation_of_variation_bound_sup_modulus}
  \uses{def:Martingale, lem:Martingale.integral_sq_terminal_eq_zero_of_bounded_continuous_finiteVariation_of_variation_bound_eventual_modulus}
  \lean{MeasureTheory.Martingale.integral_sq_terminal_eq_zero_of_bounded_continuous_finiteVariation_of_variation_bound_sup_modulus}
Let \(M\) be a real-valued true martingale with \(M_\bot=0\), and fix
\(t>\bot\).  Suppose \(M\) is bounded by a deterministic constant \(C\) on
\([\bot,t]\), and suppose the total variation of each path on \([\bot,t]\) is
almost surely bounded by a deterministic constant \(V\).  Let
\(\bot=u^n_0\le\cdots\le u^n_{m_n}=t\) be deterministic monotone partitions.
Assume that, for almost every sample point \(\omega\), the finite maximum
\[
  \max_{0\le i<m_n}
  |M_{u^n_{i+1}}(\omega)-M_{u^n_i}(\omega)|
\]
tends to \(0\) as \(n\to\infty\).  Then \(M_t^2\) is integrable and
\[
  \int M_t^2\,dP=0 .
\]
\end{lemma}

\begin{proof}
For each sample point in the almost-sure event where the maximum-increment
convergence holds, use that finite maximum as the modulus
\(\delta^\omega_n\).  It is nonnegative, tends to \(0\) by hypothesis, and
every partition increment is bounded by it by the defining property of a
finite maximum.  Hence the hypotheses of
Lemma~\ref{lem:Martingale.integral_sq_terminal_eq_zero_of_bounded_continuous_finiteVariation_of_variation_bound_eventual_modulus}
hold, and that lemma gives the required integrability and zero integral.
\end{proof}

\begin{lemma}\label{lem:UniformContinuousOn.finite_partition_sup_nnnorm_sub_tendsto_zero}
  \lean{UniformContinuousOn.finite_partition_sup_nnnorm_sub_tendsto_zero}
Let \(S\) be a subset of a uniform time space, and let \(f:S\to\mathbb R\)
be represented by a function \(f\) that is uniformly continuous on \(S\).
Let \(u^n_0,\ldots,u^n_{m_n}\) be finite deterministic point families in
\(S\).  Assume their adjacent-pair mesh tends to zero in the entourage sense:
for every entourage \(U\) of the time space, eventually in \(n\),
\[
  (u^n_i,u^n_{i+1})\in U
  \qquad (0\le i<m_n).
\]
Then
\[
  \max_{0\le i<m_n}|f(u^n_{i+1})-f(u^n_i)| \longrightarrow 0 .
\]
\end{lemma}

\begin{proof}
Fix \(\varepsilon>0\).  Uniform continuity of \(f\) on \(S\) gives an
entourage \(U\) of the time space such that whenever \(x,y\in S\) and
\((x,y)\in U\), one has \(|f(x)-f(y)|<\varepsilon\).  By the mesh hypothesis,
eventually every adjacent pair of the \(n\)-th partition lies in \(U\).  Since
all partition points lie in \(S\), every adjacent increment of \(f\) is then
strictly smaller than \(\varepsilon\).  The finite maximum of those increments
is therefore at most \(\varepsilon\) eventually.  This is exactly convergence
of the finite maximum to \(0\).
\end{proof}

\begin{lemma}\label{lem:ContinuousOn.uniformContinuousOn_Icc}
  \lean{ContinuousOn.uniformContinuousOn_Icc}
Let \(f\) be a real-valued function on a compact ordered time interval
\([a,b]\).  If \(f\) is continuous on \([a,b]\), then \(f\) is uniformly
continuous on \([a,b]\).
\end{lemma}

\begin{proof}
The interval \([a,b]\) is compact.  Apply the Heine--Cantor theorem for
continuous maps on compact sets: the restriction of \(f\) to \([a,b]\) is
uniformly continuous.
\end{proof}

\begin{lemma}\label{lem:Martingale.integral_sq_terminal_eq_zero_of_bounded_continuous_finiteVariation_of_variation_bound_uniformContinuousOn}
  \uses{def:Martingale, lem:UniformContinuousOn.finite_partition_sup_nnnorm_sub_tendsto_zero, lem:Martingale.integral_sq_terminal_eq_zero_of_bounded_continuous_finiteVariation_of_variation_bound_sup_modulus}
  \lean{MeasureTheory.Martingale.integral_sq_terminal_eq_zero_of_bounded_continuous_finiteVariation_of_variation_bound_uniformContinuousOn}
Let \(M\) be a real-valued true martingale with \(M_\bot=0\), and fix
\(t>\bot\).  Suppose \(M\) is bounded by a deterministic constant \(C\) on
\([\bot,t]\), and suppose the total variation of each path on \([\bot,t]\) is
almost surely bounded by a deterministic constant \(V\).  Let
\(\bot=u^n_0\le\cdots\le u^n_{m_n}=t\) be deterministic monotone partitions
whose adjacent-pair mesh tends to zero in the entourage sense.  If every path
\(s\mapsto M_s(\omega)\) is uniformly continuous on \([\bot,t]\), then
\(M_t^2\) is integrable and
\[
  \int M_t^2\,dP=0 .
\]
\end{lemma}

\begin{proof}
Apply
Lemma~\ref{lem:UniformContinuousOn.finite_partition_sup_nnnorm_sub_tendsto_zero}
to each sample path \(s\mapsto M_s(\omega)\) on \(S=[\bot,t]\).  The
deterministic mesh hypothesis is independent of \(\omega\), while uniform
continuity may vary with the sample point.  Thus, for every \(\omega\), the
finite maximum of the partition increments tends to \(0\).  This gives the
almost-sure maximum-increment hypothesis required by
Lemma~\ref{lem:Martingale.integral_sq_terminal_eq_zero_of_bounded_continuous_finiteVariation_of_variation_bound_sup_modulus}.
The remaining boundedness, variation, partition, and martingale hypotheses
are unchanged, so that lemma yields the desired terminal square-integral
identity.
\end{proof}

\begin{lemma}\label{lem:Martingale.integral_sq_terminal_eq_zero_of_bounded_continuous_finiteVariation_of_variation_bound_continuousOn}
  \uses{def:Martingale, lem:ContinuousOn.uniformContinuousOn_Icc, lem:Martingale.integral_sq_terminal_eq_zero_of_bounded_continuous_finiteVariation_of_variation_bound_uniformContinuousOn}
  \lean{MeasureTheory.Martingale.integral_sq_terminal_eq_zero_of_bounded_continuous_finiteVariation_of_variation_bound_continuousOn}
Let \(M\) be a real-valued true martingale with \(M_\bot=0\), and fix
\(t>\bot\).  Suppose the time interval \([\bot,t]\) is compact in the time
uniformity.  Suppose \(M\) is bounded by a deterministic constant \(C\) on
\([\bot,t]\), and suppose the total variation of each path on \([\bot,t]\) is
almost surely bounded by a deterministic constant \(V\).  Let
\(\bot=u^n_0\le\cdots\le u^n_{m_n}=t\) be deterministic monotone partitions
whose adjacent-pair mesh tends to zero in the entourage sense.  If every path
\(s\mapsto M_s(\omega)\) is continuous on \([\bot,t]\), then \(M_t^2\) is
integrable and
\[
  \int M_t^2\,dP=0 .
\]
\end{lemma}

\begin{proof}
For each sample point \(\omega\), apply
Lemma~\ref{lem:ContinuousOn.uniformContinuousOn_Icc} to the continuous path
\(s\mapsto M_s(\omega)\) on \([\bot,t]\).  This gives the pathwise uniform
continuity hypothesis required by
Lemma~\ref{lem:Martingale.integral_sq_terminal_eq_zero_of_bounded_continuous_finiteVariation_of_variation_bound_uniformContinuousOn}.
All partition, mesh, boundedness, and deterministic variation-bound
hypotheses are unchanged, so that lemma gives the stated integrability and
zero square-integral identity.
\end{proof}

\begin{lemma}\label{lem:Martingale.eq_zero_of_bounded_continuous_finiteVariation_of_variation_bound_continuousOn}
  \uses{def:Martingale, lem:MeasureTheory.ae_eq_zero_of_integral_sq_eq_zero, lem:Martingale.integral_sq_terminal_eq_zero_of_bounded_continuous_finiteVariation_of_variation_bound_continuousOn}
  \lean{MeasureTheory.Martingale.eq_zero_of_bounded_continuous_finiteVariation_of_variation_bound_continuousOn}
Under the hypotheses of
Lemma~\ref{lem:Martingale.integral_sq_terminal_eq_zero_of_bounded_continuous_finiteVariation_of_variation_bound_continuousOn},
the terminal value vanishes almost surely:
\[
  M_t=0 \quad\text{almost surely.}
\]
\end{lemma}

\begin{proof}
First apply
Lemma~\ref{lem:Martingale.integral_sq_terminal_eq_zero_of_bounded_continuous_finiteVariation_of_variation_bound_continuousOn}
to obtain integrability of \(M_t^2\) and the identity
\(\int M_t^2\,dP=0\).  Then apply
Lemma~\ref{lem:MeasureTheory.ae_eq_zero_of_integral_sq_eq_zero} to the
integrable terminal random variable \(M_t\).
\end{proof}

\begin{lemma}\label{lem:Martingale.eq_zero_on_event_of_bounded_continuous_finiteVariation_of_variation_bound_continuousOn}
  \uses{def:Martingale, lem:Martingale.eq_zero_of_bounded_continuous_finiteVariation_of_variation_bound_continuousOn}
  \lean{MeasureTheory.Martingale.eq_zero_on_event_of_bounded_continuous_finiteVariation_of_variation_bound_continuousOn}
Let \(Z\) be a real-valued true martingale with \(Z_\bot=0\).  Fix
\(t>\bot\), deterministic monotone partitions of \([\bot,t]\) whose
adjacent-pair mesh tends to zero, a deterministic bound \(C\), and a
deterministic variation bound \(V\).  Assume \(Z\) satisfies the hypotheses of
Lemma~\ref{lem:Martingale.eq_zero_of_bounded_continuous_finiteVariation_of_variation_bound_continuousOn}:
it is pathwise continuous on \([\bot,t]\), bounded by \(C\) there almost
surely, and has total variation at most \(V\) there almost surely.

Let \(M\) be another process and let \(E\) be an event.  If, almost surely,
\(Z_t(\omega)=M_t(\omega)\) for every \(\omega\in E\), then
\[
  \omega\in E \Longrightarrow M_t(\omega)=0
\]
almost surely.
\end{lemma}

\begin{proof}
Apply
Lemma~\ref{lem:Martingale.eq_zero_of_bounded_continuous_finiteVariation_of_variation_bound_continuousOn}
to the localized martingale \(Z\).  This gives \(Z_t=0\) almost surely.  On
the almost-sure event where \(Z_t=0\) and where \(Z_t=M_t\) on \(E\), every
\(\omega\in E\) satisfies \(M_t(\omega)=Z_t(\omega)=0\).
\end{proof}

\begin{lemma}\label{lem:MeasureTheory.ae_eq_zero_of_eventually_event_zero_exhaustion}
  \lean{MeasureTheory.ae_eq_zero_of_eventually_event_zero_exhaustion}
Let \(f\) be a real-valued random variable and let \((E_n)_{n\ge 0}\) be
events.  Suppose that almost every sample point belongs to some \(E_n\), and
suppose that for every \(n\), \(f=0\) almost surely on \(E_n\).  Then
\(f=0\) almost surely.
\end{lemma}

\begin{proof}
Intersect the countably many almost-sure statements \(f=0\) on \(E_n\) over
all \(n\), and also intersect the almost-sure exhaustion event
\(\{\omega:\exists n,\omega\in E_n\}\).  For any sample point in this
intersection choose an index \(n\) with \(\omega\in E_n\), then apply the
\(n\)-th local zero statement.
\end{proof}

\begin{lemma}\label{lem:Martingale.eq_zero_of_localized_bounded_continuous_finiteVariation_of_variation_bound_continuousOn}
  \uses{def:Martingale, lem:Martingale.eq_zero_on_event_of_bounded_continuous_finiteVariation_of_variation_bound_continuousOn, lem:MeasureTheory.ae_eq_zero_of_eventually_event_zero_exhaustion}
  \lean{MeasureTheory.Martingale.eq_zero_of_localized_bounded_continuous_finiteVariation_of_variation_bound_continuousOn}
Let \(N\) be a process and let \((Z_n)_{n\ge 0}\) be a countable family of
real-valued true martingales with \(Z_n(\bot)=0\).  Fix a deterministic time
\(t\), deterministic monotone partitions of \([\bot,t]\) whose adjacent-pair
mesh tends to zero, and events \(E_n\) which cover almost every sample point.
For each \(n\), assume \(Z_n\) is bounded on \([\bot,t]\) by a deterministic
constant \(C_n\), has total variation on \([\bot,t]\) bounded by a deterministic
constant \(V_n\), and is pathwise continuous on \([\bot,t]\).  Assume also that,
on \(E_n\), \(Z_n(t)\) agrees with \(N(t)\) almost surely.  Then
\[
  N_t=0
\]
almost surely.
\end{lemma}

\begin{proof}
For each fixed \(n\), apply
Lemma~\ref{lem:Martingale.eq_zero_on_event_of_bounded_continuous_finiteVariation_of_variation_bound_continuousOn}
to the martingale \(Z_n\), the event \(E_n\), and the deterministic bounds
\(C_n,V_n\).  This gives
\[
  \omega\in E_n \Longrightarrow N_t(\omega)=0
\]
almost surely, for every \(n\).  Since the events \(E_n\) cover almost every
sample point, Lemma~\ref{lem:MeasureTheory.ae_eq_zero_of_eventually_event_zero_exhaustion}
applied to the terminal random variable \(N_t\) gives \(N_t=0\) almost surely.
\end{proof}

\begin{lemma}\label{lem:Martingale.eq_zero_of_stoppedProcess_bounded_continuous_finiteVariation_of_variation_bound_continuousOn}
  \uses{def:Martingale, lem:Martingale.eq_zero_of_localized_bounded_continuous_finiteVariation_of_variation_bound_continuousOn}
  \lean{MeasureTheory.Martingale.eq_zero_of_stoppedProcess_bounded_continuous_finiteVariation_of_variation_bound_continuousOn}
Let \(N\) be a process with \(N_\bot=0\), and let \(\tau_n\) be a sequence of
extended-time cutoffs such that for almost every sample point there is an
index \(n\) with \(t<\tau_n(\omega)\).  Put
\[
  Z^n_s(\omega)
  =
  \mathbf 1_{\{\bot<\tau_n(\omega)\}}
  N_{s\wedge \tau_n(\omega)}(\omega).
\]
Assume each \(Z^n\) is a true martingale and satisfies the bounded,
continuous, deterministic-variation-bound hypotheses of
Lemma~\ref{lem:Martingale.eq_zero_of_localized_bounded_continuous_finiteVariation_of_variation_bound_continuousOn}
on \([\bot,t]\), with the same deterministic partitions and mesh hypothesis.
Then \(N_t=0\) almost surely.
\end{lemma}

\begin{proof}
Apply
Lemma~\ref{lem:Martingale.eq_zero_of_localized_bounded_continuous_finiteVariation_of_variation_bound_continuousOn}
to the localized family \(Z^n\) and to the events
\[
  E_n=\{\omega : t<\tau_n(\omega)\}.
\]
These events cover almost every sample point by hypothesis.  On \(E_n\) we have
\(\bot<\tau_n(\omega)\) and the stopping has not yet occurred by time \(t\), so
the stopped process satisfies \(Z^n_t(\omega)=N_t(\omega)\).  The normalization
\(Z^n_\bot=0\) follows from \(N_\bot=0\) on \(\{\bot<\tau_n\}\) and from the
indicator being zero off that event.  The localized-family lemma then gives the
desired terminal zero statement for \(N_t\).
\end{proof}

\begin{lemma}\label{lem:Martingale.eq_zero_of_localizingSequence_bounded_continuous_finiteVariation_of_variation_bound_continuousOn}
  \uses{def:Martingale, lem:Martingale.eq_zero_of_stoppedProcess_bounded_continuous_finiteVariation_of_variation_bound_continuousOn}
  \lean{MeasureTheory.Martingale.eq_zero_of_localizingSequence_bounded_continuous_finiteVariation_of_variation_bound_continuousOn}
Let \(N\) be a real-valued true martingale with \(N_\bot=0\), and suppose
the paths of \(N\) are cadlag.  Let \((\tau_n)\) be a localizing sequence for
the filtration.  Fix \(t\), deterministic monotone partitions of
\([\bot,t]\) whose adjacent-pair mesh tends to zero, and deterministic bounds
\(C_n,V_n\).  Assume that, for each \(n\), the stopped/indicator process
\[
  Z^n_s(\omega)
  =
  \mathbf 1_{\{\bot<\tau_n(\omega)\}}
  N_{s\wedge\tau_n(\omega)}(\omega)
\]
is bounded by \(C_n\) on \([\bot,t]\) almost surely, has total variation at
most \(V_n\) on \([\bot,t]\) almost surely, and is continuous on
\([\bot,t]\) for every sample point.  Then \(N_t=0\) almost surely.
\end{lemma}

\begin{proof}
The stopped/indicator processes \(Z^n\) are true martingales because \(N\) is
a true martingale, the paths of \(N\) are right-continuous, and each
\(\tau_n\) is a stopping time.  The stopping-time fact is part of the
localizing-sequence hypothesis, and right-continuity follows from the cadlag
assumption.

It remains only to supply the covering event required by
Lemma~\ref{lem:Martingale.eq_zero_of_stoppedProcess_bounded_continuous_finiteVariation_of_variation_bound_continuousOn}.
Since \(\tau_n(\omega)\to\top\) almost surely, for almost every \(\omega\) the
open neighbourhood \((t,\top]\) of \(\top\) eventually contains
\(\tau_n(\omega)\).  In particular, for such an \(\omega\) there is an index
\(n\) with \(t<\tau_n(\omega)\).  Thus the events
\(\{\omega:t<\tau_n(\omega)\}\) cover almost surely.  The stopped-process
lemma now applies with the already assumed deterministic bounds, variation
bounds, continuity, partition, and mesh hypotheses.
\end{proof}

\begin{lemma}\label{lem:Martingale.eq_zero_of_localizingSequence_of_bound_variation_bound_continuousOn}
  \uses{def:Martingale, lem:ae_stoppedProcess_indicator_bound_variation_on_Icc, lem:stoppedProcess_indicator_continuousOn_Icc, lem:Martingale.eq_zero_of_localizingSequence_bounded_continuous_finiteVariation_of_variation_bound_continuousOn}
  \lean{MeasureTheory.Martingale.eq_zero_of_localizingSequence_of_bound_variation_bound_continuousOn}
Let \(N\) be a real-valued true martingale with \(N_\bot=0\), and suppose the
paths of \(N\) are cadlag.  Let \((\tau_n)\) be a localizing sequence.  Fix a
deterministic horizon \([\bot,t]\), deterministic monotone partitions of that
horizon with adjacent-pair mesh tending to zero, and deterministic bounds
\(C_n,V_n\).  Assume that, for each \(n\), the original process \(N\) is
bounded by \(C_n\) on \([\bot,t]\) almost surely and has total variation at
most \(V_n\) on \([\bot,t]\) almost surely.  Assume also that the stopped
pieces
\[
  Z^n_s(\omega)=
  \mathbf 1_{\{\bot<\tau_n(\omega)\}}\,
  N_{s\wedge\tau_n(\omega)}(\omega)
\]
are continuous on \([\bot,t]\) for every sample point.  Then \(N_t=0\) almost
surely.
\end{lemma}

\begin{proof}
For each \(n\), apply
Lemma~\ref{lem:ae_stoppedProcess_indicator_bound_variation_on_Icc} with
\(\tau=\tau_n\), \(C=C_n\), and \(V=V_n\).  This converts the almost-sure
bound and variation hypotheses for \(N\) into the corresponding almost-sure
bound and variation hypotheses for the stopped/indicator process \(Z^n\).
The present Lean statement still assumes the continuity of the stopped
pieces.  The next lemma packages the common situation where continuity of the
original path has already been established; in that case
Lemma~\ref{lem:stoppedProcess_indicator_continuousOn_Icc} supplies exactly
the stopped-piece continuity.  The deterministic partitions and mesh
hypothesis are still assumptions.  Therefore all hypotheses of
Lemma~\ref{lem:Martingale.eq_zero_of_localizingSequence_bounded_continuous_finiteVariation_of_variation_bound_continuousOn}
are available, and that lemma gives \(N_t=0\) almost surely.
\end{proof}

\begin{lemma}\label{lem:Martingale.eq_zero_of_localizingSequence_of_bound_variation_bound_original_continuousOn}
  \uses{def:Martingale, lem:stoppedProcess_indicator_continuousOn_Icc, lem:Martingale.eq_zero_of_localizingSequence_of_bound_variation_bound_continuousOn}
  \lean{MeasureTheory.Martingale.eq_zero_of_localizingSequence_of_bound_variation_bound_original_continuousOn}
Let \(N\) be a real-valued true martingale with \(N_\bot=0\), and suppose the
paths of \(N\) are cadlag.  Let \((\tau_n)\) be a localizing sequence.  Fix a
deterministic horizon \([\bot,t]\), deterministic monotone partitions of that
horizon with adjacent-pair mesh tending to zero, and deterministic bounds
\(C_n,V_n\).  Assume that, for each \(n\), the original process \(N\) is
bounded by \(C_n\) on \([\bot,t]\) almost surely and has total variation at
most \(V_n\) on \([\bot,t]\) almost surely.  Assume in addition that the
original paths \(s\mapsto N_s(\omega)\) are continuous on \([\bot,t]\) for
every sample point \(\omega\).  Then \(N_t=0\) almost surely.
\end{lemma}

\begin{proof}
For each \(n\) and \(\omega\), apply
Lemma~\ref{lem:stoppedProcess_indicator_continuousOn_Icc} to the original
path-continuity hypothesis for \(N\), with the stopping time specialized to
\(\tau_n\).  This proves that the stopped/indicator path
\[
  s\longmapsto
  \mathbf 1_{\{\bot<\tau_n(\omega)\}}\,
  N_{s\wedge\tau_n(\omega)}(\omega)
\]
is continuous on \([\bot,t]\).  The original a.e. bound and variation-bound
hypotheses, the deterministic partitions, and the mesh hypothesis are
unchanged.  Hence
Lemma~\ref{lem:Martingale.eq_zero_of_localizingSequence_of_bound_variation_bound_continuousOn}
applies and gives \(N_t=0\) almost surely.
\end{proof}

\begin{lemma}\label{lem:Martingale.eq_zero_and_leftLim_eq_zero_of_localizingSequence_of_bound_variation_bound_original_continuousOn}
  \uses{def:Martingale, def:LocallyBoundedVariationOn, lem:Martingale.ae_eq_leftLim_localizingSequence_of_left_approach_of_bound, lem:Martingale.eq_zero_of_localizingSequence_of_bound_variation_bound_original_continuousOn}
  \lean{MeasureTheory.Martingale.eq_zero_and_leftLim_eq_zero_of_localizingSequence_of_bound_variation_bound_original_continuousOn}
Let \(N\) be a strongly predictable real-valued true martingale with
\(N_\bot=0\), cadlag paths, and locally bounded variation.  Let
\((\tau_n)\) be a localizing sequence.  Fix \(t>\bot\), a deterministic
left-approaching sequence \(v_k<t\) with \(v_k\to t\) from the left, and
deterministic horizon bounds \(C_n\).  Assume also the explicit hypotheses
needed by the localized continuous finite-variation square-integral wrapper:
deterministic monotone partitions of \([\bot,t]\) with mesh tending to zero,
deterministic variation bounds \(V_n\), almost-sure \(C_n\)-bounds and
\(V_n\)-variation bounds on \([\bot,t]\), and pathwise continuity on the
deterministic horizon.  Then both
\[
  N_t=0
  \quad\text{and}\quad
  N_{t-}=0
\]
almost surely.
\end{lemma}

\begin{proof}
Apply
Lemma~\ref{lem:Martingale.eq_zero_of_localizingSequence_of_bound_variation_bound_original_continuousOn}
to the explicit localized continuous finite-variation inputs.  This gives
\(N_t=0\) almost surely.

Independently, apply
Lemma~\ref{lem:Martingale.ae_eq_leftLim_localizingSequence_of_left_approach_of_bound}
to the same martingale, predictability and locally bounded-variation
hypotheses, using the explicit left-approaching sequence and the same
almost-sure horizon bounds.  This gives \(N_t=N_{t-}\) almost surely.  On the
intersection of these two full-measure events, symmetry of the equality
\(N_t=N_{t-}\) and transitivity with \(N_t=0\) give \(N_{t-}=0\).  Thus the
terminal value and the chosen left limit are both zero almost surely.
\end{proof}

\begin{lemma}\label{lem:Martingale.eq_zero_and_leftLim_eq_zero_of_localizingSequence_of_bound_variation_bound_original_continuousOn_of_neBot_left}
  \uses{def:Martingale, def:LocallyBoundedVariationOn, lem:Filter.exists_seq_lt_tendsto_nhdsWithin_Iio_of_neBot, lem:Martingale.eq_zero_and_leftLim_eq_zero_of_localizingSequence_of_bound_variation_bound_original_continuousOn}
  \lean{MeasureTheory.Martingale.eq_zero_and_leftLim_eq_zero_of_localizingSequence_of_bound_variation_bound_original_continuousOn_of_neBot_left}
Let \(N\) be a strongly predictable real-valued true martingale with
\(N_\bot=0\), cadlag paths, and locally bounded variation.  Let
\((\tau_n)\) be a localizing sequence and fix \(t>\bot\).  Assume the
left-neighborhood filter at \(t\) is nontrivial.  Assume the same original-path
inputs as in
Lemma~\ref{lem:Martingale.eq_zero_and_leftLim_eq_zero_of_localizingSequence_of_bound_variation_bound_original_continuousOn}:
deterministic monotone partitions of \([\bot,t]\) with mesh tending to zero,
deterministic constants \(C_n,V_n\ge0\), almost-sure \(C_n\)-bounds and
\(V_n\)-variation bounds for \(N\) on \([\bot,t]\), and pathwise continuity of
\(N\) on \([\bot,t]\).  Then \(N_t=0\) and \(N_{t-}=0\) almost surely.
\end{lemma}

\begin{proof}
Use
Lemma~\ref{lem:Filter.exists_seq_lt_tendsto_nhdsWithin_Iio_of_neBot}
to choose a deterministic sequence \(v_k<t\) tending to \(t\) from the left.
With this sequence supplied, every remaining hypothesis is exactly one of the
hypotheses of
Lemma~\ref{lem:Martingale.eq_zero_and_leftLim_eq_zero_of_localizingSequence_of_bound_variation_bound_original_continuousOn}.
Apply that lemma.  This wrapper derives only the left-approaching sequence; it
does not derive the deterministic bounds, variation bounds, original-path
continuity, partitions, or mesh.
\end{proof}

\begin{lemma}\label{lem:Martingale.ae_eq_leftLim_stoppedProcess_indicator_of_left_approach_of_stopped_bound}
  \uses{def:Martingale, def:LocallyBoundedVariationOn, lem:Martingale.ae_eq_leftLim_of_left_approach_past_setIntegral_zero_of_bound_no_integrability}
  \lean{MeasureTheory.Martingale.ae_eq_leftLim_stoppedProcess_indicator_of_left_approach_of_stopped_bound}
Let \(N\) be a strongly predictable real-valued true martingale with cadlag
paths and locally bounded variation.  Let \(\tau\) be a stopping time and put
\[
  Z_s(\omega)=
  \mathbf 1_{\{\bot<\tau(\omega)\}}\,
  N_{s\wedge\tau(\omega)}(\omega).
\]
Fix \(t>\bot\), a deterministic left-approaching sequence \(v_k<t\) with
\(v_k\to t\) from the left, and a deterministic constant \(C\ge 0\).  Assume
directly that the stopped process \(Z\) is bounded by \(C\) on
\([\bot,t]\) almost surely.  Then \(Z_t=Z_{t-}\) almost surely.
\end{lemma}

\begin{proof}
The proof is the stopped-process jump-removal argument without first deriving
the stopped bound from an original-process bound.  The stopped/indicator
process \(Z\) is a true martingale by stopped martingale stability, is
strongly predictable by stopped predictability, and has locally bounded
variation by the stopped finite-variation stability lemma.  The assumed
almost-sure bound is already the bounded-horizon hypothesis for \(Z\).  Hence
Lemma~\ref{lem:Martingale.ae_eq_leftLim_of_left_approach_past_setIntegral_zero_of_bound_no_integrability}
applies to \(Z\) and gives \(Z_t=Z_{t-}\) almost surely.
\end{proof}

\begin{lemma}\label{lem:Martingale.ae_eq_leftLim_on_event_of_left_approach_of_stopped_bound}
  \uses{def:Martingale, def:LocallyBoundedVariationOn, lem:Martingale.ae_eq_leftLim_stoppedProcess_indicator_of_left_approach_of_stopped_bound, lem:leftLim_stoppedProcess_indicator_eq_of_lt}
  \lean{MeasureTheory.Martingale.ae_eq_leftLim_on_event_of_left_approach_of_stopped_bound}
In the setting of the preceding lemma, suppose \(E=\{\omega:t<\tau(\omega)\}\).
If the stopped process \(Z\) is bounded by \(C\) on \([\bot,t]\) almost
surely, then on \(E\) the original process has no jump at \(t\):
\[
  N_t(\omega)=N_{t-}(\omega)
\]
for almost every \(\omega\in E\).
\end{lemma}

\begin{proof}
Apply
Lemma~\ref{lem:Martingale.ae_eq_leftLim_stoppedProcess_indicator_of_left_approach_of_stopped_bound}
to obtain \(Z_t=Z_{t-}\) almost surely.  On the event \(t<\tau(\omega)\),
the terminal values agree, \(Z_t(\omega)=N_t(\omega)\), by the definition of
stopping.  The same event puts the stopped and original paths equal on a
left-neighborhood of \(t\), so
Lemma~\ref{lem:leftLim_stoppedProcess_indicator_eq_of_lt} gives
\(Z_{t-}(\omega)=N_{t-}(\omega)\).  Combining these equalities proves the
event-localized original jump-removal statement.
\end{proof}

\begin{lemma}\label{lem:Martingale.ae_eq_leftLim_localizingSequence_of_left_approach_of_stopped_bound}
  \uses{def:Martingale, def:LocallyBoundedVariationOn, lem:Martingale.ae_eq_leftLim_on_event_of_left_approach_of_stopped_bound, lem:ProbabilityTheory.IsLocalizingSequence.eventually_exists_gt}
  \lean{MeasureTheory.Martingale.ae_eq_leftLim_localizingSequence_of_left_approach_of_stopped_bound}
Let \(N\) be a strongly predictable real-valued true martingale with cadlag
paths and locally bounded variation.  Let \((\tau_n)\) be a localizing
sequence.  Fix \(t>\bot\), a deterministic left-approaching sequence
\(v_k<t\), and constants \(C_n\ge0\).  Assume that each stopped/indicator
process
\[
  Z^n_s(\omega)=
  \mathbf 1_{\{\bot<\tau_n(\omega)\}}\,
  N_{s\wedge\tau_n(\omega)}(\omega)
\]
is bounded by \(C_n\) on \([\bot,t]\) almost surely.  Then
\[
  N_t=N_{t-}
\]
almost surely.
\end{lemma}

\begin{proof}
For each \(n\), apply
Lemma~\ref{lem:Martingale.ae_eq_leftLim_on_event_of_left_approach_of_stopped_bound}
to \(\tau_n\) and \(C_n\).  This gives \(N_t=N_{t-}\) almost surely on the
event \(\{t<\tau_n\}\).  Intersect the resulting countably many full-measure
sets.  The localizing-sequence cover
Lemma~\ref{lem:ProbabilityTheory.IsLocalizingSequence.eventually_exists_gt}
says that for almost every \(\omega\) some \(n\) satisfies
\(t<\tau_n(\omega)\).  The corresponding event-local equality then proves the
global almost-sure jump-removal statement.
\end{proof}

\begin{lemma}\label{lem:Martingale.eq_zero_and_leftLim_eq_zero_of_localizingSequence_bounded_continuous_finiteVariation_of_variation_bound_continuousOn}
  \uses{def:Martingale, def:LocallyBoundedVariationOn, lem:Martingale.eq_zero_of_localizingSequence_bounded_continuous_finiteVariation_of_variation_bound_continuousOn, lem:Martingale.ae_eq_leftLim_localizingSequence_of_left_approach_of_stopped_bound}
  \lean{MeasureTheory.Martingale.eq_zero_and_leftLim_eq_zero_of_localizingSequence_bounded_continuous_finiteVariation_of_variation_bound_continuousOn}
Let \(N\) be a strongly predictable real-valued true martingale with
\(N_\bot=0\), cadlag paths, and locally bounded variation.  Let
\((\tau_n)\) be a localizing sequence.  Fix \(t>\bot\), a deterministic
left-approaching sequence \(v_k<t\), deterministic monotone partitions of
\([\bot,t]\) whose mesh tends to zero, and constants \(C_n,V_n\ge0\).
Assume directly that every stopped/indicator process \(Z^n\) is bounded by
\(C_n\) on \([\bot,t]\) almost surely, has total variation at most \(V_n\)
on \([\bot,t]\) almost surely, and is pathwise continuous on \([\bot,t]\).
Then both \(N_t=0\) and \(N_{t-}=0\) almost surely.
\end{lemma}

\begin{proof}
Apply
Lemma~\ref{lem:Martingale.eq_zero_of_localizingSequence_bounded_continuous_finiteVariation_of_variation_bound_continuousOn}
to the stopped/indicator martingales, using the stopped-piece bound,
variation, continuity, partition, and mesh hypotheses.  This gives
\(N_t=0\) almost surely.  Separately apply
Lemma~\ref{lem:Martingale.ae_eq_leftLim_localizingSequence_of_left_approach_of_stopped_bound}
with the same stopped-piece horizon bounds to get \(N_t=N_{t-}\) almost
surely.  As before, symmetry and transitivity of almost-sure equality give
\(N_{t-}=0\).  This is the stopped-bound analogue of
Lemma~\ref{lem:Martingale.eq_zero_and_leftLim_eq_zero_of_localizingSequence_of_bound_variation_bound_original_continuousOn}.
\end{proof}

\begin{lemma}\label{lem:Filter.exists_seq_lt_tendsto_nhdsWithin_Iio_of_neBot}
  \lean{Filter.exists_seq_lt_tendsto_nhdsWithin_Iio_of_neBot}
Let \(\kappa\) be a second-countable topological preorder and fix a point
\(t\).  If the left-neighborhood filter
\[
  \mathcal N[t^-] := \mathcal N(t)\cap\{s:s<t\}
\]
is nontrivial, then there is a deterministic sequence \(v_n<t\) for every
\(n\), with \(v_n\to t\) along the left-neighborhood filter.
\end{lemma}

\begin{proof}
The filter \(\mathcal N[t^-]\) is countably generated: second countability
gives first countability of neighborhoods, and Mathlib's
countable-generation theorem for within-neighborhood filters transfers this
to \(\mathcal N(t)\cap\{s:s<t\}\).  The set \(\{s:s<t\}\) itself belongs to
the within filter, so under the nontriviality hypothesis the predicate
\(s<t\) is frequent in that filter.  Apply the countably-generated filter
sequence extraction theorem with the frequent predicate \(s<t\).  It returns
a sequence tending to the within filter and satisfying \(v_n<t\) for every
index \(n\).
\end{proof}

\begin{lemma}\label{lem:Martingale.eq_zero_and_leftLim_eq_zero_of_localizingSequence_bounded_continuous_finiteVariation_of_variation_bound_continuousOn_of_neBot_left}
  \uses{def:Martingale, def:LocallyBoundedVariationOn, lem:Filter.exists_seq_lt_tendsto_nhdsWithin_Iio_of_neBot, lem:Martingale.eq_zero_and_leftLim_eq_zero_of_localizingSequence_bounded_continuous_finiteVariation_of_variation_bound_continuousOn}
  \lean{MeasureTheory.Martingale.eq_zero_and_leftLim_eq_zero_of_localizingSequence_bounded_continuous_finiteVariation_of_variation_bound_continuousOn_of_neBot_left}
Let \(N\) be a strongly predictable real-valued true martingale with
\(N_\bot=0\), cadlag paths, and locally bounded variation.  Let
\((\tau_n)\) be a localizing sequence.  Fix \(t>\bot\).  Assume the
left-neighborhood filter at \(t\) is nontrivial.  Assume also the stopped-piece
inputs of
Lemma~\ref{lem:Martingale.eq_zero_and_leftLim_eq_zero_of_localizingSequence_bounded_continuous_finiteVariation_of_variation_bound_continuousOn}:
deterministic monotone partitions of \([\bot,t]\) with mesh tending to zero,
deterministic constants \(C_n,V_n\ge0\), almost-sure \(C_n\)-bounds and
\(V_n\)-variation bounds for every stopped/indicator process, and pathwise
continuity of every stopped/indicator piece on \([\bot,t]\).  Then
\[
  N_t=0
  \quad\text{and}\quad
  N_{t-}=0
\]
almost surely.
\end{lemma}

\begin{proof}
Use
Lemma~\ref{lem:Filter.exists_seq_lt_tendsto_nhdsWithin_Iio_of_neBot}
to choose a deterministic sequence \(v_k<t\) tending to \(t\) from the left.
All remaining hypotheses are exactly the hypotheses of
Lemma~\ref{lem:Martingale.eq_zero_and_leftLim_eq_zero_of_localizingSequence_bounded_continuous_finiteVariation_of_variation_bound_continuousOn}.
Apply that lemma with the constructed sequence \(v\).  This gives the terminal
zero and left-limit zero conclusions simultaneously.  No boundedness,
variation, continuity, partition, or mesh input is derived in this lemma; only
the explicit left-approaching sequence is constructed.
\end{proof}

\begin{lemma}\label{lem:IsStronglyPredictable.stronglyMeasurable_left_isolated_past}
  \uses{lem:IsStronglyPredictable.stronglyMeasurable_past}
  \lean{MeasureTheory.IsStronglyPredictable.stronglyMeasurable_left_isolated_past}
Let \(U\) be a strongly predictable real-valued process and let \(t>\bot\).
Assume \(s<t\) is a greatest strict predecessor of \(t\): every \(r<t\)
satisfies \(r\le s\).  Then the random variable \(U_t\) is strongly
measurable with respect to \(\mathcal F_s\).
\end{lemma}

\begin{proof}
Strong predictability gives
\[
  U_t \quad\text{measurable with respect to}\quad
  \bigvee_{r<t}\mathcal F_r
\]
by Lemma~\ref{lem:IsStronglyPredictable.stronglyMeasurable_past}.  The
greatest-predecessor hypothesis implies
\(\mathcal F_r\le \mathcal F_s\) for every \(r<t\), hence
\(\bigvee_{r<t}\mathcal F_r\le \mathcal F_s\).  Monotonicity of strong
measurability with respect to the ambient sigma-algebra then upgrades the
past measurability of \(U_t\) to \(\mathcal F_s\)-measurability.
\end{proof}

\begin{lemma}\label{lem:Martingale.eq_zero_of_predictable_left_isolated_of_previous}
  \uses{def:Martingale, lem:IsStronglyPredictable.stronglyMeasurable_left_isolated_past, lem:Martingale.eq_zero_of_predictable_finiteVariation_past_measurable_zero_increment}
  \lean{MeasureTheory.Martingale.eq_zero_of_predictable_left_isolated_of_previous}
Let \(N\) be a real-valued martingale and a strongly predictable process.
Assume \(s<t\) is a greatest strict predecessor of \(t\), and assume
\(N_s=0\) almost surely.  Then \(N_t=0\) almost surely.
\end{lemma}

\begin{proof}
By
Lemma~\ref{lem:IsStronglyPredictable.stronglyMeasurable_left_isolated_past},
the random variable \(N_t\) is \(\mathcal F_s\)-measurable.  The martingale
adaptedness gives \(\mathcal F_s\)-measurability of \(N_s\), so the increment
\(N_t-N_s\) is also \(\mathcal F_s\)-measurable.  Apply
Lemma~\ref{lem:Martingale.eq_zero_of_predictable_finiteVariation_past_measurable_zero_increment}
to this increment and the inequality \(s\le t\); it yields
\(N_t-N_s=0\) almost surely.  Intersect this full-measure event with the
hypothesis \(N_s=0\) almost surely.  On the intersection, \(N_t=N_s=0\), so
\(N_t=0\) almost surely.
\end{proof}

\begin{lemma}\label{lem:Martingale.eq_zero_of_predictable_bottom_immediate}
  \uses{def:Martingale, lem:Martingale.eq_zero_of_predictable_left_isolated_of_previous}
  \lean{MeasureTheory.Martingale.eq_zero_of_predictable_bottom_immediate}
Let \(N\) be a real-valued martingale and a strongly predictable process.
Let \(t>\bot\), and assume that \(t\) is immediate above the bottom time:
every \(r<t\) satisfies \(r\le \bot\).  If \(N_\bot=0\) pointwise, then
\(N_t=0\) almost surely.
\end{lemma}

\begin{proof}
Apply
Lemma~\ref{lem:Martingale.eq_zero_of_predictable_left_isolated_of_previous}
with the predecessor \(s=\bot\).  The strict predecessor inequality is exactly
\(\bot<t\), and the bottom-immediacy hypothesis says that every strict
predecessor of \(t\) is at most \(s\).  The pointwise normalization
\(N_\bot(\omega)=0\) for every \(\omega\) gives the required almost-sure
previous-time zero statement.  The predecessor lemma then yields
\(N_t=0\) almost surely.
\end{proof}

\begin{lemma}\label{lem:Martingale.eq_zero_and_leftLim_eq_zero_of_localizingSequence_bounded_continuous_finiteVariation_of_variation_bound_original_continuousOn_of_neBot_left}
  \uses{def:Martingale, def:LocallyBoundedVariationOn, lem:stoppedProcess_indicator_continuousOn_Icc, lem:Martingale.eq_zero_and_leftLim_eq_zero_of_localizingSequence_bounded_continuous_finiteVariation_of_variation_bound_continuousOn_of_neBot_left}
  \lean{MeasureTheory.Martingale.eq_zero_and_leftLim_eq_zero_of_localizingSequence_bounded_continuous_finiteVariation_of_variation_bound_original_continuousOn_of_neBot_left}
Let \(N\) be a strongly predictable real-valued true martingale with
\(N_\bot=0\), cadlag paths, and locally bounded variation.  Let
\((\tau_n)\) be a localizing sequence and fix \(t>\bot\).  Assume the
left-neighborhood filter at \(t\) is nontrivial.  Assume the stopped-bound
inputs of
Lemma~\ref{lem:Martingale.eq_zero_and_leftLim_eq_zero_of_localizingSequence_bounded_continuous_finiteVariation_of_variation_bound_continuousOn_of_neBot_left},
except replace the stopped-piece continuity hypothesis by the original-path
hypothesis that every path \(s\mapsto N_s(\omega)\) is continuous on
\([\bot,t]\).  Thus the stopped/indicator pieces still have deterministic
\(C_n\)-bounds and \(V_n\)-variation bounds almost surely, and the deterministic
partitions and mesh hypotheses remain explicit.  Then \(N_t=0\) and
\(N_{t-}=0\) almost surely.
\end{lemma}

\begin{proof}
For each \(n\) and \(\omega\), apply
Lemma~\ref{lem:stoppedProcess_indicator_continuousOn_Icc} to the original
continuity of \(s\mapsto N_s(\omega)\) on \([\bot,t]\), with stopping time
\(\tau_n\).  This supplies the stopped-piece continuity hypothesis required by
Lemma~\ref{lem:Martingale.eq_zero_and_leftLim_eq_zero_of_localizingSequence_bounded_continuous_finiteVariation_of_variation_bound_continuousOn_of_neBot_left}.
All stopped-piece bounds, stopped-piece variation bounds, partitions, mesh, and
the nontrivial-left-filter hypothesis are unchanged.  Applying that lemma gives
the terminal zero and left-limit zero conclusions.
\end{proof}

\begin{lemma}\label{lem:Martingale.eq_zero_of_localizingSequence_bounded_continuous_finiteVariation_of_variation_bound_original_continuousOn_of_left_branch}
  \uses{def:Martingale, def:LocallyBoundedVariationOn, lem:Martingale.eq_zero_and_leftLim_eq_zero_of_localizingSequence_bounded_continuous_finiteVariation_of_variation_bound_original_continuousOn_of_neBot_left, lem:Martingale.eq_zero_of_predictable_bottom_immediate, lem:Martingale.eq_zero_of_predictable_left_isolated_of_previous}
  \lean{MeasureTheory.Martingale.eq_zero_of_localizingSequence_bounded_continuous_finiteVariation_of_variation_bound_original_continuousOn_of_left_branch}
Let \(N\) be a strongly predictable real-valued true martingale with
\(N_\bot=0\), cadlag paths, and locally bounded variation.  Let
\((\tau_n)\) be a localizing sequence and fix \(t>\bot\).  Assume the
stopped-bound inputs of
Lemma~\ref{lem:Martingale.eq_zero_and_leftLim_eq_zero_of_localizingSequence_bounded_continuous_finiteVariation_of_variation_bound_original_continuousOn_of_neBot_left}:
deterministic monotone partitions of \([\bot,t]\) with mesh tending to zero,
deterministic constants \(C_n,V_n\ge0\), stopped/indicator \(C_n\)-bounds and
\(V_n\)-variation bounds almost surely, and original-path continuity on
\([\bot,t]\).  Assume also one of the following explicit left-branch
hypotheses holds:
\[
  \mathcal N[t^-]\neq\bot,
  \qquad
  \forall r<t,\ r\le\bot,
  \qquad\text{or}\qquad
  \exists s<t,\ \bigl(\forall r<t,\ r\le s\bigr)\ \text{and}\ N_s=0
  \text{ a.s.}
\]
Then \(N_t=0\) almost surely.
\end{lemma}

\begin{proof}
Split into the three stated branches.  In the nontrivial-left-filter branch,
apply
Lemma~\ref{lem:Martingale.eq_zero_and_leftLim_eq_zero_of_localizingSequence_bounded_continuous_finiteVariation_of_variation_bound_original_continuousOn_of_neBot_left}
with the stopped-piece bounds, stopped-piece variation bounds, original
continuity, partitions and mesh, and take its first conclusion \(N_t=0\).

In the bottom-immediate branch, apply
Lemma~\ref{lem:Martingale.eq_zero_of_predictable_bottom_immediate}.  The
normalization \(N_\bot=0\) supplies the previous zero statement, and the
branch hypothesis supplies the explicit bottom-immediacy condition.

In the left-isolated successor branch, choose the stated predecessor \(s\).
The branch hypothesis supplies \(s<t\), the greatest-predecessor condition,
and \(N_s=0\) almost surely.  Apply
Lemma~\ref{lem:Martingale.eq_zero_of_predictable_left_isolated_of_previous}
to propagate zero from \(s\) to \(t\).  No branch is inferred from topology;
the trichotomy is an explicit hypothesis of this connector.
\end{proof}

\begin{lemma}\label{lem:Martingale.eq_zero_of_localizingSequence_of_bound_variation_bound_original_continuousOn_of_left_branch}
  \uses{def:Martingale, def:LocallyBoundedVariationOn, lem:ae_stoppedProcess_indicator_bound_variation_on_Icc, lem:Martingale.eq_zero_of_localizingSequence_bounded_continuous_finiteVariation_of_variation_bound_original_continuousOn_of_left_branch}
  \lean{MeasureTheory.Martingale.eq_zero_of_localizingSequence_of_bound_variation_bound_original_continuousOn_of_left_branch}
Let \(N\) be a strongly predictable real-valued true martingale with
\(N_\bot=0\), cadlag paths, and locally bounded variation.  Let
\((\tau_n)\) be a localizing sequence and fix \(t>\bot\).  Assume deterministic
monotone partitions of \([\bot,t]\) with mesh tending to zero, original-path
continuity on \([\bot,t]\), and deterministic constants \(C_n,V_n\ge0\).
Assume that, for each \(n\), the original process satisfies the almost-sure
horizon bound
\[
  \forall s\in[\bot,t],\quad |N_s|\le C_n
\]
and the almost-sure variation bound
\[
  \operatorname{Var}_{[\bot,t]}(N)\le V_n .
\]
Assume also the explicit left-branch disjunction from
Lemma~\ref{lem:Martingale.eq_zero_of_localizingSequence_bounded_continuous_finiteVariation_of_variation_bound_original_continuousOn_of_left_branch}:
nontrivial left-neighborhood filter, bottom-immediacy, or a greatest strict
predecessor \(s<t\) with \(N_s=0\) almost surely.  Then \(N_t=0\) almost surely.
\end{lemma}

\begin{proof}
For each \(n\), apply
Lemma~\ref{lem:ae_stoppedProcess_indicator_bound_variation_on_Icc} to the
stopping time \(\tau_n\).  The original horizon bound and original variation
bound on \([\bot,t]\), together with \(C_n\ge0\), give the stopped/indicator
horizon bound and stopped/indicator variation bound required by
Lemma~\ref{lem:Martingale.eq_zero_of_localizingSequence_bounded_continuous_finiteVariation_of_variation_bound_original_continuousOn_of_left_branch}.
Now invoke that branch connector with the original-path continuity, the same
deterministic partitions and mesh, and the supplied branch disjunction.  The
conclusion is exactly \(N_t=0\) almost surely.  This lemma does not construct
the deterministic bounds, variation levels, partitions, mesh, or branch
trichotomy; they remain explicit hypotheses.
\end{proof}

\begin{lemma}\label{lem:Martingale.eq_zero_of_localizingSequence_of_pre_stop_bound_variation_bound_original_continuousOn_of_left_branch}
  \uses{def:Martingale, def:LocallyBoundedVariationOn, lem:ae_stoppedProcess_indicator_bound_variation_on_Icc_of_pre_stop_bound, lem:Martingale.eq_zero_of_localizingSequence_bounded_continuous_finiteVariation_of_variation_bound_original_continuousOn_of_left_branch}
  \lean{MeasureTheory.Martingale.eq_zero_of_localizingSequence_of_pre_stop_bound_variation_bound_original_continuousOn_of_left_branch}
Let \(N\) be a strongly predictable real-valued true martingale with
\(N_\bot=0\), cadlag paths, and locally bounded variation.  Let
\((\tau_n)\) be a localizing sequence and fix \(t>\bot\).  Assume deterministic
monotone partitions of \([\bot,t]\) with mesh tending to zero, original-path
continuity on \([\bot,t]\), and deterministic constants \(C_n,V_n\ge0\).
For each \(n\), assume the following three almost-sure inputs:
\[
  r\in[\bot,t],\quad r<\tau_n(\omega)
  \quad\Longrightarrow\quad
  |N_r(\omega)|\le C_n,
\]
\[
  \bot<\tau_n(\omega),\quad \tau_n(\omega)\ne\top,\quad
  \tau_n(\omega)\le (t:\iota_\infty)
  \quad\Longrightarrow\quad
  |N_{\tau_n(\omega)}(\omega)|\le C_n,
\]
and \(N(\cdot,\omega)\) has bounded variation at most \(V_n\) on the closed
pre-stop horizon set
\[
  S_{\tau_n,t}(\omega)
  =
  \{r\in[\bot,t] : (r:\iota_\infty)\le \tau_n(\omega)\}.
\]
Assume also the explicit left-branch disjunction from
Lemma~\ref{lem:Martingale.eq_zero_of_localizingSequence_bounded_continuous_finiteVariation_of_variation_bound_original_continuousOn_of_left_branch}:
nontrivial left-neighborhood filter, bottom-immediacy, or a greatest strict
predecessor \(s<t\) with \(N_s=0\) almost surely.  Then \(N_t=0\) almost surely.
\end{lemma}

\begin{proof}
For each \(n\), apply
Lemma~\ref{lem:ae_stoppedProcess_indicator_bound_variation_on_Icc_of_pre_stop_bound}
to the stopping time \(\tau_n\).  The pre-stop horizon bound, finite
stop-value bound, and closed-pre-stop variation bound give the stopped/indicator
\(C_n\)-horizon bound and \(V_n\)-variation bound on \([\bot,t]\).  These are
exactly the stopped-bound inputs required by
Lemma~\ref{lem:Martingale.eq_zero_of_localizingSequence_bounded_continuous_finiteVariation_of_variation_bound_original_continuousOn_of_left_branch}.
Invoke that branch connector with the same original-path continuity,
deterministic partitions, mesh hypothesis, and supplied branch disjunction.
This proves \(N_t=0\) almost surely.  No deterministic bound, variation level,
partition, mesh, branch alternative, or no-overshoot property is inferred; all
such data is assumed explicitly.
\end{proof}

\begin{lemma}\label{lem:Martingale.eq_zero_of_localizingSequence_of_pre_stop_bound_variation_bound_original_continuousOn_of_dense_left}
  \uses{def:Martingale, def:LocallyBoundedVariationOn, lem:Filter.nhdsWithin_Iio_self_neBot_of_bot_lt, lem:Martingale.eq_zero_of_localizingSequence_of_pre_stop_bound_variation_bound_original_continuousOn_of_left_branch}
  \lean{MeasureTheory.Martingale.eq_zero_of_localizingSequence_of_pre_stop_bound_variation_bound_original_continuousOn_of_dense_left}
Let \(N\) be a strongly predictable real-valued true martingale with
\(N_\bot=0\), cadlag paths, and locally bounded variation, indexed by a dense
order with bottom.  Let \((\tau_n)\) be a localizing sequence and fix
\(t>\bot\).  Assume deterministic monotone partitions of \([\bot,t]\) with
mesh tending to zero, original-path continuity on \([\bot,t]\), and
deterministic constants \(C_n,V_n\ge0\).  For each \(n\), assume exactly the
same three almost-sure pre-stop inputs as in
Lemma~\ref{lem:Martingale.eq_zero_of_localizingSequence_of_pre_stop_bound_variation_bound_original_continuousOn_of_left_branch}:
a \(C_n\)-bound before \(\tau_n\), a \(C_n\)-bound at the finite stopped value
when \(\tau_n(\omega)\le t\), and a \(V_n\)-bound on the variation over the
closed pre-stop horizon set
\[
  \{r\in[\bot,t] : (r:\iota_\infty)\le \tau_n(\omega)\}.
\]
Then \(N_t=0\) almost surely.
\end{lemma}

\begin{proof}
Because the time order is dense and \(\bot<t\),
Lemma~\ref{lem:Filter.nhdsWithin_Iio_self_neBot_of_bot_lt} gives a
nontrivial strict-left neighborhood filter at \(t\).  Therefore the explicit
branch disjunction required by
Lemma~\ref{lem:Martingale.eq_zero_of_localizingSequence_of_pre_stop_bound_variation_bound_original_continuousOn_of_left_branch}
holds by its nontrivial-left alternative.  Apply that pre-stop localizing
left-branch wrapper with the same localizing sequence, pre-stop bounds,
finite stop-value bounds, closed-pre-stop variation bounds, original-path
continuity, deterministic partitions, and mesh.  This dense-left
specialization constructs only the branch datum; it does not infer
no-overshoot, deterministic bounds, variation levels, continuity, partitions,
mesh, or localizing stopping times.
\end{proof}

\begin{lemma}\label{lem:Martingale.eq_zero_of_bound_variation_bound_original_continuousOn_of_left_branch}
  \uses{def:Martingale, def:LocallyBoundedVariationOn, lem:Martingale.eq_zero_of_localizingSequence_of_bound_variation_bound_original_continuousOn_of_left_branch}
  \lean{MeasureTheory.Martingale.eq_zero_of_bound_variation_bound_original_continuousOn_of_left_branch}
Let \(N\) be a strongly predictable real-valued true martingale with
\(N_\bot=0\), cadlag paths, and locally bounded variation.  Fix
\(t>\bot\).  Assume deterministic monotone partitions of \([\bot,t]\) whose
mesh tends to zero, original-path continuity on \([\bot,t]\), and single
deterministic constants \(C,V\ge0\).  Assume that, almost surely,
\[
  \forall s\in[\bot,t],\quad |N_s|\le C,
\]
and that, almost surely, \(N\) has bounded variation on \([\bot,t]\) with
total variation at most \(V\).  Assume also the same explicit left-branch
disjunction as in
Lemma~\ref{lem:Martingale.eq_zero_of_localizingSequence_of_bound_variation_bound_original_continuousOn_of_left_branch}.
Then \(N_t=0\) almost surely.
\end{lemma}

\begin{proof}
Use the constant localizing sequence \(\tau_n\equiv\top\), which is a
localizing sequence because every stopping time is identically infinite and
the sequence trivially tends to \(\top\).  Define the deterministic level
families \(C_n=C\) and \(V_n=V\).  The single almost-sure horizon bound and
the single almost-sure variation bound provide the \(n\)-indexed hypotheses
required by
Lemma~\ref{lem:Martingale.eq_zero_of_localizingSequence_of_bound_variation_bound_original_continuousOn_of_left_branch}
for every \(n\).  Apply that lemma with the same deterministic partitions,
mesh hypothesis, original-path continuity, and branch disjunction.  This
argument only packages the already supplied deterministic levels; it does not
construct bounds, variation levels, partitions, mesh, or branch alternatives.
\end{proof}

\begin{lemma}\label{lem:Filter.exists_greatest_lt_of_not_neBot_nhdsWithin_Iio}
  \lean{Filter.exists_greatest_lt_of_not_neBot_nhdsWithin_Iio}
Let \(t>\bot\) in a linearly ordered topological space with the order
topology.  If the left-neighborhood filter
\(\mathcal N[t^-]=\mathcal N_t\!\restriction(-\infty,t)\) is not nontrivial,
then \(t\) has a greatest strict predecessor: there exists \(s<t\) such that
every \(r<t\) satisfies \(r\le s\).
\end{lemma}

\begin{proof}
Since the within-filter is not nontrivial, it is the bottom filter; equivalently
the empty set is a left-neighborhood of \(t\).  Thus there is an ordinary
neighborhood \(U\) of \(t\) such that \(U\cap(-\infty,t)\) is empty.  Because
\(\bot<t\), the order-topology neighborhood-basis lemma
\texttt{exists\_Ioc\_subset\_of\_mem\_nhds} gives some \(s<t\) with
\((s,t]\subseteq U\).  If \(r<t\) and \(s<r\), then \(r\in(s,t]\subseteq U\)
and also \(r\in(-\infty,t)\), contradicting the emptiness of
\(U\cap(-\infty,t)\).  Therefore every \(r<t\) has \(r\le s\), so \(s\) is the
greatest strict predecessor of \(t\).
\end{proof}

\begin{lemma}\label{lem:Martingale.eq_zero_of_bound_variation_bound_original_continuousOn_of_strictPast_zero}
  \uses{def:Martingale, def:LocallyBoundedVariationOn, lem:Filter.exists_greatest_lt_of_not_neBot_nhdsWithin_Iio, lem:Martingale.eq_zero_of_bound_variation_bound_original_continuousOn_of_left_branch}
  \lean{MeasureTheory.Martingale.eq_zero_of_bound_variation_bound_original_continuousOn_of_strictPast_zero}
Let \(N\) be a strongly predictable real-valued true martingale with
\(N_\bot=0\), cadlag paths, and locally bounded variation.  Fix \(t>\bot\).
Assume the same bounded-level analytic inputs as in
Lemma~\ref{lem:Martingale.eq_zero_of_bound_variation_bound_original_continuousOn_of_left_branch}:
deterministic monotone partitions of \([\bot,t]\) with mesh tending to zero,
original-path continuity on \([\bot,t]\), constants \(C,V\ge0\), an
almost-sure \(C\)-bound on \(N\) over \([\bot,t]\), and an almost-sure
\(V\)-bound on the variation over \([\bot,t]\).  Finally assume the explicit
strict-past zero hypothesis
\[
  \forall s<t,\qquad N_s=0\quad\text{almost surely}.
\]
Then \(N_t=0\) almost surely.
\end{lemma}

\begin{proof}
Split on whether the left-neighborhood filter at \(t\) is nontrivial.  In the
nontrivial-left case, the explicit branch disjunction required by
Lemma~\ref{lem:Martingale.eq_zero_of_bound_variation_bound_original_continuousOn_of_left_branch}
holds by its first alternative.

In the complementary case, apply
Lemma~\ref{lem:Filter.exists_greatest_lt_of_not_neBot_nhdsWithin_Iio} to obtain
a greatest strict predecessor \(s<t\).  The strict-past zero hypothesis gives
\(N_s=0\) almost surely.  Hence the explicit branch disjunction required by
Lemma~\ref{lem:Martingale.eq_zero_of_bound_variation_bound_original_continuousOn_of_left_branch}
holds by its predecessor alternative.  Apply that bounded-level endpoint with
the same bounds, variation bound, continuity, partitions, and mesh.  This
lemma constructs only the branch data from the extra strict-past zero
hypothesis; it does not construct deterministic bounds, variation levels,
partitions, or mesh, and it is not a full induction principle.
\end{proof}

\begin{lemma}\label{lem:Filter.nhdsWithin_Iio_self_neBot_of_bot_lt}
  \lean{Filter.nhdsWithin_Iio_self_neBot_of_bot_lt}
Let \(t>\bot\) in a densely ordered linearly ordered topological space with
the order topology and a bottom element.  Then the strict-left neighborhood
filter at \(t\),
\[
  \mathcal N_t\!\restriction(-\infty,t),
\]
is nontrivial.
\end{lemma}

\begin{proof}
Since \(\bot<t\), the interval \((-\infty,t)\) is nonempty.  Mathlib's
dense-order closure lemma `nhdsWithin_Iio_neBot'` says that, in a densely
ordered order-topological space, if \((-\infty,c)\) is nonempty and \(b\le c\),
then \(\mathcal N_b\!\restriction(-\infty,c)\) is nontrivial.  Apply it with
\(b=c=t\), use the witness \(\bot\in(-\infty,t)\), and use \(t\le t\).
Equivalently, one may use the specialized Mathlib lemma
`nhdsLT_neBot_of_exists_lt` with the witness \(\bot<t\).
\end{proof}

\begin{lemma}\label{lem:Martingale.eq_zero_of_localizingSequence_of_bound_variation_bound_original_continuousOn_of_dense_left}
  \uses{def:Martingale, def:LocallyBoundedVariationOn, lem:Filter.nhdsWithin_Iio_self_neBot_of_bot_lt, lem:Martingale.eq_zero_of_localizingSequence_of_bound_variation_bound_original_continuousOn_of_left_branch}
  \lean{MeasureTheory.Martingale.eq_zero_of_localizingSequence_of_bound_variation_bound_original_continuousOn_of_dense_left}
Let \(N\) be a strongly predictable real-valued true martingale with
\(N_\bot=0\), cadlag paths, and locally bounded variation, indexed by a
dense order with bottom.  Fix \(t>\bot\) and a localizing sequence.  Assume
the same original-process indexed horizon bounds, indexed variation bounds,
original-path continuity on \([\bot,t]\), deterministic monotone partitions,
endpoints, interval membership, and mesh as in
Lemma~\ref{lem:Martingale.eq_zero_of_localizingSequence_of_bound_variation_bound_original_continuousOn_of_left_branch}.
Then \(N_t=0\) almost surely.
\end{lemma}

\begin{proof}
By Lemma~\ref{lem:Filter.nhdsWithin_Iio_self_neBot_of_bot_lt}, the
strict-left neighborhood filter at \(t\) is nontrivial.  Therefore the branch
disjunction required by
Lemma~\ref{lem:Martingale.eq_zero_of_localizingSequence_of_bound_variation_bound_original_continuousOn_of_left_branch}
holds by its first alternative.  Apply that lemma with the same localizing
sequence, indexed deterministic levels, bounds, variation bounds, continuity,
partitions, and mesh.  This result constructs only the dense-left branch
input; it does not derive any deterministic bound, variation level, partition,
mesh, or zero-propagation fact.
\end{proof}

\begin{lemma}\label{lem:Martingale.eq_zero_of_bound_variation_bound_original_continuousOn_of_dense_left}
  \uses{def:Martingale, def:LocallyBoundedVariationOn, lem:Filter.nhdsWithin_Iio_self_neBot_of_bot_lt, lem:Martingale.eq_zero_of_bound_variation_bound_original_continuousOn_of_left_branch}
  \lean{MeasureTheory.Martingale.eq_zero_of_bound_variation_bound_original_continuousOn_of_dense_left}
Let \(N\) be a strongly predictable real-valued true martingale with
\(N_\bot=0\), cadlag paths, and locally bounded variation, indexed by a
dense order with bottom.  Fix \(t>\bot\).  Assume the same bounded-level
analytic inputs as in
Lemma~\ref{lem:Martingale.eq_zero_of_bound_variation_bound_original_continuousOn_of_left_branch}:
constants \(C,V\ge0\), an almost-sure \(C\)-bound on \(N\) over
\([\bot,t]\), an almost-sure \(V\)-bound on variation over \([\bot,t]\),
original-path continuity on \([\bot,t]\), deterministic monotone partitions,
endpoints, interval membership, and mesh.  Then \(N_t=0\) almost surely.
\end{lemma}

\begin{proof}
By Lemma~\ref{lem:Filter.nhdsWithin_Iio_self_neBot_of_bot_lt}, the
strict-left neighborhood filter at \(t\) is nontrivial.  Hence the explicit
branch disjunction required by
Lemma~\ref{lem:Martingale.eq_zero_of_bound_variation_bound_original_continuousOn_of_left_branch}
holds by the nontrivial-left alternative.  Apply that bounded-level endpoint
with the same bounds, variation bound, continuity, partitions, and mesh.
Again, this is only branch-data packaging for dense time indices.
\end{proof}

\begin{lemma}\label{lem:Martingale.integral_sq_terminal_eq_zero_of_bounded_continuous_finiteVariation}
  \uses{def:Martingale, def:LocallyBoundedVariationOn, lem:Martingale.integral_sq_terminal_eq_zero_of_bounded_continuous_finiteVariation_of_variation_bound_continuousOn, lem:Martingale.eq_zero_of_bounded_continuous_finiteVariation_of_variation_bound_continuousOn}
  \lean{MeasureTheory.Martingale.integral_sq_terminal_eq_zero_of_bounded_continuous_finiteVariation}
Let \(M\) be a real-valued true martingale with \(M_\bot=0\).  Fix a
deterministic time \(t>\bot\).  Assume that on the deterministic horizon
\([\bot,t]\) the paths are continuous, have bounded variation, and are bounded
by a deterministic integrable constant.  Assume moreover that the proof has
chosen deterministic partitions of \([\bot,t]\) with adjacent-pair mesh tending
to zero in the time uniformity.  Then \(M_t^2\) is integrable and
\[
  \int M_t^2\,dP=0 .
\]
\end{lemma}

\begin{proof}
Choose deterministic finite partitions
\(\pi_n=\{\bot=s^n_0\le \cdots \le s^n_{k_n}=t\}\) whose mesh tends to zero.
For each partition, the telescoping identity gives
\[
  M_t-M_\bot
  =
  \sum_i (M_{s^n_{i+1}}-M_{s^n_i}).
\]

For each sample point, continuity on \([\bot,t]\) implies that the largest
increment along \(\pi_n\) tends to zero.  This modulus depends on the sample
point; no deterministic modulus is available from pathwise continuity alone.
First prove the variation-bounded stopped subcase, where the total variation on
\([\bot,t]\) is also bounded by a deterministic constant.  In that subcase,
convert continuity on the compact interval to uniform continuity and use the
chosen deterministic partitions with adjacent-pair mesh tending to zero.
Lemma~\ref{lem:Martingale.integral_sq_terminal_eq_zero_of_bounded_continuous_finiteVariation_of_variation_bound_continuousOn}
then turns the continuous-path mesh control and deterministic variation bound
into \(\int M_t^2\,dP=0\).  The full locally bounded-variation statement is
obtained by stopping or localizing at increasing variation levels and then
letting the level tend to infinity on the event where the localized process
agrees with the original process up to \(t\).  The deterministic bound on
\(M_t\) gives integrability of \(M_t^2\).
\end{proof}

\begin{lemma}\label{lem:Martingale.eq_zero_of_bounded_continuous_finiteVariation_core}
  \uses{def:Martingale, lem:MeasureTheory.ae_eq_zero_of_integral_sq_eq_zero}
  \lean{MeasureTheory.Martingale.eq_zero_of_bounded_continuous_finiteVariation_core}
Let \(M\) be a real-valued true martingale with \(M_\bot=0\).  Fix a
deterministic time \(t>\bot\).  Assume that on the deterministic horizon
\([\bot,t]\) the paths are continuous, have bounded variation, and are bounded
by a deterministic integrable constant.  If \(M_t^2\) is integrable and the
square-integral identity \(\int M_t^2\,dP=0\) has been proved, then \(M_t=0\)
almost surely.
\end{lemma}

\begin{proof}
Apply
Lemma~\ref{lem:MeasureTheory.ae_eq_zero_of_integral_sq_eq_zero} to the
integrable random variable \(M_t\).  The continuity, boundedness, martingale,
normalization, and finite-variation hypotheses are retained here because this
terminal lemma is called from the bounded-continuous finite-variation bridge;
the actual proof work left to that bridge is to derive the square-integral
identity from those hypotheses.
\end{proof}


\begin{lemma}\label{lem:Martingale.eq_zero_of_predictable_finiteVariation_bounded_continuous_reduction}
  % iter-062 PIVOT: \uses trimmed to the endpoints the assembly proof actually
  % invokes (dense-order specialization). The jump-removal pipeline is reached
  % through the localizing dense-left endpoint, not lemma-by-lemma here.
  \uses{def:Martingale, def:LocallyBoundedVariationOn, lem:Martingale.eq_zero_of_localizingSequence_of_bound_variation_bound_original_continuousOn_of_dense_left, lem:Martingale.eq_zero_of_bound_variation_bound_original_continuousOn_of_dense_left, lem:Martingale.eq_zero_and_leftLim_eq_zero_of_localizingSequence_of_bound_variation_bound_original_continuousOn, lem:Martingale.integral_increment_mul_increment_eq_zero, lem:Martingale.eq_zero_of_bounded_continuous_finiteVariation_core}
  \lean{MeasureTheory.Martingale.eq_zero_of_predictable_finiteVariation_bounded_continuous_reduction}
Let \(M\) be a strongly predictable càdlàg true martingale with
\(M_\bot=0\), and suppose every path has locally bounded variation.  For
every deterministic time \(t>\bot\), the remaining analytic reduction proves
\(M_t=0\) almost surely.
\end{lemma}

% PIVOT (iter-062): the route stopped adding general-`κ` branch/pre-stop helpers.
% A strategy-critic + mathlib-analogist reassessment found that this `_reduction`
% is NOT provable over an arbitrary linearly ordered `κ`: closing it requires a
% deterministic refining partition of `[⊥,t]` with vanishing mesh, which does not
% exist in a general order. Every downstream endpoint already takes the partition
% and mesh as EXPLICIT hypotheses precisely to defer this construction; `_reduction`
% is where the buck stops. The fix is to SPECIALIZE to densely-ordered time
% (`[DenselyOrdered κ]` with the topological instances the endpoints already carry,
% e.g. `ℝ≥0`, where dyadic partitions + Heine–Cantor supply the mesh) and ASSEMBLE
% the already-closed pieces — not to prove more helpers. The jump-removal and
% continuous-localizing endpoints below are all closed; the only remaining content
% is the partition construction and wiring.
\begin{proof}
Specialize to densely-ordered time so that a deterministic refining partition of
the compact horizon \([\bot,t]\) with vanishing mesh can be constructed (on
\(\mathbb R_{\ge 0}\) take dyadic subdivisions; uniform continuity of the paths on
the compact interval gives the vanishing-mesh control by Heine–Cantor).  With that
partition and the deterministic a.e.\ horizon and variation bounds available, the
proof is pure assembly of closed lemmas; no new branch, pre-stop, or mesh helper is
needed.

\emph{Step 1 — bounds and partition.}  From the càdlàg paths on the compact
horizon \([\bot,t]\) obtain an a.e.\ deterministic horizon bound \(\sup_{s\in[\bot,t]}\|M_s\|\le C\),
and from locally bounded variation obtain an a.e.\ deterministic variation bound
\((\operatorname{Var}_{[\bot,t]} M)\le V\).  Construct the deterministic refining
partition \(u_n\) of \([\bot,t]\) and verify the vanishing-mesh hypothesis.

\emph{Step 2 — jump removal gives continuity on the localized pieces.}  The
predictable-jump removal pipeline is already closed and supplies, along a localizing
sequence, that the stopped/indicator pieces have no predictable jump, i.e.\ are
a.e.\ continuous on the horizon.  This is exactly the input the continuous-localizing
endpoint consumes; the existing jump-removal wrappers (terminal-and-left-limit zero
packages) feed it.

\emph{Step 3 — apply the closed continuous dense-left endpoint.}  Feed the bounds,
variation bound, partition, mesh, and continuity into
Lemma~\ref{lem:Martingale.eq_zero_of_localizingSequence_of_bound_variation_bound_original_continuousOn_of_dense_left}
(or, with a single localization level, its fixed-level dense-left companion
Lemma~\ref{lem:Martingale.eq_zero_of_bound_variation_bound_original_continuousOn_of_dense_left}).
Under \([\text{DenselyOrdered }\kappa]\) the branch disjunction these endpoints
internally require is always the nontrivial strict-left branch, so the
bottom-immediate and left-isolated cases — and the entire `leastGT` pre-stop layer
the route was about to extend — are vacuous and unused.  The endpoint already runs
the deterministic-partition square-increment estimate
\(\mathbb E[M_t^2]=\sum_i \mathbb E[(M_{t_{i+1}}-M_{t_i})^2]\le \mathbb E[\max_i|\Delta M|\cdot V]\to 0\)
(increment orthogonality from
Lemma~\ref{lem:Martingale.integral_increment_mul_increment_eq_zero}) and the terminal
core Lemma~\ref{lem:Martingale.eq_zero_of_bounded_continuous_finiteVariation_core},
yielding \(M_t=0\) almost surely.
\end{proof}


\begin{lemma}\label{lem:Martingale.eq_zero_of_predictable_finiteVariation_bounded_continuous}
  \uses{def:Martingale, def:LocallyBoundedVariationOn, lem:Martingale.eq_zero_of_predictable_finiteVariation_bounded_continuous_reduction}
  \lean{MeasureTheory.Martingale.eq_zero_of_predictable_finiteVariation_bounded_continuous}
Let \(M\) be a strongly predictable càdlàg true martingale with \(M_0=0\),
and suppose every path has locally bounded variation.  For every
deterministic time \(t>0\), the bounded-continuous finite-variation bridge
shows that \(M_t=0\) almost surely.
\end{lemma}

\begin{proof}
The Lean statement still carries the full predictable càdlàg finite-variation
hypotheses; the phrase ``bounded-continuous'' names the core bridge reached
inside the proof, not an extra assumption in the statement.
It is now a wrapper around
Lemma~\ref{lem:Martingale.eq_zero_of_predictable_finiteVariation_bounded_continuous_reduction}.
\end{proof}


\begin{lemma}\label{lem:Martingale.eq_zero_of_predictable_finiteVariation_value_zero}
  \uses{def:Martingale, def:LocallyBoundedVariationOn, lem:Martingale.eq_zero_of_predictable_finiteVariation_past_measurable_zero_increment, lem:LocallyBoundedVariationOn.exists_tendsto_left_univ, lem:Martingale.stronglyMeasurable_leftLim_past_of_left_approach, lem:Martingale.stronglyMeasurable_jump_leftLim_past_of_left_approach, lem:Filtration.ae_eq_zero_of_past_setIntegral_eq_zero, lem:Martingale.ae_eq_leftLim_of_left_approach_past_setIntegral_zero, lem:Martingale.setIntegral_increment_eq_zero_of_measurableSet, lem:Martingale.setIntegral_jump_leftLim_eq_zero_of_left_approach_of_dominated, lem:Martingale.ae_eq_leftLim_of_left_approach_past_measurable_increments, lem:Martingale.eq_zero_of_predictable_finiteVariation_bounded_continuous}
  \lean{MeasureTheory.Martingale.eq_zero_of_predictable_finiteVariation_value_zero}
Let \(M\) be a strongly predictable càdlàg true martingale with \(M_0=0\),
and suppose every path has locally bounded variation.  For every
deterministic time \(t>0\), \(M_t=0\) almost surely.
\end{lemma}

\begin{proof}
First remove deterministic predictable jumps.  By
Lemma~\ref{lem:LocallyBoundedVariationOn.exists_tendsto_left_univ}, each
sample path has a left limit at \(t\).  The closed martingale increment
Lemma~\ref{lem:Martingale.eq_zero_of_predictable_finiteVariation_past_measurable_zero_increment}
says that an increment \(M_t-M_s\) vanishes almost surely whenever it is
already \(\mathcal F_s\)-measurable.  Use strong predictability only through
the actual predictable-section or past-sigma-algebra measurability it gives;
do not replace it by measurability with respect to an arbitrary fixed
\(\mathcal F_s\) without proving that measurability.  Instead, along a
deterministic cofinal left-approaching sequence, prove
\(M_{t-}\) is past-measurable by
Lemma~\ref{lem:Martingale.stronglyMeasurable_leftLim_past_of_left_approach}
and combine it with strong predictability to make the jump
\(M_t-M_{t-}\) past-measurable by
Lemma~\ref{lem:Martingale.stronglyMeasurable_jump_leftLim_past_of_left_approach}.
The remaining jump-removal step is then a conditional-expectation argument
against the past sigma-algebra: it is enough, by
Lemma~\ref{lem:Filtration.ae_eq_zero_of_past_setIntegral_eq_zero}, to show
that the jump has zero integral on every earlier-filtration set, and then
apply
Lemma~\ref{lem:Martingale.ae_eq_leftLim_of_left_approach_past_setIntegral_zero}.
The stronger
Lemma~\ref{lem:Martingale.ae_eq_leftLim_of_left_approach_past_measurable_increments}
is available only when the increments are genuinely
\(\mathcal F_{u_n}\)-measurable.

It remains to handle the continuous finite-variation part on a fixed finite
horizon.  Stop the martingale so it is bounded on the horizon, and work with
a deterministic partition
\(\pi=\{0=s_0<\cdots<s_n=t\}\).  Orthogonality of martingale increments gives
\[
  \mathbb E[(M_t-M_0)^2]
  =
  \sum_{k<n}\mathbb E[(M_{s_{k+1}}-M_{s_k})^2].
\]
For each path of the continuous finite-variation process,
\[
  \sum_{k<n} (M_{s_{k+1}}-M_{s_k})^2
  \le
  \max_{k<n}|M_{s_{k+1}}-M_{s_k}|\,
  \operatorname{Var}_{[0,t]}(M).
\]
Along refining deterministic partitions whose mesh tends to zero, continuity
forces the maximum increment to tend to zero, while finite variation gives a
finite pathwise variation bound.  Hence the square-increment sums tend to
zero pathwise.  In the bounded stopped case, dominated convergence sends the
expectations of these sums to zero, so
\(\mathbb E[(M_t-M_0)^2]=0\).  Since \(M_0=0\), \(M_t=0\) almost surely.
Finally remove the bounded stopping by localization.
\end{proof}


\begin{lemma}\label{lem:Martingale.eq_zero_of_predictable_finiteVariation_past_condExp_zero}
  \uses{def:Martingale, def:LocallyBoundedVariationOn, lem:condExp_ae_eq_zero_of_ae_eq_zero, lem:Martingale.eq_zero_of_predictable_finiteVariation_value_zero}
  \lean{MeasureTheory.Martingale.eq_zero_of_predictable_finiteVariation_past_condExp_zero}
Let \(M\) be a strongly predictable càdlàg true martingale with \(M_0=0\),
and suppose every path has locally bounded variation.  For every
deterministic time \(t>0\),
\[
  \mathbb E[M_t\mid \bigvee_{s<t}\mathcal F_s]=0
\]
almost surely.
\end{lemma}

\begin{proof}
By
Lemma~\ref{lem:Martingale.eq_zero_of_predictable_finiteVariation_value_zero},
\(M_t=0\) almost surely.  Apply
Lemma~\ref{lem:condExp_ae_eq_zero_of_ae_eq_zero} to this almost-sure equality
with the past sigma-algebra \(\bigvee_{s<t}\mathcal F_s\).  The conditional
expectation of \(M_t\) onto that sigma-algebra is therefore zero almost
surely.
\end{proof}


\begin{theorem}\label{thm:Martingale.eq_zero_of_predictable_finiteVariation_noninitial_analytic}
  \uses{def:Martingale, def:LocallyBoundedVariationOn, lem:Martingale.eq_zero_of_predictable_finiteVariation_past_measurable_value, lem:Martingale.eq_zero_of_predictable_finiteVariation_past_condExp_zero}
  \lean{MeasureTheory.Martingale.eq_zero_of_predictable_finiteVariation_noninitial_analytic}
Let \(M\) be a strongly predictable càdlàg martingale with \(M_0=0\).
Assume every path has locally bounded variation, and fix a deterministic time
\(t>0\).  Suppose in addition that \(M_t\) is strongly measurable with
respect to the past sigma-algebra \(\bigvee_{s<t}\mathcal F_s\).  Then
\(M_t=0\) almost surely.
\end{theorem}

\begin{proof}
By
Lemma~\ref{lem:Martingale.eq_zero_of_predictable_finiteVariation_past_condExp_zero},
the conditional expectation of \(M_t\) onto the past sigma-algebra
\(\bigvee_{s<t}\mathcal F_s\) is zero almost surely.  Since the theorem
assumes that \(M_t\) is strongly measurable with respect to this same
past sigma-algebra, apply
Lemma~\ref{lem:Martingale.eq_zero_of_predictable_finiteVariation_past_measurable_value}
to conclude \(M_t=0\) almost surely.
\end{proof}


\begin{theorem}\label{thm:Martingale.eq_zero_of_predictable_finiteVariation_noninitial}
  \uses{def:Martingale, def:LocallyBoundedVariationOn, lem:IsStronglyPredictable.stronglyMeasurable_past, thm:Martingale.eq_zero_of_predictable_finiteVariation_noninitial_analytic}
  \lean{MeasureTheory.Martingale.eq_zero_of_predictable_finiteVariation_noninitial}
Let \(M\) be a càdlàg martingale with \(M_0=0\).  Assume that \(M\) is
strongly predictable and that every path has locally bounded variation.
Then, for every deterministic time \(t\) with \(0<t\), \(M_t=0\) almost
surely.
\end{theorem}

\begin{proof}
Fix \(t>0\).  By
Lemma~\ref{lem:IsStronglyPredictable.stronglyMeasurable_past}, strong
predictability gives exactly the hypothesis that \(M_t\) is strongly
measurable with respect to \(\bigvee_{s<t}\mathcal F_s\).  Apply
Theorem~\ref{thm:Martingale.eq_zero_of_predictable_finiteVariation_noninitial_analytic}
with this past-measurability input.
\end{proof}


\begin{theorem}\label{thm:Martingale.eq_zero_of_predictable_finiteVariation}
  \uses{def:Martingale, def:LocallyBoundedVariationOn, thm:Martingale.eq_zero_of_predictable_finiteVariation_noninitial}
  \lean{MeasureTheory.Martingale.eq_zero_of_predictable_finiteVariation}
Let \(M\) be a càdlàg martingale with \(M_0=0\).  Assume that \(M\) is
strongly predictable and that every path has locally bounded variation.
Then, for every deterministic time \(t\), \(M_t=0\) almost surely.
\end{theorem}

\begin{proof}
Fix \(t\).  If \(t=0\), the conclusion is exactly the assumed normalization
\(M_0=0\).  If \(0<t\), apply
Theorem~\ref{thm:Martingale.eq_zero_of_predictable_finiteVariation_noninitial}
at this time.  These two cases give the fixed-time conclusion for every
deterministic \(t\).
\end{proof}


\begin{theorem}\label{thm:IsLocalMartingale.eq_zero_of_predictable_finiteVariation}
  \uses{def:IsLocalMartingale, def:LocallyBoundedVariationOn, lem:IsStronglyPredictable_stoppedProcess_indicator, lem:LocallyBoundedVariationOn_stoppedProcess_indicator, thm:Martingale.eq_zero_of_predictable_finiteVariation}
  \lean{ProbabilityTheory.IsLocalMartingale.eq_zero_of_predictable_finiteVariation}
  \leanok
Let \(N\) be a càdlàg local martingale with \(N_0=0\).  Assume that \(N\)
is strongly predictable and that each path of \(N\) has locally bounded
variation on compact time intervals.  Then, for every deterministic time
\(t\), \(N_t=0\) almost surely.
\end{theorem}

\begin{proof}

\leanok
Fix a deterministic time \(t\).  If \(t=0\), the claim is the normalization
assumption \(N_0=0\).  Otherwise choose a localizing sequence \((\tau_n)\)
for \(N\).  For each \(n\), the stopped and indicator-truncated process
\[
  N^n_s
  =
  N_{s\wedge\tau_n}\mathbf 1_{\{\tau_n>0\}}
\]
is a càdlàg martingale.  Strong predictability is stable under stopping and
indicator truncation by a stopping-time event, and local bounded variation is
pathwise stable under stopping.  Thus it suffices to prove the theorem for a
true càdlàg predictable martingale of locally bounded variation, covered by
Theorem~\ref{thm:Martingale.eq_zero_of_predictable_finiteVariation}.

After applying that theorem to each \(N^n\), take the countable intersection
of the resulting a.e. equalities and the a.e. event on which
\(\tau_n\uparrow\infty\).  For such an \(\omega\) and fixed \(t>0\), there is
some \(n\) with \(t<\tau_n(\omega)\), so the stopped-indicator process agrees
with \(N_t(\omega)\).  Hence \(N_t=0\) almost surely.  If \(t=0\), the claim
is exactly the normalization \(N_0=0\).

The finite-variation and predictability assumptions are both essential:
continuous martingales such as Brownian motion are predictable but not of
finite variation, and càdlàg finite-variation martingales may have totally
inaccessible jumps when predictability is omitted.
\end{proof}


\begin{lemma}\label{lem:predictablePart_eq_of_normalized_decomposition}
  \uses{thm:local_doobMeyer_normalized, thm:IsLocalMartingale.eq_zero_of_predictable_finiteVariation}
  \lean{ProbabilityTheory.IsLocalSubmartingale.predictablePart_eq_of_normalized_decomposition}
  \leanok
Suppose \(X=M+A\), where \(M\) is a càdlàg local martingale and \(A\) is
strongly predictable, strongly progressive, càdlàg, locally integrable,
nondecreasing, and \(A_0=0\).  Then, for every deterministic time \(t\),
the choice-based predictable part of the normalized Doob--Meyer
decomposition of \(X\) is almost surely equal to \(A_t\).
\end{lemma}

\begin{proof}

\leanok
Let \(X=M'+A'\) be the normalized chosen decomposition.  Since both
decompositions have the same sum, \(A'-A=M-M'\).  The right-hand side is
a local martingale, being the difference of two local martingales.  The
left-hand side is strongly predictable, because both \(A'\) and \(A\) are
strongly predictable.  It is of finite variation on compact intervals,
since each nondecreasing real-valued path has locally bounded variation and
locally bounded variation is stable under subtraction.  It also starts at
zero because both decompositions are normalized.  Therefore
Theorem~\ref{thm:IsLocalMartingale.eq_zero_of_predictable_finiteVariation}
implies \(A'_t-A_t=0\) almost surely for every deterministic time \(t\),
and hence \(A'_t=A_t\) almost surely for every \(t\).
\end{proof}
